% Economic life and occupations — Anglais Tle Tech — Cours signé OnBuch+
% Programme officiel MINESEC. Compiler avec : tectonic N02-economic-life-and-occupations.tex
\newcommand{\DOCMATIERE}{Anglais}
\newcommand{\DOCNIVEAU}{Tle Tech}
\newcommand{\DOCMODULE}{Anglais – classe de Terminale technique}
\newcommand{\DOCLECON}{N02}
\newcommand{\DOCTITRE}{Economic life and occupations}
% OnBuch+ — préambule « cours signés » ENRICHI (compilé avec tectonic / XeLaTeX)
% Système visuel « Pop » d'OnBuch+ : fond crème (jamais blanc), bordures encre
% épaisses, ombres « dures » décalées, titres Archivo Black en capitales.
% Version MATHÉMATIQUES : graphes (pgfplots/tikz), boîtes pédagogiques
% (définition, propriété, méthode, attention, le savais-tu, exercice résolu,
% exercice, TP, bilan), page de garde et signature OnBuch+ ancrée en bas.
%
% Macros attendues dans main.tex AVANT \input{preamble} :
%   \DOCMATIERE  \DOCNIVEAU  \DOCMODULE  \DOCLECON (numéro)  \DOCTITRE
\documentclass[11pt]{article}

\usepackage{fontspec}
\usepackage[french]{babel}
\usepackage[a4paper,margin=2cm,top=3.4cm,bottom=2.8cm,headheight=1.4cm,footskip=1.3cm]{geometry}
\usepackage{amsmath,amssymb,amsfonts}
\usepackage{mathtools} % \xmapsto, \coloneqq... souvent utilisés par les modèles
\usepackage{cancel} % simplifications barrées (bilans de forces)
% \abs{z} (module, valeur absolue) et \norm{u} (norme), utilisés par les modèles
\providecommand{\abs}{}\let\abs\relax\DeclarePairedDelimiter{\abs}{\lvert}{\rvert}
\providecommand{\norm}{}\let\norm\relax\DeclarePairedDelimiter{\norm}{\lVert}{\rVert}
\usepackage[table]{xcolor} % [table] : fournit aussi \rowcolors (alternance de lignes)
\usepackage{pagecolor}
\usepackage{tikz}
\usetikzlibrary{arrows.meta,positioning,calc,shapes.geometric,shapes.misc,shapes.arrows,decorations.pathmorphing,decorations.pathreplacing,decorations.markings,patterns,shadows,mindmap,trees,fit,backgrounds,matrix,intersections,angles,quotes}
\usepackage{pgfplots}
\usepackage[version=4]{mhchem} % équations nucléaires \ce{^{235}_{92}U}
\usepackage[european,siunitx]{circuitikz} % schémas électriques
\pgfplotsset{compat=1.18}
\usepgfplotslibrary{fillbetween,groupplots} % aires entre courbes (calcul intégral)
\usepackage{enumitem}
\usepackage{fancyhdr}
% [most] charge toutes les bibliothèques tcolorbox (listings, minted, raster,
% documentation...) et gonfle inutilement la mémoire TeX sur un document long
% (~300 boîtes) : on ne charge que ce qui est réellement utilisé.
\usepackage[skins,breakable]{tcolorbox}
\usepackage{graphicx}
\usepackage{colortbl}
\usepackage{booktabs}
\usepackage{tabularx}
\usepackage{multirow}
\usepackage{diagbox} % case d'en-tête diagonale des tableaux à double entrée
% Colonnes extensibles centrée / gauche / droite (C, L, R), souvent utilisées par les modèles
\newcolumntype{C}{>{\centering\arraybackslash}X}
\newcolumntype{L}{>{\raggedright\arraybackslash}X}
\newcolumntype{R}{>{\raggedleft\arraybackslash}X}
\usepackage{array}
\usepackage{multicol}
\usepackage{siunitx}
% parse-numbers=false : accepte aussi \SI{2,5 \times 10^{-3}}{...}
\sisetup{locale=FR,output-decimal-marker={,},group-minimum-digits=4,parse-numbers=false}
\DeclareSIUnit\ohm{\ensuremath{\symup{\Omega}}} % U+2126 absent de la police
\DeclareSIUnit\division{div} % réglages d'oscilloscope (ms/div, V/div)
\DeclareSIUnit\tr{tr} % tours (vitesse de rotation, ex. \SI{1500}{\tr\per\minute})
\DeclareSIUnit\molar{M} % concentration molaire (ex. \SI{3}{\molar}, \SI{1}{\micro\molar})
\DeclareSIUnit\an{an} % année (le modèle confond parfois avec \year, un compteur TeX) — ex. \SI{5}{\kilo\meter\per\an}
\DeclareSIUnit\FCFA{FCFA} % franc CFA (ex. \SI{500}{\FCFA\per\kilo\gram})
\DeclareSIUnit\degreeBrix{\textdegree Brix} % teneur en sucre d'un sirop/jus (réfractométrie)
\DeclareSIUnit\month{mois} % le modèle utilise parfois \month, un compteur TeX, comme unité de durée
\DeclareSIUnit\microliter{\micro\liter} % le modèle écrit parfois \microliter en un seul token
\DeclareSIUnit\million{millions} % dénombrement (ex. \SI{15}{\million\per\mL} spermatozoïdes)
\usepackage{qrcode}
% \degree : commande fréquemment utilisée par les modèles pour un angle ou une
% température en dehors de \SI{}{} (non fournie par les paquets chargés).
\providecommand{\degree}{^{\circ}}
% \qed (fin de démonstration), utilisé par les modèles sans amsthm
\providecommand{\qed}{\hfill$\square$}
% \barre : double barre des valeurs interdites dans un tableau de variations (tkz-tab)
\providecommand{\barre}{\ensuremath{\|}}
% environnement proof (amsthm non chargé) utilisé par les modèles
\ifdefined\proof\else\newenvironment{proof}[1][Démonstration]{\par\noindent\textit{#1.}\ }{\qed\par}\fi
\usepackage{lastpage}
\usepackage{hyperref}
\hypersetup{hidelinks}

% ============================================================
% Palette « Pop » OnBuch+ — mêmes hex que la classe P (plus_ui.dart)
% ============================================================
\definecolor{poporange}{HTML}{F59321}
\definecolor{popdark}{HTML}{E07A0C}
\definecolor{popcream}{HTML}{FFF6E8}
\definecolor{popink}{HTML}{1C1714}
\definecolor{poppurple}{HTML}{7A5AE0}
\definecolor{popgreen}{HTML}{2E9E6A}
\definecolor{poppink}{HTML}{E0457B}
\definecolor{popblue}{HTML}{2F7DE1}
\definecolor{popgold}{HTML}{C98A12}
\definecolor{popmuted}{HTML}{8A7F76}
\definecolor{popline}{HTML}{EADFD2}
\definecolor{popgray}{HTML}{8A7F76}
\colorlet{popgrey}{popgray}
% Alias tolérants (le modèle nomme parfois les couleurs différemment)
\colorlet{popwhite}{white}
\colorlet{popblack}{popink}
\colorlet{popred}{poppink}
\colorlet{popyellow}{popgold}
% Teintes claires (fonds de tableaux/encadrés) — le modèle nomme parfois
% les couleurs avec un suffixe L (ex. popblueL) sans les avoir définies.
\colorlet{poporangeL}{poporange!20}
\colorlet{popdarkL}{popdark!20}
\colorlet{poppurpleL}{poppurple!20}
\colorlet{popgreenL}{popgreen!20}
\colorlet{poppinkL}{poppink!20}
\colorlet{popblueL}{popblue!20}
\colorlet{popgoldL}{popgold!20}
\colorlet{popredL}{popred!20}
\colorlet{popgrayL}{popgray!20}
% Teintes claires pour fonds de tableaux / zones
\colorlet{popblueL}{popblue!10!white}
\colorlet{poporangeL}{poporange!14!white}
\colorlet{popgreenL}{popgreen!12!white}
\colorlet{poppurpleL}{poppurple!12!white}
\colorlet{poppinkL}{poppink!12!white}
\colorlet{popgoldL}{popgold!14!white}

\pagecolor{popcream}

% ============================================================
% Typographie
% ============================================================
\defaultfontfeatures{Ligatures=TeX}
\setmainfont{PlusJakartaSans-Medium.ttf}[Path=../../../fonts/,BoldFont=PlusJakartaSans-Bold.ttf]
% Maths en sans-serif (Fira Math) avec chiffres et lettres droites de Plus
% Jakarta, pour que les formules se fondent dans le texte.
\usepackage[math-style=ISO,bold-style=ISO]{unicode-math}
\setmathfont{FiraMath-Regular.otf}[Path=../../../fonts/]
\setmathfont{PlusJakartaSans-Medium.ttf}[Path=../../../fonts/,range={up/num,up/{Latin,latin},\mathup}]
\usepackage{icomma} % virgule décimale sans espace en mode math
% Caractères mathématiques Unicode absents des polices de texte (Plus Jakarta,
% Archivo Black) : rendus par leur équivalent mathématique, en texte comme en
% formule — sinon ils s'affichent en carré vide (ex. « Arithmétique dans ℤ »).
\usepackage{newunicodechar}
\newunicodechar{ℝ}{\ensuremath{\mathbb{R}}}
\newunicodechar{ℤ}{\ensuremath{\mathbb{Z}}}
\newunicodechar{ℂ}{\ensuremath{\mathbb{C}}}
\newunicodechar{ℕ}{\ensuremath{\mathbb{N}}}
\newunicodechar{ℚ}{\ensuremath{\mathbb{Q}}}
\newunicodechar{𝔻}{\ensuremath{\mathbb{D}}}
\newunicodechar{α}{\ensuremath{\alpha}}
\newunicodechar{β}{\ensuremath{\beta}}
\newunicodechar{γ}{\ensuremath{\gamma}}
\newunicodechar{δ}{\ensuremath{\delta}}
\newunicodechar{ε}{\ensuremath{\varepsilon}}
\newunicodechar{θ}{\ensuremath{\theta}}
\newunicodechar{λ}{\ensuremath{\lambda}}
\newunicodechar{μ}{\ensuremath{\mu}}
\newunicodechar{σ}{\ensuremath{\sigma}}
\newunicodechar{τ}{\ensuremath{\tau}}
\newunicodechar{φ}{\ensuremath{\varphi}}
\newunicodechar{ψ}{\ensuremath{\psi}}
\newunicodechar{ω}{\ensuremath{\omega}}
\newunicodechar{Σ}{\ensuremath{\Sigma}}
% majuscules produites par \MakeUppercase dans les titres (α-aminés -> Α)
\newunicodechar{Α}{\ensuremath{\alpha}}
\newunicodechar{Β}{\ensuremath{\beta}}
\newunicodechar{Γ}{\ensuremath{\Gamma}}
\newunicodechar{Δ}{\ensuremath{\Delta}}
\newunicodechar{Ε}{\ensuremath{\varepsilon}}
\newunicodechar{Θ}{\ensuremath{\Theta}}
\newunicodechar{Λ}{\ensuremath{\Lambda}}
\newunicodechar{Μ}{\ensuremath{\mu}}
\newunicodechar{Τ}{\ensuremath{\tau}}
\newunicodechar{Φ}{\ensuremath{\Phi}}
\newunicodechar{Ψ}{\ensuremath{\Psi}}
\newunicodechar{Ω}{\ensuremath{\Omega}}
\newunicodechar{↦}{\ensuremath{\mapsto}}
\newunicodechar{∈}{\ensuremath{\in}}
\newunicodechar{∉}{\ensuremath{\notin}}
\newunicodechar{∧}{\ensuremath{\wedge}}
\newunicodechar{⇔}{\ensuremath{\Leftrightarrow}}
\newunicodechar{⟺}{\ensuremath{\Longleftrightarrow}}
\newunicodechar{⇒}{\ensuremath{\Rightarrow}}
\newunicodechar{⟹}{\ensuremath{\Longrightarrow}}
\newunicodechar{∘}{\ensuremath{\circ}}
\newunicodechar{∪}{\ensuremath{\cup}}
\newunicodechar{∩}{\ensuremath{\cap}}
\newunicodechar{≡}{\ensuremath{\equiv}}
\newunicodechar{⊂}{\ensuremath{\subset}}
\newunicodechar{⊥}{\ensuremath{\perp}}
\newunicodechar{∃}{\ensuremath{\exists}}
\newunicodechar{∀}{\ensuremath{\forall}}
\newunicodechar{∛}{\ensuremath{\sqrt[3]{\,}}}
\newunicodechar{∜}{\ensuremath{\sqrt[4]{\,}}}
\newunicodechar{⃗}{\ensuremath{{}^{\to}}}
\newunicodechar{ⁿ}{\ifmmode^{n}\else\textsuperscript{\ensuremath{n}}\fi}
\newunicodechar{ˣ}{\ifmmode^{x}\else\textsuperscript{\ensuremath{x}}\fi}
\newunicodechar{ᵇ}{\ifmmode^{b}\else\textsuperscript{\ensuremath{b}}\fi}
\newunicodechar{ʳ}{\ifmmode^{r}\else\textsuperscript{\ensuremath{r}}\fi}
\newunicodechar{ᵘ}{\ifmmode^{u}\else\textsuperscript{\ensuremath{u}}\fi}
\newunicodechar{ʸ}{\ifmmode^{y}\else\textsuperscript{\ensuremath{y}}\fi}
\newunicodechar{ᵃ}{\ifmmode^{a}\else\textsuperscript{\ensuremath{a}}\fi}
\newunicodechar{ᵉ}{\ifmmode^{e}\else\textsuperscript{\ensuremath{e}}\fi}
\newunicodechar{ᵏ}{\ifmmode^{k}\else\textsuperscript{\ensuremath{k}}\fi}
\newunicodechar{ᵖ}{\ifmmode^{p}\else\textsuperscript{\ensuremath{p}}\fi}
\newunicodechar{ᴺ}{\ifmmode^{N}\else\textsuperscript{\ensuremath{N}}\fi}
\newunicodechar{ᶜ}{\ifmmode^{c}\else\textsuperscript{\ensuremath{c}}\fi}
\newunicodechar{ʰ}{\ifmmode^{h}\else\textsuperscript{\ensuremath{h}}\fi}
\newunicodechar{ᵠ}{\ifmmode^{q}\else\textsuperscript{\ensuremath{q}}\fi}
\newunicodechar{ₙ}{\ifmmode_{n}\else\textsubscript{\ensuremath{n}}\fi}
\newunicodechar{ₐ}{\ifmmode_{a}\else\textsubscript{\ensuremath{a}}\fi}
\newunicodechar{ᵢ}{\ifmmode_{i}\else\textsubscript{\ensuremath{i}}\fi}
\newunicodechar{ᵣ}{\ifmmode_{r}\else\textsubscript{\ensuremath{r}}\fi}
\newunicodechar{ₖ}{\ifmmode_{k}\else\textsubscript{\ensuremath{k}}\fi}
\newunicodechar{ᵦ}{\ifmmode_{\beta}\else\textsubscript{\ensuremath{\beta}}\fi}
\newunicodechar{ₓ}{\ifmmode_{x}\else\textsubscript{\ensuremath{x}}\fi}
\newfontfamily\popdisplay{ArchivoBlack-Regular.ttf}[Path=../../../fonts/]
\newfontfamily\poplogo{SpaceGrotesk-Bold.ttf}[Path=../../../fonts/]
\newfontfamily\popsemi{PlusJakartaSans-SemiBold.ttf}[Path=../../../fonts/]
\newfontfamily\popxbold{PlusJakartaSans-ExtraBold.ttf}[Path=../../../fonts/]
\newfontfamily\popxboldls{PlusJakartaSans-ExtraBold.ttf}[Path=../../../fonts/,LetterSpace=3.0]

\setlist[enumerate]{leftmargin=1.7em,itemsep=5pt,topsep=4pt}
\setlist[itemize]{leftmargin=1.7em,itemsep=4pt,topsep=4pt,label={\color{poporange}\textbullet}}
\setlength{\parindent}{0pt}
\setlength{\parskip}{5pt}
\renewcommand{\arraystretch}{1.25}

% pgfplots : style Pop par défaut
\pgfplotsset{
  every axis/.append style={axis line style={popink,line width=1pt},tick style={popink},
    grid style={popline,line width=0.5pt},label style={font=\small},tick label style={font=\footnotesize}},
  popcurve/.style={poporange,line width=1.8pt,no markers,smooth},
}
% Même style disponible dans un \draw TikZ simple (hors pgfplots)
\tikzset{popcurve/.style={poporange,line width=1.8pt}}
\tikzset{popgrid/.style={popline,very thin}}
\colorlet{popcurve}{poporange} % le nom du style est parfois utilisé comme couleur

% ============================================================
% Pill / squiggle
% ============================================================
\newcommand{\poppill}[3]{%
  \tikz[baseline=-0.55ex]\node[draw=popink,line width=1.2pt,rounded corners=2.6pt,%
    fill=#2,inner xsep=6.5pt,inner ysep=3pt]{%
    \popxboldls\fontsize{7}{7}\selectfont\color{#3}\MakeUppercase{#1}};%
}
\newcommand{\popsquiggle}[1]{%
  \tikz[baseline=-0.4ex]{
    \draw[#1,line width=2.6pt,line cap=round]
      plot [smooth] coordinates {(0,0.09) (0.34,0.01) (0.68,0.21) (1.02,0.01) (1.36,0.10)};
  }%
}
% Mot-clé surligné (marqueur orange sous le texte)
\newcommand{\cle}[1]{{\popxbold\color{popdark}#1}}

% ============================================================
% En-tête / pied
% ============================================================
\pagestyle{fancy}
\fancyhf{}
\renewcommand{\headrulewidth}{0pt}
\renewcommand{\footrulewidth}{0pt}
\lhead{\raisebox{-0.15ex}{\poplogo\fontsize{13}{13}\selectfont\color{popink}OnBuch\color{poporange}+}}
\rhead{\poppill{\DOCMATIERE\ \textbullet\ \DOCNIVEAU\ \textbullet\ Leçon \DOCLECON}{white}{popink}}
\lfoot{\popxbold\fontsize{7.6}{7.6}\selectfont\color{popmuted}\MakeUppercase{Cours signé OnBuch+}\ \textbullet\ {\popsemi\DOCTITRE}}
\rfoot{\tikz[baseline=-0.6ex]\node[rounded corners=8.5pt,fill=popink,minimum height=17pt,minimum width=17pt,inner xsep=5pt,inner sep=0]{\popxbold\fontsize{8}{8}\selectfont\color{popcream}\thepage\,/\,\pageref*{LastPage}};}

% ============================================================
% Titres
% ============================================================
\newcommand{\chaptitre}[2]{%
  \begin{tcolorbox}[enhanced,colback=poporange,colframe=popink,boxrule=2.6pt,arc=11pt,
      drop shadow=popink,left=15pt,right=15pt,top=12pt,bottom=11pt,before skip=3pt,after skip=18pt]
    \poppill{#2}{popink}{popcream}\par\vspace{9pt}
    {\raggedright\hyphenpenalty=10000\exhyphenpenalty=10000\popdisplay\fontsize{22}{24}\selectfont\color{popcream}\MakeUppercase{#1}\par}
  \end{tcolorbox}
}
% Section numérotée : pastille orange avec le numéro + titre Archivo
\newcounter{popsec}
\newcommand{\coursec}[1]{%
  \stepcounter{popsec}\setcounter{popsub}{0}%
  \par\vspace{16pt}\noindent
  \begin{minipage}{\linewidth}
  \tikz[baseline=-0.7ex]\node[draw=popink,line width=1.6pt,fill=poporange,rounded corners=4pt,minimum size=22pt,inner sep=0]{\popdisplay\fontsize{12}{12}\selectfont\color{popcream}\thepopsec};%
  \hspace{8pt}{\popdisplay\fontsize{15}{17}\selectfont\color{popink}\MakeUppercase{#1}}\par
  \vspace{4pt}\noindent{\color{poporange}\rule{36pt}{3.6pt}}
  \end{minipage}\par\nopagebreak\vspace{8pt}
}
\newcounter{popsub}
\newcommand{\courssub}[1]{%
  \stepcounter{popsub}%
  \par\vspace{9pt}\noindent{\popxbold\fontsize{11.5}{13}\selectfont\color{popink}{\color{poporange}\thepopsec.\thepopsub}\hspace{6pt}#1}\par\nopagebreak\vspace{4pt}
}

% ============================================================
% Boîtes pédagogiques — même recette : fond blanc, bordure épaisse colorée,
% ombre dure encre, badge plein. Titre optionnel [#1].
% ============================================================
\tcbset{popbase/.style={breakable,enhanced,colback=white,boxrule=2.3pt,arc=9pt,
  shadow={3pt}{-3pt}{0pt}{popink},left=14pt,right=14pt,top=11pt,bottom=11pt,before skip=11pt,after skip=15pt,
  fontupper=\popsemi\fontsize{10.6}{14.5}\selectfont\color{popink},
  fontlower=\popsemi\fontsize{10.6}{14.5}\selectfont\color{popink}}}
% Coche dessinée (le glyphe ✓ n'existe pas dans les polices du document)
\newcommand{\popcheck}[1]{\tikz[baseline=-0.5ex]\draw[#1,line width=1.6pt,line cap=round,line join=round](-0.13,0.0)--(-0.04,-0.09)--(0.13,0.1);}
\providecommand{\coche}{\popcheck{popgreen}}
\providecommand{\resultat}[1]{\par\smallskip\noindent\fcolorbox{popgreen}{popgreenL}{\parbox{\dimexpr\linewidth-2\fboxsep-2\fboxrule\relax}{\textbf{Résultat.} #1}}\par\smallskip}
\newcommand{\popboxhead}[3]{\poppill{#1}{#2}{white}\ifx\relax#3\relax\else\hspace{6pt}{\popxbold\fontsize{10.6}{12}\selectfont\color{popink}#3}\fi\par\vspace{7pt}}

\newtcolorbox{aretenir}[1][]{popbase,colframe=popgreen,before upper={\popboxhead{À retenir}{popgreen}{#1}}}
\newtcolorbox{exemplebox}[1][]{popbase,colframe=poporange,before upper={\popboxhead{Exemple}{poporange}{#1}}}
\newtcolorbox{definition}[1][]{popbase,colframe=popblue,before upper={\popboxhead{Définition}{popblue}{#1}}}
\newtcolorbox{propriete}[1][]{popbase,colframe=poppurple,before upper={\popboxhead{Propriété}{poppurple}{#1}}}
\newtcolorbox{methode}[1][]{popbase,colframe=popgold,before upper={\popboxhead{Méthode}{popgold}{#1}}}
\newtcolorbox{attention}[1][]{popbase,colframe=poppink,before upper={\popboxhead{Attention}{poppink}{#1}}}
\newtcolorbox{experience}[1][]{popbase,colframe=popblue,colback=popblueL,before upper={\popboxhead{Expérience}{popblue}{#1}}}
% « Le savais-tu ? » : carte violette pleine (contexte, culture, Cameroun)
\newtcolorbox{savaistu}[1][]{popbase,colback=poppurple,colframe=popink,
  fontupper=\popsemi\fontsize{10.4}{14.2}\selectfont\color{white},
  before upper={\poppill{Le savais-tu\,?}{white}{poppurple}\ifx\relax#1\relax\else\hspace{6pt}{\popxbold\color{white}#1}\fi\par\vspace{7pt}}}
% Exercice résolu : énoncé en haut, \tcblower puis la solution sur fond crème
\newcounter{popexo}\newcounter{popexr}
\newtcolorbox{exoresolu}[1][]{popbase,colframe=poporange,colbacklower=popcream,
  segmentation style={popink,line width=1.2pt,dashed},
  before upper={\stepcounter{popexr}\popboxhead{Exercice résolu \thepopexr}{poporange}{#1}},
  before lower={\poppill{Solution}{popink}{popcream}\par\vspace{6pt}}}
% Exercice (non résolu en place) — la correction est dans la partie Corrigés
\newtcolorbox{exercice}[1][]{popbase,colframe=popink,
  before upper={\stepcounter{popexo}\popboxhead{Exercice \thepopexo}{popink}{#1}}}
% Corrigé
\newtcolorbox{corrige}[1][]{popbase,colframe=popgreen,colback=popgreenL,
  before upper={\popboxhead{Corrigé}{popgreen}{#1}}}
% Objectifs / prérequis / bilan
\newtcolorbox{objectifs}{popbase,colframe=popink,colback=poporangeL,
  before upper={\popboxhead{Objectifs de la leçon}{popink}{}}}
\newtcolorbox{prerequis}{popbase,colframe=popink,colback=popblueL,
  before upper={\popboxhead{Prérequis}{popblue}{}}}
\newtcolorbox{bilan}{popbase,colframe=popink,colback=white,boxrule=2.8pt,
  before upper={\popboxhead{Fiche bilan}{poporange}{L'essentiel à connaître pour l'examen}}}
% Figure encadrée avec légende
\newtcolorbox{popfigure}[1][]{enhanced,breakable=false,colback=white,colframe=popline,boxrule=1.4pt,arc=8pt,
  left=10pt,right=10pt,top=10pt,bottom=8pt,before skip=10pt,after skip=12pt,halign=center,
  lower separated=false,halign lower=center,
  fontlower=\popsemi\fontsize{9}{11.5}\selectfont\color{popmuted},
  #1}
\newcommand{\legende}[1]{\par\vspace{6pt}{\popsemi\fontsize{9}{11.5}\selectfont\color{popmuted}#1}}

% ============================================================
% Signature OnBuch+
% ============================================================
\newcommand{\poplogoinline}[1]{{\poplogo\fontsize{#1}{#1}\selectfont\color{popink}OnBuch\color{poporange}+}}
\newcommand{\signatureonbuch}{%
  \begin{tcolorbox}[enhanced,colback=popink,colframe=popink,boxrule=2.2pt,arc=10pt,
      drop shadow=poporange,left=14pt,right=14pt,top=10pt,bottom=10pt,before skip=0pt,after skip=0pt]
    \tikz[baseline=-0.7ex]\node[circle,fill=popgreen,draw=popcream,line width=1.3pt,minimum size=22pt,inner sep=0]{\popcheck{white}};%
    \hspace{9pt}\begin{minipage}[c]{0.62\linewidth}
      {\popxbold\fontsize{12}{13}\selectfont\color{popcream}Signé\ \textbullet\ {\poplogo OnBuch\color{poporange}+}}\par\vspace{2pt}
      {\raggedright\popsemi\fontsize{8}{10.5}\selectfont\color{popline}\MakeUppercase{Équipe pédagogique OnBuch+}\\\MakeUppercase{Conforme au programme officiel MINESEC}\par}
    \end{minipage}\hfill
    {\poplogo\fontsize{22}{22}\selectfont\color{popcream}OnBuch\color{poporange}+}
  \end{tcolorbox}%
}

% ============================================================
% Page de garde
% #1 titre · #2 matière · #3 niveau · #4 module · #5 n° leçon · #6 durée ·
% #7 description / objectifs (texte ou itemize)
% ============================================================
\newcommand{\pagegarde}[7]{%
  \thispagestyle{empty}
  \newgeometry{a4paper,margin=1.7cm,top=1.6cm,bottom=1.4cm}%
  \begin{tikzpicture}[remember picture,overlay]
    \fill[poporange!18] ([xshift=-3.2cm,yshift=-3.6cm]current page.north east) circle (4.6cm);
    \fill[poppurple!14] ([xshift=2.4cm,yshift=7.6cm]current page.south west) circle (3.8cm);
  \end{tikzpicture}%
  {\poplogo\fontsize{30}{30}\selectfont\color{popink}OnBuch\color{poporange}+}\hfill
  \raisebox{0.6ex}{\poppill{Cours signé \textbullet\ Premium}{popink}{popcream}}\par
  \vspace{1pt}\popsquiggle{poppurple}\par
  \vspace{8pt}
  \begin{tcolorbox}[enhanced,colback=poporange,colframe=popink,boxrule=2.8pt,arc=13pt,
      drop shadow=popink,left=18pt,right=18pt,top=12pt,bottom=13pt,after skip=10pt]
    \poppill{#2 \textbullet\ #3}{popink}{popcream}\hspace{5pt}\poppill{#4}{white}{popink}\par\vspace{8pt}
    {\popdisplay\fontsize{14}{14}\selectfont\color{popink}LEÇON #5}\par\vspace{4pt}
    {\raggedright\hyphenpenalty=10000\exhyphenpenalty=10000\popdisplay\fontsize{25}{27}\selectfont\color{popcream}\MakeUppercase{#1}\par}
    \vspace{8pt}
    {\popxbold\fontsize{9.5}{9.5}\selectfont\color{popink}\MakeUppercase{Durée conseillée : #6}}
  \end{tcolorbox}
  \begin{tcolorbox}[enhanced,colback=white,colframe=popink,boxrule=2.2pt,arc=10pt,
      drop shadow=popink,left=16pt,right=16pt,top=10pt,bottom=10pt,after skip=8pt]
    \poppill{Dans cette leçon}{poppurple}{white}\par\vspace{5pt}
    \popsemi\fontsize{9.3}{12.6}\selectfont\color{popink}#7
  \end{tcolorbox}
  \noindent\begin{minipage}[t]{0.487\linewidth}
    \begin{tcolorbox}[enhanced,colback=popgreen,colframe=popink,boxrule=2.2pt,arc=10pt,
        drop shadow=popink,left=11pt,right=11pt,top=10pt,bottom=10pt,height=3.1cm,valign=center]
      \tcbox[colback=white,colframe=popink,boxrule=1.3pt,arc=5pt,boxsep=0pt,nobeforeafter,
        left=3pt,right=3pt,top=3pt,bottom=3pt,valign=center]{\qrcode[height=1.9cm]{https://play.google.com/store/apps/details?id=com.luvvix.onbuch&hl=fr}}%
      \hspace{8pt}\begin{minipage}[c]{0.5\linewidth}
        {\popdisplay\fontsize{11}{12.5}\selectfont\color{white}\MakeUppercase{Télécharge OnBuch}}\par\vspace{4pt}
        {\raggedright\popsemi\fontsize{8.4}{11}\selectfont\color{white}Gratuit \textbullet\ Google Play\\Tuteur IA, annales, concours, résultats\par}
      \end{minipage}
    \end{tcolorbox}
  \end{minipage}\hfill
  \begin{minipage}[t]{0.487\linewidth}
    \begin{tcolorbox}[enhanced,colback=poppurple,colframe=popink,boxrule=2.2pt,arc=10pt,
        drop shadow=popink,left=11pt,right=11pt,top=10pt,bottom=10pt,height=3.1cm,valign=center]
      \tikz[baseline=-0.7ex]\node[circle,fill=white,draw=popink,line width=1.3pt,minimum size=1.6cm,inner sep=0]{\popdisplay\fontsize{24}{24}\selectfont\color{poppurple}+};%
      \hspace{8pt}\begin{minipage}[c]{0.6\linewidth}
        {\popdisplay\fontsize{11}{12.5}\selectfont\color{white}\MakeUppercase{Passe à OnBuch+}}\par\vspace{4pt}
        {\raggedright\popsemi\fontsize{8.4}{11}\selectfont\color{white}Tous les cours signés, Tuteur IA illimité, fiches PDF — onglet OnBuch+ de l'app\par}
      \end{minipage}
    \end{tcolorbox}
  \end{minipage}
  \vspace{8pt}
  \popcontenu
  \vfill
  \signatureonbuch
  \restoregeometry
  \newpage
}

% Rangée « Ce cours contient » (page de garde)
\newcommand{\poptile}[3]{% #1 couleur #2 gros texte #3 légende
  \begin{tcolorbox}[enhanced,colback=#1,colframe=popink,boxrule=2pt,arc=9pt,shadow={2.5pt}{-2.5pt}{0pt}{popink},
      left=6pt,right=6pt,top=9pt,bottom=9pt,halign=center,valign=center,height=1.75cm,width=\linewidth,nobeforeafter]
    {\popdisplay\fontsize{12.5}{13}\selectfont\color{white}\MakeUppercase{#2}}\par\vspace{3pt}
    {\centering\popsemi\fontsize{7.8}{9.5}\selectfont\color{white}#3\par}
  \end{tcolorbox}}
\newcommand{\popcontenu}{%
  \par\noindent{\popxboldls\fontsize{7.5}{7.5}\selectfont\color{popmuted}CE COURS SIGNÉ CONTIENT}\par\vspace{7pt}
  \noindent\begin{minipage}[t]{0.235\linewidth}\poptile{popblue}{Cours}{complet, illustré, pas à pas}\end{minipage}\hfill
  \begin{minipage}[t]{0.235\linewidth}\poptile{poporange}{Méthodes}{et exercices résolus}\end{minipage}\hfill
  \begin{minipage}[t]{0.235\linewidth}\poptile{poppink}{Exercices}{type examen + corrigés}\end{minipage}\hfill
  \begin{minipage}[t]{0.235\linewidth}\poptile{popgreen}{Fiche bilan}{l'essentiel à retenir}\end{minipage}\par}

% Sommaire
\newtcolorbox{sommaire}{enhanced,colback=white,colframe=popink,boxrule=2.2pt,arc=10pt,
  drop shadow=popink,left=18pt,right=18pt,top=13pt,bottom=13pt,before skip=6pt,after skip=20pt,
  before upper={\poppill{Sommaire}{poppurple}{white}\par\vspace{9pt}\popsemi\fontsize{10.6}{14.5}\selectfont\color{popink}}}

% Signature de fin de cours — poussée tout en bas de la dernière page
\newcommand{\findecours}{%
  \par\vfill\nopagebreak
  \begin{center}\popsquiggle{poporange}\end{center}\vspace{4pt}
  \signatureonbuch
}
\AtBeginDocument{\def\llbracket{\mathopen{[\mkern-3.5mu[}}\def\rrbracket{\mathclose{]\mkern-3.5mu]}}} % intervalles d'entiers (glyphe ⟦ absent de la police)
% Glyphes absents de Fira Math : redéfinis à partir de glyphes disponibles
\AtBeginDocument{%
  \def\setminus{\mathbin{\backslash}}\let\smallsetminus\setminus
  \def\vdots{\mathord{\vbox{\baselineskip3.2pt\lineskiplimit0pt\kern4pt\hbox{.}\hbox{.}\hbox{.}}}}%
  \def\ddots{\mathinner{\mkern1mu\raise6.5pt\hbox{.}\mkern2mu\raise3.6pt\hbox{.}\mkern2mu\raise0.7pt\hbox{.}\mkern1mu}}%
  \def\triangle{\mathord{\scalebox{0.9}{$\symup{\Delta}$}}}%
  \def\nabla{\mathord{\raisebox{\depth}{\scalebox{1}[-1]{$\symup{\Delta}$}}}}%
  \def\checkmark{\popcheck{popdark}}%
  \def\star{\mathbin{\ast}}\def\bigstar{\mathord{\ast}}%
  \def\diamond{\mathbin{\scalebox{0.6}{$\Diamond$}}}%
  \def\bigcup{\mathop{\mathchoice{\raisebox{-0.3em}{\scalebox{1.7}{$\cup$}}}{\scalebox{1.25}{$\cup$}}{\cup}{\cup}}\displaylimits}%
  \def\bigcap{\mathop{\mathchoice{\raisebox{-0.3em}{\scalebox{1.7}{$\cap$}}}{\scalebox{1.25}{$\cap$}}{\cap}{\cap}}\displaylimits}%
  \def\lesseqqgtr{\mathrel{\lessgtr}}\def\bigoplus{\mathop{\oplus}\displaylimits}%
}
\providecommand{\ssi}{si et seulement si} % abréviation fréquente
\AtBeginDocument{\providecommand{\wideparen}{\overparen}} % arc de cercle

\begin{document}

\pagegarde{Economic life and occupations}{Anglais}{Tle Tech}{Anglais – classe de Terminale technique}{N02}{13 séances (fiche de progression harmonisée)}{Cette leçon développe les compétences en anglais technique pour la vie économique : évaluer les besoins des consommateurs, comparer prix et qualité, maîtriser les services bancaires, et rédiger des textes argumentés. Elle intègre grammaire (comparaisons, adverbes temporels, phrases complexes, préférences) et vocabulaire spécialisé à travers des contextes camerounais concrets.
\vspace{4pt}
\begin{multicols}{2}\small
\begin{itemize}
\item À la fin de cette leçon, je sais évaluer les besoins des consommateurs et présenter les ressources communautaires en anglais en utilisant le vocabulaire approprié.
\item Je sais comparer des prix et des qualités pour déterminer les meilleurs achats en employant les formes de comparaison (double, égalité, degré) et les adverbes temporels (for, since, ago, after, afterwards).
\item Je sais comprendre des textes oraux et écrits sur l'économie de consommation et les services bancaires au Cameroun.
\item Je sais construire des phrases complexes par coordination et subordination pour exprimer des idées nuancées sur la gestion financière.
\item Je sais exprimer mes préférences bancaires et justifier mes choix avec des arguments structurés.
\item Je sais rédiger une composition cohérente sur l'utilisation des services bancaires (ouverture de compte, crédit, épargne) en respectant les conventions de l'écrit formel.
\end{itemize}\end{multicols}}

\begin{sommaire}
\begin{multicols}{2}
\begin{enumerate}
\item 1. Assessing Consumer Needs \& Community Resources
\item 2. Smart Shopping: Comparing Prices \& Quality for Best Buys
\item 3. Banking Services in Cameroon: Reading \& Vocabulary
\item 4. Complex Sentences: Subordination \& Expressing Preferences
\item 5. Writing Workshop: Composing About Banking Services
\item 6. Integration, Evaluation \& Remediation
\item Activité d'intégration
\item Méthodes et exercices résolus
\item Exercices
\item Corrigés des exercices
\item Fiche bilan
\end{enumerate}
\end{multicols}
\end{sommaire}

\begin{objectifs}
\begin{itemize}
\item À la fin de cette leçon, je sais évaluer les besoins des consommateurs et présenter les ressources communautaires en anglais en utilisant le vocabulaire approprié.
\item Je sais comparer des prix et des qualités pour déterminer les meilleurs achats en employant les formes de comparaison (double, égalité, degré) et les adverbes temporels (for, since, ago, after, afterwards).
\item Je sais comprendre des textes oraux et écrits sur l'économie de consommation et les services bancaires au Cameroun.
\item Je sais construire des phrases complexes par coordination et subordination pour exprimer des idées nuancées sur la gestion financière.
\item Je sais exprimer mes préférences bancaires et justifier mes choix avec des arguments structurés.
\item Je sais rédiger une composition cohérente sur l'utilisation des services bancaires (ouverture de compte, crédit, épargne) en respectant les conventions de l'écrit formel.
\item Je sais appliquer ces notions à des situations réelles de la vie courante camerounaise (marché, banque, coopérative) et justifier ma démarche.
\end{itemize}
\end{objectifs}

\begin{prerequis}
\begin{itemize}
\item Maîtrise des comparatifs de supériorité, infériorité et égalité de base (4ème/3ème).
\item Utilisation du present perfect avec for/since vs past simple avec ago (début 1ère).
\item Vocabulaire de base sur l'argent, les achats et la banque (1ère).
\item Construction de phrases simples et coordonnées (and, but, or) (1ère).
\item Rédaction d'un paragraphe structuré (introduction, développement, conclusion) (1ère).
\end{itemize}
\end{prerequis}

\newpage

\coursec{Assessing Consumer Needs \& Community Resources}

\courssub{Introduction: Marie's Dilemma}

Imagine Douala, on a busy Monday morning. \textbf{Marie} wants to open a small neighbourhood grocery shop (\emph{une épicerie de quartier}) in Bonabéri. She has a little capital, big dreams, and many questions. She hesitates between two banks for a microcredit (\emph{un microcrédit}), she does not know how to assess the real needs of her future customers, and she cannot compare the prices of wholesalers (\emph{grossistes}) efficiently. Her cousin \textbf{Jean}, on the other hand, has just opened a savings account (\emph{un compte épargne}) without understanding the hidden fees (\emph{les frais cachés}).

\textbf{Problematic:} How can we help Marie and Jean take the best economic decisions for their community? This lesson gives you the linguistic and practical tools to \cle{analyse}, \cle{compare}, \cle{argue} and \cle{write} in English in authentic Cameroonian financial situations. In this first section, we focus on the very first step of any economic project: \textbf{assessing consumer needs and community resources}.

\courssub{What Is a Needs Assessment?}

\textbf{Intuition first.} Before you buy stock for a shop, you must know what people around you actually want to buy. If Marie fills her shelves with expensive imported rice but her neighbours can only afford local rice, her shop will fail. A good businessperson first \emph{listens} to the community, then \emph{acts}.

\begin{definition}[Needs assessment]
A \cle{needs assessment} (\emph{évaluation des besoins}) is a \textbf{systematic process} of identifying the gaps (\emph{écarts}) between the \textbf{current situation} and the \textbf{desired situation} for a \textbf{target population} (\emph{population cible}). It answers three questions: What do people have? What do people need? What is missing?
\end{definition}

In Cameroon, the main consumer needs concern four sectors:
\begin{itemize}
\item \textbf{Food} (\emph{alimentation}): rice, oil, soap, cassava, plantain, fish;
\item \textbf{Health} (\emph{santé}): medicines, health centres, clean water;
\item \textbf{Education} (\emph{éducation}): school fees, uniforms, textbooks;
\item \textbf{Transport} (\emph{transport}): taxis, motorbikes (\emph{bendskins}), bus fares.
\end{itemize}

\begin{methode}[How to conduct a needs assessment in 5 steps]
\begin{enumerate}
\item \textbf{Design a survey} (\emph{concevoir une enquête}): write a short \cle{questionnaire} with clear questions (age, income, daily expenses, preferred products).
\item \textbf{Identify the target group} (\emph{groupe cible}): for Marie, the households within 500 metres of her shop.
\item \textbf{Collect data} (\emph{collecter les données}): interview people at the \cle{local market}, at home, or at the GIE meeting.
\item \textbf{Analyse the results} (\emph{analyser}): calculate percentages, compare prices, rank the \cle{priorities}.
\item \textbf{Write a report and act} (\emph{rédiger un rapport et agir}): present conclusions and take decisions (what to stock, at what price).
\end{enumerate}
\end{methode}

\begin{popfigure}
\begin{center}
\begin{tikzpicture}[node distance=0.55cm, every node/.style={font=\small}]
\node[draw=popblue, fill=popblueL, rounded corners, minimum width=3.1cm, minimum height=0.95cm, align=center] (s1) {\textbf{1. Survey}\\ questionnaire};
\node[draw=poporange, fill=poporangeL, rounded corners, minimum width=3.1cm, minimum height=0.95cm, align=center, right=of s1] (s2) {\textbf{2. Data collection}\\ interviews};
\node[draw=popgreen, fill=popgreenL, rounded corners, minimum width=3.1cm, minimum height=0.95cm, align=center, right=of s2] (s3) {\textbf{3. Analysis}\\ percentages};
\node[draw=poppurple, fill=poppurpleL, rounded corners, minimum width=3.1cm, minimum height=0.95cm, align=center, right=of s3] (s4) {\textbf{4. Report}\\ conclusions};
\node[draw=popgold, fill=popgoldL, rounded corners, minimum width=3.1cm, minimum height=0.95cm, align=center, right=of s4] (s5) {\textbf{5. Action}\\ decisions};
\draw[-{Stealth}, thick, popdark] (s1) -- (s2);
\draw[-{Stealth}, thick, popdark] (s2) -- (s3);
\draw[-{Stealth}, thick, popdark] (s3) -- (s4);
\draw[-{Stealth}, thick, popdark] (s4) -- (s5);
\end{tikzpicture}
\end{center}
\legende{The needs assessment flow: from the field survey at the local market to concrete economic action for the community.}
\end{popfigure}

\courssub{Key Vocabulary: Assessing Needs and Community Resources}

\begin{center}
\begin{tabularx}{\linewidth}{|l|X|l|}
\hline
\rowcolor{popblueL} \textbf{English} & \textbf{Meaning / example} & \textbf{French} \\
\hline
needs assessment & a needs assessment of the neighbourhood & évaluation des besoins \\
\hline
survey / questionnaire & to carry out a survey in the market & enquête / questionnaire \\
\hline
target group & women aged 25--45 are our target group & groupe cible \\
\hline
community resources & tontines, GIEs and cooperatives & ressources communautaires \\
\hline
local market & the Mokolo market in Yaoundé & marché local \\
\hline
purchasing power & low purchasing power of households & pouvoir d'achat \\
\hline
priority & Food is the first priority. & priorité \\
\hline
affordability & the affordability of basic goods & caractère abordable \\
\hline
accessibility & accessibility of health services & accessibilité \\
\hline
household & an average household of six people & ménage \\
\hline
\end{tabularx}
\end{center}

\begin{attention}[False friends and tricky words]
Do not confuse \cle{affordability} (can people pay?) with \cle{accessibility} (can people reach it?). A product can be affordable but not accessible if the shop is too far. Also, \emph{to achieve} means \og réaliser un objectif \fg{}, not \og acheter \fg{} (to buy / to purchase).
\end{attention}

\courssub{Grammar Focus: Advanced Forms of Comparison}

To analyse survey data, you must compare prices, incomes and needs precisely. English offers three powerful tools beyond the simple comparative.

\textbf{1. Equality: as \ldots\ as / the same \ldots\ as.} Use \emph{as + adjective + as} to say two things are equal: \og Local rice is \textbf{as expensive as} imported rice in some markets. \fg{} Use \emph{the same + noun + as}: \og This wholesaler offers \textbf{the same quality as} the other one. \fg{}

\textbf{2. Double comparatives: the more \ldots, the more \ldots} This structure shows that two things change together.

\begin{methode}[Building a double comparative]
Use: \textbf{The + comparative + subject + verb, the + comparative + subject + verb.}\\
Example: \emph{The higher the price, the lower the demand.} (\emph{Plus le prix est élevé, plus la demande est faible.})
\end{methode}

\begin{exemplebox}[Double comparatives in context]
\begin{itemize}
\item \emph{The more customers Marie interviews, the better she understands their needs.}
\item \emph{The cheaper the goods, the faster they sell.}
\item \emph{The more you save, the safer your business becomes.}
\end{itemize}
\end{exemplebox}

\textbf{3. Degree: far more, slightly less, much better.} Intensify or soften a comparison with adverbs of degree:

\begin{center}
\begin{tabularx}{\linewidth}{|l|X|}
\hline
\rowcolor{popgreenL} \textbf{Degree word} & \textbf{Example} \\
\hline
far / much + comparative & Imported rice is \emph{far more expensive} than local rice. \\
\hline
slightly / a little + comparative & Soap is \emph{slightly cheaper} in Douala than in Yaoundé. \\
\hline
much / far + better/worse & The cooperative's oil is \emph{much better} than the market oil. \\
\hline
not nearly as \ldots\ as & The new shop is \emph{not nearly as busy as} the Central Market. \\
\hline
\end{tabularx}
\end{center}

\begin{attention}[as \ldots\ as or so \ldots\ as?]
A frequent error: confusing \emph{as \ldots\ as} (equality, affirmative) and \emph{so \ldots\ as} (used only in negative sentences). In the affirmative, \textbf{always} use \emph{as \ldots\ as}: \og Local rice is \textbf{as} good \textbf{as} imported rice \fg{} --- never \og so good as \fg{} in an affirmative sentence. Negative: \og It is \textbf{not as/so} expensive \textbf{as} before. \fg{}
\end{attention}

\courssub{Analysing Real Data: Prices in Douala Markets}

Numbers speak. According to INS and MINCOMMERCE data, the prices of basic goods (\emph{denrées de base}) vary strongly from one market to another. Let us study the price of rice --- the most consumed cereal in urban Cameroon.

\begin{exemplebox}[A key figure]
In 2023, the average price of local rice in Cameroon was \textbf{350 FCFA/kg}, against \textbf{450 FCFA/kg} for imported rice (source: MINCOMMERCE). Local rice is therefore \emph{much cheaper than} imported rice, and the gap explains why low-income households prefer it.
\end{exemplebox}

\begin{popfigure}
\begin{center}
\begin{tikzpicture}
\begin{axis}[
  ybar, bar width=10pt,
  width=0.92\linewidth, height=6.2cm,
  symbolic x coords={Marché Central,Mboppi,Deïdo,Sandaga,Nkouloulou},
  xtick=data, x tick label style={rotate=20, anchor=east, font=\footnotesize},
  ymin=0, ymax=560,
  ylabel={Price (FCFA/kg)}, ylabel style={font=\small},
  legend style={at={(0.02,0.97)}, anchor=north west, font=\small},
  ymajorgrids, grid style={popline!40},
  every axis/.append style={font=\small}
]
\addplot+[ybar, fill=popgreen, draw=popdark] coordinates {(Marché Central,340) (Mboppi,330) (Deïdo,350) (Sandaga,360) (Nkouloulou,370)};
\addplot+[ybar, fill=poporange, draw=popdark] coordinates {(Marché Central,440) (Mboppi,430) (Deïdo,450) (Sandaga,470) (Nkouloulou,480)};
\legend{Local rice, Imported rice}
\end{axis}
\end{tikzpicture}
\end{center}
\legende{Comparative bar chart: prices of local vs imported rice (FCFA/kg) in five Douala markets (INS 2023, indicative values). Imported rice is consistently far more expensive.}
\end{popfigure}

\textbf{Reading the chart.} In every market, imported rice is \emph{far more expensive than} local rice. The gap is largest in Nkouloulou (about 110 FCFA/kg) and smallest in Mboppi. We can say: \emph{The poorer the neighbourhood, the wider the price gap matters for the household budget.} For Marie, the conclusion is clear: stocking affordable local rice should be a \cle{priority}.

\courssub{Community Resources: The Informal Economy at Work}

A community is not only consumers; it is also \textbf{resources}. In Cameroon, economic life relies on powerful solidarity networks:

\begin{itemize}
\item \textbf{Tontines (njangi)}: rotating savings groups where members contribute monthly and receive the pot in turn;
\item \textbf{GIE} (Groupement d'Intérêt Économique): legal economic interest groupings, often run by women (shea butter, cassava processing, tailoring);
\item \textbf{Cooperatives} (\emph{coopératives}): farmers or traders pooling production and sales;
\item \textbf{Vocational training centres} (\emph{centres de formation professionnelle}): they train young people in trades and entrepreneurship.
\end{itemize}

\begin{savaistu}[Tontines: giants of informal finance]
Tontines (\emph{njangi}) manage more than \textbf{500 billion FCFA per year} in Cameroon. They are the first source of informal financing for small and medium enterprises (SMEs), far ahead of microfinance institutions. For many traders at the Central Market of Yaoundé, the tontine \emph{is} the bank.
\end{savaistu}

\begin{popfigure}
\begin{center}
\begin{tikzpicture}[
  root/.style={draw=popblue, fill=popblueL, rounded corners, font=\small\bfseries, align=center, minimum width=3.4cm, minimum height=0.9cm},
  child node/.style={draw=popdark, fill=popcream, rounded corners, font=\footnotesize, align=center, minimum width=2.6cm, minimum height=0.8cm},
  edge from parent/.style={draw=popmuted, thick},
  level distance=2.1cm, sibling distance=3.1cm,
  grow=right
]
\node[root, align=center]{Community\\ Resources}
  child { node[child node, align=center]{Vocational\\ training centres} }
  child { node[child node, align=center]{Cooperatives\\ (farmers, traders)} }
  child { node[child node, align=center]{GIE\\ (women's groups)} }
  child { node[child node, align=center]{Tontines\\ (njangi)} };
\end{tikzpicture}
\end{center}
\legende{Mind map of community resources in Cameroon: four pillars that finance, train and organise local economic life.}
\end{popfigure}

\courssub{Speaking Practice: Role-Plays at the Market}

\begin{exemplebox}[Model dialogue: interviewing a trader at the Central Market, Yaoundé]
\textbf{Enumerator (E):} Good morning, Madam. I am carrying out a \emph{survey} on consumer needs in this neighbourhood. May I ask you a few questions?\\
\textbf{Trader (T):} Of course, young lady.\\
\textbf{E:} What products do your customers buy most?\\
\textbf{T:} Rice, oil and soap. \emph{The cheaper the product, the faster it sells.}\\
\textbf{E:} Is local rice \emph{as popular as} imported rice?\\
\textbf{T:} Much more popular! It is \emph{far cheaper}, and the quality is now \emph{as good as} the imported one.\\
\textbf{E:} What is your customers' main difficulty?\\
\textbf{T:} Their \emph{purchasing power}. Salaries are low, so \emph{affordability} is the first \emph{priority}.
\end{exemplebox}

\textbf{Your turn.} In pairs, prepare a role-play: one student is an enumerator, the other is either a trader at the Marché Central or the leader of a women's GIE. Use at least three target vocabulary words and two comparison structures. Then, each group presents a \textbf{3-minute oral synthesis} (\emph{synthèse orale}) of its findings with a visual support (table or chart): state the target group, the three main needs, the price data, and one recommendation.

\courssub{Practice}

\begin{exoresolu}[Comparing prices with degree words]
Complete with the correct form: \emph{Imported rice is \ldots\ (expensive) local rice, and it sells \ldots\ (slowly).}\\
\tcblower
\textbf{Solution.} We compare two items with a degree adverb: \emph{Imported rice is \textbf{far more expensive than} local rice, and it sells \textbf{much more slowly}.} Note: \emph{expensive} is a long adjective, so the comparative is \emph{more expensive}; \emph{slowly} takes \emph{more slowly}.
\end{exoresolu}

\begin{exercice}[Grammar: comparisons]
\begin{enumerate}
\item Rewrite with a double comparative: \og When prices rise, demand falls. \fg
\item Correct the mistake: \og Local rice is so cheap as imported rice. \fg
\item Complete: \og The new market is \ldots\ (big) the old one, but \ldots\ (not / clean). \fg
\item Translate: \og Plus on épargne tôt, plus le projet devient possible. \fg
\end{enumerate}
\end{exercice}

\begin{corrige}[Exercice : comparisons]
\begin{enumerate}
\item \emph{The higher the prices, the lower the demand.}
\item \emph{Local rice is \textbf{as} cheap \textbf{as} imported rice.} (In the affirmative, use \emph{as \ldots\ as}, not \emph{so \ldots\ as}.)
\item \emph{The new market is \textbf{bigger than} the old one, but \textbf{not as clean as} it.}
\item \emph{The earlier you save, the more possible the project becomes.}
\end{enumerate}
\end{corrige}

\begin{aretenir}[Key points of this section]
\begin{itemize}
\item A \cle{needs assessment} identifies the gap between the current and the desired situation for a target group, in five steps: survey, data collection, analysis, report, action.
\item Master the three comparison tools: equality (\emph{as \ldots\ as}), double comparatives (\emph{the more \ldots, the more \ldots}) and degree (\emph{far more, slightly less, much better}).
\item Community resources --- tontines, GIEs, cooperatives, training centres --- are the backbone of Cameroon's local economy.
\item Always support your oral synthesis with data: prices, percentages, priorities.
\end{itemize}
\end{aretenir}

\coursec{Smart Shopping: Comparing Prices \& Quality for Best Buys}

In the previous section, we learned how to \cle{assess consumer needs} and identify \cle{community resources}. Now that we know \emph{what} we need and \emph{where} to find it, we must master the skill of \cle{smart shopping}: getting the best \cle{value for money}. Whether you are buying maize flour for your family in Yaoundé or spare parts for a workshop in Douala, the principles remain the same: \textbf{compare unit prices}, \textbf{check quality}, and \textbf{negotiate wisely}.

\courssub{Listening: Radio Ads \& Price Comparisons}

\begin{experience}[Listening to Radio Advertisements]
\textbf{Context:} You hear two radio advertisements. One is from a supermarket promoting a \cle{discount} on wholesale packs. The other is from a traditional market vendor (for example at Marché Central in Yaoundé or Marché Mboppi in Douala) advertising fresh produce at \cle{retail} prices.

\textbf{Task:} Listen carefully and fill in the table below with the \cle{key information}.
\begin{center}
\begin{tabularx}{\linewidth}{|l|X|X|}\hline
\rowcolor{popblueL}\textbf{Criteria} & \textbf{Supermarket Ad} & \textbf{Market Vendor Ad} \\ \hline
Product & & \\ \hline
Price (FCFA) & & \\ \hline
Quantity / Pack size & & \\ \hline
\cle{Unit price} (FCFA/kg or FCFA/L) & & \\ \hline
\cle{Expiry date} mentioned? & & \\ \hline
\cle{Quality check} / Certification (ANOR?) & & \\ \hline
Special offer (\cle{bargain}, loyalty card?) & & \\ \hline
\end{tabularx}
\end{center}

\textbf{Follow-up:} Which source offers the most \cle{cost-effective} solution for a family of five for one month? Justify your answer using the information in your table.
\end{experience}

\begin{definition}[Unit Price]
The \cle{unit price} is the cost per standard unit of measure (kilogram, litre, metre, piece). It allows a \textbf{real comparison} between different pack sizes (e.g. a 500 g box versus a 2 kg bag) or different sellers (supermarket versus market).
\[
\text{Unit Price} = \frac{\text{Total Price}}{\text{Total Quantity (in standard units)}}
\]
\textbf{Example:} A 2,5 kg bag of rice costs 2\,500 FCFA. Unit price $= 2\,500 \div 2,5 = \textbf{1\,000 FCFA/kg}$.
\end{definition}

\begin{attention}[Packaging Can Be Misleading]
A bigger pack is \textbf{not always} cheaper per kilogram. Some sellers use large packaging to attract buyers while the \cle{unit price} remains high. \textbf{Always calculate the unit price before deciding.} Also check the \cle{expiry date}: a cheap product that expires tomorrow is not a \cle{bargain}.
\end{attention}

\courssub{Vocabulary: Smart Shopping Terms}

\begin{center}
\begin{tabularx}{\linewidth}{|l|X|l|}\hline
\rowcolor{popblueL}\textbf{English Term} & \textbf{Definition / Context} & \textbf{French Gloss} \\ \hline
\cle{Bargain} & A thing bought for less than its usual price. & Bonne affaire \\ \hline
\cle{Discount} & A reduction in the usual price. & Réduction / Remise \\ \hline
\cle{Wholesale} & Selling goods in large quantities at low prices (to retailers). & Vente en gros \\ \hline
\cle{Retail} & Selling goods in small quantities directly to consumers. & Vente au détail \\ \hline
\cle{Unit price} & Price per kg, litre, metre... (see Definition box). & Prix unitaire \\ \hline
\cle{Quality check} & Inspecting goods before buying (freshness, seal, label). & Contrôle qualité \\ \hline
\cle{Expiry date} / \cle{Best before} & Date after which a product should not be eaten or used. & Date de péremption \\ \hline
\cle{Value for money} & Good quality relative to the price paid. & Rapport qualité-prix \\ \hline
\cle{Cost-effective} & Producing good results without costing too much. & Rentable / Économique \\ \hline
\cle{Budget-friendly} & Cheap enough not to exceed planned spending. & Abordable \\ \hline
\cle{Bulk purchase} & Buying large quantities at once (often with a discount). & Achat en gros \\ \hline
\cle{Receipt} & Paper proof of payment given by the seller. & Reçu / Facture \\ \hline
\cle{To negotiate} / \cle{To bargain} & To discuss the price with the seller to reduce it. & Négocier / Marchander \\ \hline
\end{tabularx}
\end{center}

\courssub{Grammar: Time Adverbs (for, since, ago, after, afterwards)}

These adverbs situate actions in time. Mastering them is crucial for describing \emph{price trends}, \emph{shopping habits}, or \emph{past transactions}.

\begin{propriete}[Time Adverbs: Forms \& Uses]
\begin{itemize}
\item \textbf{FOR} + \emph{duration} (How long?): \textit{I have shopped at this market \textbf{for 10 years}.} (depuis / pendant 10 ans). Used with the present perfect for an action continuing until now.
\item \textbf{SINCE} + \emph{starting point} (When did it start?): \textit{Prices have risen \textbf{since 2020}.} (depuis 2020). Also used with the present perfect.
\item \textbf{AGO} + \emph{past moment}, placed \textbf{after} the duration, with the Past Simple: \textit{I bought this chicken \textbf{two days ago}.} (il y a deux jours).
\item \textbf{AFTER} + \emph{noun / gerund / clause} (preposition or conjunction): \textit{\textbf{After the purchase}, I checked the receipt.} / \textit{\textbf{After paying}, I left.} / \textit{\textbf{After I paid}, I left.}
\item \textbf{AFTERWARDS} (adverb, stands alone): \textit{We compared prices. \textbf{Afterwards}, we decided.} (ensuite, par la suite).
\end{itemize}
\end{propriete}

\begin{attention}[After vs Afterwards]
\textbf{Do not confuse them!}
\begin{itemize}
\item \textbf{After} connects two events and is always followed by a noun, a gerund or a clause: \textit{After lunch, we went shopping.} / \textit{After we ate, we went shopping.}
\item \textbf{Afterwards} modifies the verb and stands alone, usually at the beginning or end of the sentence: \textit{We had lunch. Afterwards, we went shopping.}
\item \textbf{Never say:} ``Afterwards lunch'' or ``Afterwards we ate, we went shopping.''
\item \textbf{Never say:} ``I have seen him two days ago.'' With \textbf{ago}, always use the Past Simple: \textit{I saw him two days ago.}
\end{itemize}
\end{attention}

\begin{exercice}[Complete with for, since, ago, after or afterwards]
\begin{enumerate}
\item Fuel prices have been high \_\_\_\_\_\_\_\_ the increase in fuel prices in 2023.
\item I waited \_\_\_\_\_\_\_\_ 20 minutes at the checkout \_\_\_\_\_\_\_\_ paying.
\item She bought that fabric three months \_\_\_\_\_\_\_\_.
\item We compared the offers. \_\_\_\_\_\_\_\_, we chose the local brand.
\item \_\_\_\_\_\_\_\_ checking the \cle{expiry date}, put the milk in the fridge.
\item They have sold second-hand clothes in this market \_\_\_\_\_\_\_\_ fifteen years.
\end{enumerate}
\end{exercice}

\begin{corrige}[Exercice : Time Adverbs]
\begin{enumerate}
\item \textbf{since} (starting point + present perfect).
\item \textbf{for} (duration) ; \textbf{after} (after + gerund ``paying'').
\item \textbf{ago} (past moment + Past Simple).
\item \textbf{Afterwards} (adverb standing alone after a full stop).
\item \textbf{After} (after + gerund ``checking'').
\item \textbf{for} (duration + present perfect).
\end{enumerate}
\end{corrige}

\courssub{Practical Activity: Comparing Three Offers of Maize Flour}

\begin{methode}[How to Determine the \cle{Best Buy} (Quality and Price)]
Follow these steps before any important purchase:
\begin{enumerate}
\item \textbf{Identify your need:} which product, which quantity, for how many people?
\item \textbf{Collect the offers:} note the price and the quantity of each seller (supermarket, market, wholesaler).
\item \textbf{Calculate the unit price} of each offer: divide the total price by the quantity in standard units (kg or L).
\item \textbf{Check the quality:} read the label, verify the \cle{expiry date}, look for certification (for example ANOR standards in Cameroon), inspect freshness.
\item \textbf{Compare and rank} the offers from the cheapest to the most expensive per unit.
\item \textbf{Decide and justify:} the \cle{best buy} is the offer with the lowest unit price \emph{among products of acceptable quality} — not simply the cheapest product.
\end{enumerate}
\end{methode}

\begin{exoresolu}[Worked Example: Which Bag of Maize Flour Is the Best Buy?]
\textbf{Statement.} A housewife in Bafoussam finds three offers of maize flour:
\begin{center}
\begin{tabularx}{\linewidth}{|l|X|X|X|}\hline
\rowcolor{popblueL}\textbf{Offer} & \textbf{Quantity} & \textbf{Total Price} & \textbf{Quality observation} \\ \hline
A (market vendor) & 1 kg & 650 FCFA & Fresh, well dried \\ \hline
B (supermarket) & 5 kg & 3\,500 FCFA & Certified, expires in 2 months \\ \hline
C (wholesaler) & 10 kg & 6\,000 FCFA & Certified, expires in 6 months \\ \hline
\end{tabularx}
\end{center}
Calculate the unit price of each offer and determine the \cle{best buy} for a family that consumes 10 kg per month.
\tcblower
\textbf{Solution.}
\begin{align*}
\text{Offer A:} \quad & 650 \div 1 = 650 \text{ FCFA/kg}\\
\text{Offer B:} \quad & 3\,500 \div 5 = 700 \text{ FCFA/kg}\\
\text{Offer C:} \quad & 6\,000 \div 10 = 600 \text{ FCFA/kg}
\end{align*}
\textbf{Ranking (cheapest to most expensive per kg):} C (600 FCFA/kg) $<$ A (650 FCFA/kg) $<$ B (700 FCFA/kg).

\textbf{Conclusion:} Offer C is the \cle{best buy}: it has the lowest \cle{unit price} (600 FCFA/kg), it is certified, and its \cle{expiry date} (6 months) is long enough for a family consuming 10 kg per month. Offer B is the most expensive per kilogram despite its attractive packaging. Note that Offer A would be the best choice only for a buyer who cannot spend 6\,000 FCFA at once: a \cle{best buy} must also respect the buyer's \textbf{budget}.
\end{exoresolu}

\begin{exercice}[Your Turn: Best Buy at the Shop]
A school canteen in Bafoussam wants to buy cooking oil. Three offers are available:
\begin{center}
\begin{tabularx}{\linewidth}{|l|X|X|}\hline
\rowcolor{popblueL}\textbf{Offer} & \textbf{Quantity} & \textbf{Total Price} \\ \hline
X & 1 litre & 1\,300 FCFA \\ \hline
Y & 5 litres & 6\,000 FCFA \\ \hline
Z & 20 litres & 22\,000 FCFA \\ \hline
\end{tabularx}
\end{center}
\begin{enumerate}
\item Calculate the \cle{unit price} (FCFA per litre) of each offer.
\item Rank the offers from the cheapest to the most expensive per litre.
\item Which offer is the \cle{best buy} for the canteen? Justify your answer in two or three sentences using the vocabulary of the lesson (\emph{unit price, value for money, bulk purchase, cost-effective}).
\end{enumerate}
\end{exercice}

\begin{corrige}[Exercice : Best Buy at the Shop]
\begin{enumerate}
\item Offer X: $1\,300 \div 1 = 1\,300$ FCFA/L. Offer Y: $6\,000 \div 5 = 1\,200$ FCFA/L. Offer Z: $22\,000 \div 20 = 1\,100$ FCFA/L.
\item Z (1\,100 FCFA/L) $<$ Y (1\,200 FCFA/L) $<$ X (1\,300 FCFA/L).
\item Offer Z is the \cle{best buy}: it has the lowest \cle{unit price} (1\,100 FCFA/L). This \cle{bulk purchase} is \cle{cost-effective} for a canteen that uses large quantities of oil, so it offers the best \cle{value for money}.
\end{enumerate}
\end{corrige}

\begin{aretenir}[Key Points of Section 2]
\begin{itemize}
\item The \cle{unit price} $=$ total price $\div$ quantity: the essential tool for comparing offers.
\item The \cle{best buy} combines \textbf{low unit price}, \textbf{acceptable quality} (label, \cle{expiry date}, certification) and the buyer's \textbf{budget}.
\item \textbf{for} + duration, \textbf{since} + starting point (both with the present perfect); \textbf{ago} with the Past Simple; \textbf{after} + noun/gerund/clause; \textbf{afterwards} stands alone.
\item Key vocabulary: \cle{bargain}, \cle{discount}, \cle{wholesale}, \cle{retail}, \cle{value for money}, \cle{cost-effective}, \cle{bulk purchase}.
\end{itemize}
\end{aretenir}

\coursec{Banking Services in Cameroon: Reading \& Vocabulary}

\courssub{Transition from Smart Shopping to Banking Services}
In the previous section, we learned how to \cle{compare prices and quality} to make smart purchases as consumers. We used adverbs of time (\emph{for, since, ago}) and comparison structures. Now, we move a step further: once you have earned money, you need a safe place to keep it and tools to make it grow. This is where \cle{banking services} come in. In Cameroon, the landscape is evolving rapidly: traditional banks, microfinance institutions, and \cle{mobile money} operators (Orange Money, MTN MoMo) all compete for customers. Understanding their offers, reading their contracts, and writing about them are essential skills for your professional life.

\courssub{Authentic Reading: Bank Brochures \& Key Vocabulary}
Let us start by examining simplified extracts from real Cameroonian bank brochures (BICEC, Afriland First Bank, UBA, Ecobank). Read for gist first, then for detail.

\begin{exemplebox}[Brochure Extract: Opening a \emph{Compte Courant} (Current Account)]
\textbf{BICEC -- Bienvenue chez vous!}
Open your current account in 3 steps:
\begin{enumerate}
\item Provide your \textbf{KYC documents}: valid ID (CNI/Passport), proof of address (utility bill less than 3 months old), two passport photos.
\item Visit any branch for \textbf{signature} and initial \textbf{deposit} (minimum \num{5000} FCFA).
\item Receive your \textbf{debit card} and \textbf{cheque book} within 48 hours.
\end{enumerate}
\textit{Benefits:} Free SMS alerts, online banking access, overdraft facility (subject to approval).
\end{exemplebox}

\begin{exemplebox}[Brochure Extract: \emph{Épargne Jeune} / Youth Savings -- Afriland First Bank]
\textbf{Start early, grow big!}
\textbf{Interest rate:} \num{3,5}\% per annum (capitalised quarterly).
\textbf{Conditions:} Age 18--35. No monthly fees. Minimum balance: \num{1000} FCFA.
\textbf{Mobile link:} Link your account to \textbf{Orange Money} or \textbf{MTN MoMo} for instant \textbf{deposits/withdrawals}.
\textbf{Standing order:} Set up automatic monthly transfers from your salary account.
\end{exemplebox}

\begin{definition}[Financial Inclusion]
\textbf{Financial inclusion} = accès effectif aux services financiers formels (compte, crédit, assurance, paiement) à coût abordable. Au Cameroun, le taux de bancarisation était d'environ \textbf{35\%} en 2023 (estimations BEAC/INS). La \textbf{Stratégie Nationale d'Inclusion Financière (SNIF)} 2023--2030 vise \textbf{65\%} d'adultes bancarisés.
\end{definition}

\begin{center}
\begin{tabularx}{\linewidth}{|l|X|}\hline
\rowcolor{popblueL}
\textbf{English Term} & \textbf{French Gloss / Définition} \\ \hline
\cle{Account opening} & Ouverture de compte \\ \hline
\cle{Deposit} / \cle{Withdrawal} & Dépôt / Retrait \\ \hline
\cle{Interest rate} & Taux d'intérêt (épargne ou crédit) \\ \hline
\cle{Loan} / \cle{Credit} & Prêt / Crédit \\ \hline
\cle{Collateral} / \cle{Guarantee} & Garantie (bien mis en gage) \\ \hline
\cle{Overdraft} / \cle{Overdraft facility} & Découvert (autorisé) \\ \hline
\cle{Standing order} & Virement permanent (ordre permanent) \\ \hline
\cle{Bank statement} & Relevé de compte \\ \hline
\cle{KYC (Know Your Customer)} & Connaissance du client (procédure d'identification) \\ \hline
\cle{Financial inclusion} & Inclusion financière \\ \hline
\cle{Mobile money} & Monnaie mobile (Orange Money, MTN MoMo) \\ \hline
\cle{Microcredit} & Microcrédit \\ \hline
\cle{Diaspora transfer / Remittance} & Transfert de la diaspora \\ \hline
\cle{Account maintenance fees} & Frais de tenue de compte \\ \hline
\cle{Borrower's insurance} & Assurance emprunteur \\ \hline
\cle{Late payment penalty} & Pénalité de retard \\ \hline
\cle{Usury rate / Cap rate} & Taux d'usure (plafond légal) \\ \hline
\end{tabularx}
\end{center}

\courssub{Grammar Focus: Coordination -- Building Complex Sentences}
In professional writing (emails, reports, contracts), you must link ideas clearly. \cle{Coordination} joins two \emph{independent clauses} (complete sentences) of equal importance.

\begin{propriete}[Coordinating Conjunctions -- FANBOYS]
Use \textbf{FANBOYS} (\textbf{F}or, \textbf{A}nd, \textbf{N}or, \textbf{B}ut, \textbf{O}r, \textbf{Y}et, \textbf{S}o) to connect two independent clauses.
\begin{itemize}
\item \textbf{And} adds information: \emph{He opened an account \textbf{and} deposited money.}
\item \textbf{But} shows contrast: \emph{Rates are low, \textbf{but} the service is fast.}
\item \textbf{Or} gives an alternative: \emph{You can visit the branch \textbf{or} use the app.}
\item \textbf{So} shows result: \emph{She forgot her PIN, \textbf{so} she reset it.}
\end{itemize}
\emph{Punctuation rule:} Place a \textbf{comma} before the conjunction if both clauses have \textbf{different subjects}. If the subject is the same, \textbf{no comma} is needed.
\end{propriete}

\begin{propriete}[Conjunctive Adverbs -- \emph{However, Therefore, Moreover, Furthermore, Consequently}]
These are \textbf{not} conjunctions. They connect two independent clauses with a \textbf{semicolon (;)} before and a \textbf{comma (,)} after.
\[ \text{Clause 1 ; \textbf{however / therefore / moreover}, Clause 2 .} \]
\begin{itemize}
\item \emph{Interest rates are high; \textbf{however}, microfinance helps small businesses.}
\item \emph{The KYC process is strict; \textbf{therefore}, fraud is reduced.}
\item \emph{Mobile money is popular; \textbf{moreover}, it reaches rural areas.}
\end{itemize}
\end{propriete}

\begin{methode}[Coordonner deux propositions avec un adverbe de liaison]
\emph{Objectif :} Écrire une phrase complexe correcte montrant contraste, conséquence ou addition.
\begin{enumerate}
\item Identifier les deux phrases indépendantes (sujet + verbe chacune).
\item Choisir l'adverbe selon la logique : \textbf{however} (contraste), \textbf{therefore} (conséquence), \textbf{moreover} (addition).
\item Ponctuer : point-virgule + adverbe + virgule.
\item Vérifier : chaque partie a son sujet et son verbe conjugué.
\end{enumerate}
\end{methode}

\begin{attention}[Pièges fréquents]
\begin{itemize}
\item \textbf{Confusion \emph{Interest rate} vs \emph{Exchange rate}} : \emph{Interest rate} = taux d'intérêt (épargne/crédit). \emph{Exchange rate} = taux de change (devise). Ne pas confondre.
\item \textbf{Virgule manquante avec FANBOYS} : \emph{He saved money, and he bought a car.} (Sujets différents $\rightarrow$ virgule avant \emph{and}.) \emph{He saved money and bought a car.} (Même sujet $\rightarrow$ pas de virgule.)
\item \textbf{Point-virgule oublié} : \emph{Rates are high, however microfinance helps.} $\rightarrow$ \textbf{Incorrect}. Correct : \emph{Rates are high; however, microfinance helps.}
\end{itemize}
\end{attention}

\begin{exoresolu}[Coordination Practice]
\textbf{Combine each pair into one sentence using the word in brackets. Punctuate correctly.}
\begin{enumerate}
\item The bank approved the loan. The interest rate is 12\%. (\emph{but})
\item You must provide your CNI. You must provide a proof of address. (\emph{and})
\item Mobile money is convenient. It has transaction limits. (\emph{however})
\item She forgot her password. She could not access her account. (\emph{so})
\end{enumerate}
\tcblower
\textbf{Solutions :}
\begin{enumerate}
\item The bank approved the loan, \textbf{but} the interest rate is 12\%. (Two subjects $\rightarrow$ comma)
\item You must provide your CNI \textbf{and} a proof of address. (Same subject $\rightarrow$ no comma, ellipsis of the second subject.) \textit{Full form also correct:} You must provide your CNI, \textbf{and} you must provide a proof of address.
\item Mobile money is convenient\textbf{; however,} it has transaction limits. (Semicolon + comma)
\item She forgot her password\textbf{, so} she could not access her account. (Two subjects $\rightarrow$ comma)
\end{enumerate}
\end{exoresolu}

\courssub{Analysing Simplified Contract Clauses}
Before signing, a smart consumer reads the \textbf{Terms \& Conditions}. Here are key clauses to identify.

\begin{exemplebox}[Key Clauses in a \emph{Compte Épargne} Contract -- UBA Cameroon (Simplified)]
\begin{enumerate}
\item \textbf{Article 4 -- Frais de tenue de compte :} \num{1200} FCFA par trimestre (déduits automatiquement). Exonération si solde moyen mensuel $>$ \num{50000} FCFA.
\item \textbf{Article 7 -- Taux d'intérêt (épargne) :} \num{2,75}\% brut par an, capitalisé au 31/12. Prélèvement libératoire de \num{16,5}\% sur les intérêts.
\item \textbf{Article 12 -- Pénalités de retard (crédit) :} En cas d'impayé, majoration de \num{2} points du taux nominal + frais de dossier \num{5000} FCFA.
\item \textbf{Article 15 -- Assurance emprunteur (obligatoire pour crédit $>$ \num{500000} FCFA) :} Prime de \num{0,5}\% du capital restant dû par an. Couvre décès et invalidité permanente.
\item \textbf{Article 20 -- Clôture :} Préavis de 30 jours. Frais de clôture de \num{3000} FCFA si le compte a moins d'un an.
\end{enumerate}
\end{exemplebox}

\begin{exercice}[Contract Analysis]
\textbf{Read the clauses above and answer in full sentences.}
\begin{enumerate}
\item How much are the quarterly maintenance fees? When are they waived?
\item What is the net annual interest rate after the \num{16,5}\% tax on a \num{2,75}\% gross rate? (Calculate.)
\item If a borrower misses a payment on a loan at a \num{10}\% nominal rate, what is the new penalty rate?
\item Is borrower's insurance mandatory for a loan of \num{300000} FCFA? Why?
\end{enumerate}
\end{exercice}

\begin{corrige}[Exercice Contract Analysis]
\begin{enumerate}
\item Quarterly fees are \num{1200} FCFA. They are waived if the average monthly balance exceeds \num{50000} FCFA.
\item Gross rate = \num{2,75}\%. Tax = \num{16,5}\% of the interest. Net rate = $2{,}75 \times (1 - 0{,}165) = 2{,}75 \times 0{,}835 \approx \textbf{2{,}30\%}$.
\item Penalty rate = nominal rate + 2 points = \num{10}\% + \num{2}\% = \textbf{\num{12}\%} (plus \num{5000} FCFA fees).
\item No. Insurance is mandatory only for loans above \num{500000} FCFA. \num{300000} FCFA is below the threshold.
\end{enumerate}
\end{corrige}

\courssub{Case Study: \emph{Compte Jeune Entrepreneur} -- Comparing 3 Banks}
You are a young technician starting a small business (e.g., welding, IT services, agri-processing). You need a professional account. Compare these simplified offers (2024 data).

\begin{center}
\small
\begin{tabularx}{\linewidth}{|l|X|X|X|}\hline
\rowcolor{popblueL}
\textbf{Critère} & \textbf{BICEC -- \emph{Jeune Pro}} & \textbf{Afriland -- \emph{Start'Up}} & \textbf{UBA -- \emph{Young Biz}} \\ \hline
Capital social minimum & \num{100000} FCFA & \num{50000} FCFA & \num{100000} FCFA \\ \hline
Garanties exigées & Caution personnelle + domiciliation salaire & Caution solidaire (2 garants) ou nantissement matériel & Domiciliation chiffre d'affaires (min \num{200000}/mois) \\ \hline
Délai d'ouverture & 48 h (dossier complet) & 24 h (digital onboarding) & 72 h (validation risque) \\ \hline
Frais tenue compte / an & \num{6000} FCFA (gratuits la 1\iere{} année) & \num{4800} FCFA & \num{7200} FCFA \\ \hline
Plafond virement journalier & \num{5000000} FCFA & \num{3000000} FCFA & \num{10000000} FCFA \\ \hline
Accès mobile money & Oui (Orange/MTN) & Oui (Orange/MTN + AfriMoney) & Oui (UBA Moni / Orange/MTN) \\ \hline
Microcrédit associé & Jusqu'à 5 M FCFA (taux \num{14}\%) & Jusqu'à 3 M FCFA (taux \num{12}\%) & Jusqu'à 10 M FCFA (taux \num{15}\%) \\ \hline
\end{tabularx}
\end{center}

\begin{exemplebox}[Oral Interaction -- Role Play]
\textbf{Student A (Client):} Ask about opening conditions, fees, mobile money link, loan access.
\textbf{Student B (Bank Advisor):} Present your bank's offer using comparatives (\emph{lower than, higher than, the best}) and coordination (\emph{moreover, however}).
\textit{Example:} \emph{``Our fees are lower than BICEC's; moreover, we open accounts in 24 hours. However, the daily transfer limit is lower.''}
\end{exemplebox}

\courssub{Reading Comprehension: Banking in Cameroon (2023)}
Read the text and answer the questions.

\begin{exemplebox}[Article: \emph{``Cameroun : la bancarisation stagne autour de 35\%, le mobile money tire la croissance''} -- Adapted press article, 2023]
Despite a dynamic microfinance sector and the explosion of \textbf{mobile money} (over \num{10} million active wallets), Cameroon's formal \textbf{banking rate} remains stuck around \textbf{35\%} of adults (estimates BEAC/INS 2023). The \textbf{Littoral} and \textbf{Centre} regions concentrate about \num{65}\% of bank branches, leaving the \textbf{North} and \textbf{Far North} severely underserved.

Commercial banks (BICEC, Afriland, UBA, Ecobank, SCB) hold the large majority of total deposits but serve mostly corporate clients and civil servants. \textbf{Interest rate spreads} (lending rate minus deposit rate) average \num{10}--\num{12} points, among the highest in the CEMAC zone, discouraging SME borrowing. The \textbf{usury cap} (taux d'usure) is set at \num{18}\% for consumer loans (BEAC regulation).

The \textbf{National Financial Inclusion Strategy (SNIF)} 2023--2030 targets \textbf{65\%} adult inclusion by 2030. Key levers: \textbf{agent banking}, \textbf{digital onboarding} (e-KYC), \textbf{interoperability} between banks and mobile money operators, and \textbf{financial education} in technical schools.

\textit{Quote:} ``Mobile money is not a competitor; it is a gateway. The future is hybrid,'' says a Cameroonian bank CEO.
\end{exemplebox}

\begin{exercice}[Reading Comprehension -- MCQ \& Open Questions]
\textbf{A. Multiple Choice (choose the best answer):}
\begin{enumerate}
\item According to the text, the banking rate in Cameroon in 2023 was approximately:
  \begin{itemize}
  \item[a)] 18\% \item[b)] 35\% \item[c)] 65\% \item[d)] 85\%
  \end{itemize}
\item Which two regions concentrate about 65\% of bank branches?
  \begin{itemize}
  \item[a)] North \& Far North \item[b)] West \& South-West \item[c)] Littoral \& Centre \item[d)] East \& Adamawa
  \end{itemize}
\item What is the current usury cap (taux d'usure) for consumer loans in CEMAC?
  \begin{itemize}
  \item[a)] 10\% \item[b)] 12\% \item[c)] 15\% \item[d)] 18\%
  \end{itemize}
\item The SNIF 2023--2030 targets 65\% inclusion by 2030. Which is \textbf{NOT} a key lever mentioned?
  \begin{itemize}
  \item[a)] Agent banking \item[b)] Digital onboarding (e-KYC) \item[c)] Closing rural branches \item[d)] Interoperability bank/Mobile Money
  \end{itemize}
\end{enumerate}

\textbf{B. Open Questions (answer in full sentences, use coordination):}
\begin{enumerate}
\item Explain why high interest rate spreads discourage SME borrowing. Use \emph{therefore} or \emph{consequently}.
\item The CEO says ``Mobile money is a gateway.'' What does he mean? Use \emph{moreover} or \emph{furthermore} to add a reason.
\item Compare the banking situation in Littoral/Centre vs North/Far North. Use \emph{however} or \emph{whereas}.
\end{enumerate}
\end{exercice}

\begin{corrige}[Exercice Reading Comprehension]
\textbf{A. MCQ:} 1--b) 35\% \quad 2--c) Littoral \& Centre \quad 3--d) 18\% \quad 4--c) Closing rural branches

\textbf{B. Open Questions (sample answers):}
\begin{enumerate}
\item Interest rate spreads average 10--12 points; \textbf{therefore}, the cost of credit is too high for small businesses to invest profitably.
\item He means mobile money brings unbanked people into the digital financial system; \textbf{moreover}, it builds a transaction history that banks can later use to offer formal credit.
\item Littoral and Centre have about 65\% of branches; \textbf{however}, the North and Far North are severely underserved, \textbf{whereas} they have high mobile money adoption.
\end{enumerate}
\end{corrige}

\courssub{Illustration 1: Flowchart -- Account Opening Process in Cameroon}
\begin{popfigure}
\begin{tikzpicture}[
  node distance=1.2cm and 1.5cm,
  box/.style={rectangle, draw=popblue, thick, fill=popblue!10, rounded corners, minimum width=2.8cm, minimum height=1.1cm, align=center, font=\small},
  arrow/.style={-Stealth, thick, popdark},
  start/.style={ellipse, draw=popgreen, thick, fill=popgreen!15, minimum width=2.2cm, minimum height=1cm, align=center, font=\small},
  decision/.style={diamond, draw=poporange, thick, fill=poporange!15, aspect=1.5, minimum width=2.5cm, minimum height=1.2cm, align=center, font=\small}
]
\node[start] (start) {Début : projet d'ouverture};
\node[box, below=of start, align=center] (pieces){1. Rassembler pièces KYC\\CNI, justificatif domicile,\\photos, (extrait K-bis si entreprise)};
\node[box, below=of pieces, align=center] (rdv){2. Prendre rendez-vous\\Agence ou en ligne (e-KYC)};
\node[decision, below=of rdv] (check) {Dossier complet ?};
\node[box, below=of check, align=center] (signature){3. Signature contrat\\+ convention de compte};
\node[box, below=of signature, align=center] (depot){4. Premier dépôt\\Minimum requis (ex : 5 000 FCFA)};
\node[box, below=of depot, align=center] (activation){5. Activation compte\\Création codes / accès en ligne};
\node[box, below=of activation, align=center] (delivery){6. Remise carte/chéquier\\+ codes SMS/Email};
\node[start, below=of delivery] (end) {Fin : compte opérationnel};

\draw[arrow] (start) -- (pieces);
\draw[arrow] (pieces) -- (rdv);
\draw[arrow] (rdv) -- (check);
\draw[arrow] (check) -- node[left, font=\footnotesize, text=popgreen] {Oui} (signature);
\draw[arrow] (check.east) -- ++(1.5,0) |- node[near start, above, font=\footnotesize, text=poporange] {Non} (rdv.east);
\draw[arrow] (signature) -- (depot);
\draw[arrow] (depot) -- (activation);
\draw[arrow] (activation) -- (delivery);
\draw[arrow] (delivery) -- (end);
\end{tikzpicture}
\legende{Processus standard d'ouverture de compte bancaire au Cameroun (basé sur les procédures BICEC/Afriland/UBA 2024). La boucle « Non » renvoie au rendez-vous pour compléter le dossier.}
\end{popfigure}

\courssub{Illustration 2: Banking Rate by Region (2018--2023) -- pgfplots}
\begin{popfigure}
\begin{tikzpicture}
\begin{axis}[
  width=\linewidth, height=9cm,
  xlabel={Année}, ylabel={Taux de bancarisation (\%)},
  xmin=2017.5, xmax=2023.5,
  ymin=0, ymax=60,
  xtick={2018,2019,2020,2021,2022,2023},
  xticklabels={2018,2019,2020,2021,2022,2023},
  ytick={0,10,20,30,40,50,60},
  legend style={at={(0.02,0.98)}, anchor=north west, font=\small, fill=popcream, draw=popline},
  grid=major, grid style={popline!50},
  title style={font=\normalsize\bfseries, text=popdark},
  title={Évolution du taux de bancarisation par région (estimations BEAC/INS)}
]
\addplot[popcurve, mark=*, color=popblue] coordinates {(2018,28) (2019,30) (2020,31) (2021,32) (2022,34) (2023,36)}; \addlegendentry{Centre (Yaoundé)}
\addplot[popcurve, mark=square*, color=poporange] coordinates {(2018,30) (2019,32) (2020,33) (2021,35) (2022,37) (2023,39)}; \addlegendentry{Littoral (Douala)}
\addplot[popcurve, mark=triangle*, color=popgreen] coordinates {(2018,22) (2019,24) (2020,25) (2021,26) (2022,28) (2023,30)}; \addlegendentry{Ouest (Bafoussam)}
\addplot[popcurve, mark=diamond*, color=poppurple] coordinates {(2018,12) (2019,13) (2020,14) (2021,15) (2022,16) (2023,17)}; \addlegendentry{Nord / Extrême-Nord (Garoua, Maroua)}
\end{axis}
\end{tikzpicture}
\legende{Taux de bancarisation (adultes avec compte formel) par région, 2018--2023. Source : estimations BEAC/INS (valeurs indicatives). Note : le Nord et l'Extrême-Nord restent très en retrait ; le mobile money y comble partiellement l'écart.}
\end{popfigure}

\courssub{Illustration 3: Comparative Table -- Fees, Savings Rates, Mobile Money Limits}
\begin{center}
\small
\begin{tabularx}{\linewidth}{|l|c|c|c|X|}\hline
\rowcolor{popblueL}
\textbf{Institution} & \textbf{Frais annuels (FCFA)} & \textbf{Taux épargne net (\%/an)} & \textbf{Plafond Mobile Money (jour)} & \textbf{Note} \\ \hline
BICEC & 6 000 (gratuits 1\iere{} année) & 2,3\% & 1 000 000 (Orange/MTN) & Réseau d'agences dense \\ \hline
Afriland First Bank & 4 800 & 2,5\% & 1 500 000 (AfriMoney + Orange/MTN) & Digital onboarding rapide \\ \hline
UBA Cameroon & 7 200 & 2,0\% & 2 000 000 (UBA Moni + Orange/MTN) & Plafond de virement élevé \\ \hline
Ecobank & 6 500 & 2,2\% & 1 000 000 (Ecobank Mobile + Orange/MTN) & Réseau panafricain \\ \hline
SCB Cameroun & 5 500 & 2,4\% & 1 000 000 (Orange/MTN) & Focus corporate \\ \hline
\rowcolor{poporangeL}
\textbf{Orange Money} & 0 (frais par transaction) & 0\% (pas d'intérêt) & \textbf{1 000 000} (retrait) / \textbf{5 000 000} (dépôt) & Leader de part de marché \\ \hline
\rowcolor{poporangeL}
\textbf{MTN MoMo} & 0 (frais par transaction) & 0\% (pas d'intérêt) & \textbf{1 000 000} (retrait) / \textbf{5 000 000} (dépôt) & Couverture rurale forte \\ \hline
\end{tabularx}
\end{center}
\legende{Comparatif simplifié 2024 : 5 banques commerciales + 2 opérateurs Mobile Money. Frais annuels = tenue de compte. Taux épargne net = après prélèvement libératoire de 16,5\%. Plafonds Mobile Money = limites journalières réglementaires BEAC (valeurs indicatives).}

\courssub{Exemple Chiffré : Taux Directeur, Usure \& Épargne}
\begin{exemplebox}[Repères Chiffrés Clés 2024 (BEAC / Marché)]
\begin{itemize}
\item \textbf{Taux directeur BEAC (TIAO) :} environ \num{5,00}\% (taux d'injection de liquidité).
\item \textbf{Taux d'usure (plafond légal crédits conso) :} \textbf{\num{18}\%} (règlement CEMAC). Tout taux supérieur à 18\% = usure (illégal).
\item \textbf{Taux crédit conso typique banques :} \num{14}--\num{16}\% (sous le plafond).
\item \textbf{Taux crédit PME :} \num{10}--\num{13}\% (selon garanties).
\item \textbf{Épargne livret (compte sur livret) :} \num{2}--\num{3}\% net/an (après impôt de \num{16,5}\%).
\item \textbf{Épargne à terme (12 mois) :} \num{4}--\num{5,5}\% brut (négociable au-delà de 10 M FCFA).
\item \textbf{Mobile Money :} 0\% d'intérêt (porte-monnaie électronique, pas un compte d'épargne).
\end{itemize}
\end{exemplebox}

\courssub{Le Saviez-Vous ? -- National Financial Inclusion Strategy}
\begin{savaistu}[Stratégie Nationale d'Inclusion Financière (SNIF) 2023--2030]
Le Cameroun a lancé en 2023 la \textbf{SNIF} avec l'appui de la BEAC et de la Banque Mondiale. Objectif : passer d'environ \textbf{35\%} (2023) à \textbf{65\%} d'adultes ayant accès à des services financiers formels d'ici 2030. Quatre piliers :
\begin{enumerate}
\item \textbf{Accès :} agent banking, e-KYC, interopérabilité Banque/Mobile Money.
\item \textbf{Usage :} paiement numérique des salaires, subventions, factures (Eneo, Camwater).
\item \textbf{Qualité :} protection du consommateur, éducation financière (dont dans l'enseignement technique).
\item \textbf{Écosystème :} fintech, crédit numérique, assurance indicielle agricole.
\end{enumerate}
\textit{Impact pour vous :} en Terminale technique, vous êtes la cible prioritaire de l'éducation financière. Maîtriser ces termes en anglais = atout professionnel majeur.
\end{savaistu}

\courssub{Integration Exercise: Writing a Formal Email to a Bank Advisor}
\begin{exercice}[Writing Task -- Formal Email]
\textbf{Context:} You are a Terminale student in \emph{Maintenance Industrielle}. You have saved \num{150000} FCFA from holiday jobs. You want to open a \emph{Compte Épargne Jeune} at Afriland First Bank. Write a formal email (120--150 words) to the bank advisor requesting an appointment. Use:
\begin{itemize}
\item At least \textbf{2 coordinating conjunctions (FANBOYS)}.
\item At least \textbf{2 conjunctive adverbs (however, therefore, moreover)} with correct punctuation (; ,).
\item Key vocabulary: \emph{KYC, deposit, interest rate, standing order, mobile money link}.
\end{itemize}
\textit{Structure:} Subject line -- Salutation -- Purpose -- Questions (fees, rate, mobile link, documents) -- Proposed dates -- Closing.
\end{exercice}

\begin{corrige}[Exercice Writing Task -- Sample Answer]
\textbf{Subject:} Appointment Request -- Youth Savings Account Opening (Compte Épargne Jeune)

Dear Sir/Madam,

I am a final-year technical student in Industrial Maintenance. I wish to open a Youth Savings Account at your Douala Akwa branch; \textbf{moreover}, I would like to link it to my Orange Money wallet. I have saved 150,000 FCFA for the initial \textbf{deposit}.

Could you please confirm the current \textbf{interest rate} and the annual maintenance fees? \textbf{Furthermore}, what \textbf{KYC} documents are required for a student? I plan to set up a \textbf{standing order} of 10,000 FCFA monthly from my holiday job earnings.

I am available next Tuesday 14th or Thursday 16th after 3 pm. \textbf{However}, I can adapt to your schedule.

I look forward to your reply.

Yours faithfully,\\
\textit{Ngono Jean-Pierre}\\
Tel: 6xx xx xx xx\\
Email: jp.ngono@email.com
\end{corrige}

\courssub{Summary \& Key Takeaways}
\begin{aretenir}[Section 3 -- Banking Services]
\begin{itemize}
\item \textbf{Vocabulary:} Master the 15+ key terms (KYC, overdraft, standing order, collateral, usury rate, financial inclusion...).
\item \textbf{Reading:} Identify key contract clauses (fees, interest, penalties, insurance, closure). Calculate net rates.
\item \textbf{Grammar -- Coordination:} FANBOYS (comma if different subjects). Conjunctive adverbs (\emph{however, therefore, moreover}) $\rightarrow$ semicolon + adverb + comma.
\item \textbf{Case Study:} Compare \emph{Jeune Entrepreneur} offers: capital, guarantees, delay, fees, mobile link, microcredit access.
\item \textbf{Data:} Banking rate around 35\% (2023), target 65\% (2030). Regional gaps (Centre/Littoral vs North/Far North). BEAC usury cap 18\%.
\item \textbf{Writing:} Formal email structure + coordination connectors for professional clarity.
\end{itemize}
\end{aretenir}

\coursec{Complex Sentences: Subordination \& Expressing Preferences}

\noindent In the previous section (\emph{3. Banking Services in Cameroon}), we read texts about opening accounts, applying for loans, and using mobile money. We noticed that professional banking documents rarely use simple, isolated sentences. Instead, they use \textbf{complex sentences} to show precise logical relationships: cause, condition, contrast, and time. Mastering \cle{subordination} and \cle{expressing preferences} will allow you to write formal emails, loan applications, and business reports that sound professional and convincing.

\courssub{Understanding Subordination: Adverbial \& Relative Clauses}

\begin{definition}[Subordination]
\cle{Subordination} is a grammatical link where a \textbf{dependent clause (subordinate clause)} cannot stand alone and adds information to a \textbf{main clause (independent clause)}. The subordinate clause is introduced by a \cle{subordinating conjunction} (e.g., \textit{because, although, if, when}) or a \cle{relative pronoun} (e.g., \textit{who, which, where, whose}).
\begin{center}
\textit{Main Clause} + \textit{Subordinator} + \textit{Subordinate Clause}
\end{center}
\textbf{Example (Banking context):} \textit{Because interest rates are high (Subordinate Clause of Cause), I prefer a savings account (Main Clause).}
\end{definition}

\textbf{A. Adverbial Clauses (Propositions subordonnées circonstancielles)}\par
These clauses function like adverbs: they answer \textit{Why?}, \textit{When?}, \textit{Under what condition?}, \textit{Despite what?}.

\begin{center}
\begin{tabularx}{\linewidth}{|l|X|X|}\hline
\rowcolor{popblueL}
\textbf{Function} & \textbf{Key Conjunctions} & \textbf{Banking Example} \\ \hline
\textbf{Cause / Reason} & \cle{because}, \cle{since}, \cle{as} & \textit{Since the TEG exceeds the usury rate, the contract is illegal.} \\ \hline
\textbf{Concession / Contrast} & \cle{although}, \cle{even though}, \cle{while} & \textit{Although the fees are high, the service is reliable.} \\ \hline
\textbf{Condition} & \cle{if}, \cle{unless} (= if not), \cle{provided that} & \textit{Unless you provide collateral, the bank will reject the loan.} \\ \hline
\textbf{Time} & \cle{when}, \cle{while}, \cle{before}, \cle{after}, \cle{as soon as} & \textit{As soon as you receive the SMS, confirm the transaction.} \\ \hline
\end{tabularx}
\end{center}

\begin{attention}[The ``Although / But'' Trap]
\cle{Never} use \textbf{although / even though} and \textbf{but} in the same sentence. They express the same logical relation (concession).
\begin{itemize}
    \item \textbf{Incorrect:} \textit{Although the queue was long, but I waited.}
    \item \textbf{Correct:} \textit{Although the queue was long, I waited.} \quad OR \quad \textit{The queue was long, but I waited.}
\end{itemize}
\end{attention}

\textbf{B. Relative Clauses (Propositions relatives)}\par
These clauses modify a noun (the \cle{antecedent}). They are introduced by relative pronouns: \textit{who} (people), \textit{which} (things), \textit{that} (people or things, defining only), \textit{where} (places), \textit{whose} (possession).

\begin{propriete}[Defining vs. Non-Defining Relative Clauses]
\textbf{1. Defining (Restrictive) --- No commas.} Essential information. Identifies \textit{which one}.
\begin{itemize}
    \item \textit{The client \textbf{who} signed the cheque is waiting.} (Which client? The one who signed).
    \item \textit{This is the account \textbf{which/that} offers 5\,\% interest.}
\end{itemize}
\textbf{2. Non-Defining (Non-Restrictive) --- Commas required.} Extra, non-essential information.
\begin{itemize}
    \item \textit{My bank manager, \textbf{who} knows me well, approved the loan instantly.}
    \item \textit{The BEAC headquarters, \textbf{which} is in Yaoundé, regulates the usury rate.}
\end{itemize}
\textbf{Note:} In non-defining clauses, you \textbf{cannot} use \textit{that}; you must use \textit{who/which}.
\end{propriete}

\begin{popfigure}
\begin{tikzpicture}[
    level distance=1.8cm,
    sibling distance=3.5cm,
    every node/.style={font=\small, align=center},
    edge from parent/.style={draw=popdark, thick, -{Stealth[length=2mm]}},
    main/.style={rectangle, draw=popblue, thick, fill=popblueL, rounded corners=3pt, minimum width=2.5cm},
    sub/.style={rectangle, draw=poporange, thick, fill=poporangeL, rounded corners=3pt, minimum width=2.2cm},
    rel/.style={rectangle, draw=popgreen, thick, fill=popgreenL, rounded corners=3pt, minimum width=2.2cm},
    conj/.style={circle, draw=poppurple, fill=poppurpleL, minimum size=6mm, font=\scriptsize\bfseries}
]
\node[main, align=center]{Main Clause\\ \textit{I will accept the loan}}
    child { node[conj] (c1) {if}
        child { node[sub, align=center]{Adverbial Clause\\ \textit{(Condition)}\\ \textit{the rate is below 12\,\%}} }
    }
    child { node[conj] (c2) {because}
        child { node[sub, align=center]{Adverbial Clause\\ \textit{(Cause)}\\ \textit{I need capital urgently}} }
    }
    child { node[rel, align=center]{Relative Clause\\ \textit{(Defining)}\\ \textit{which the agent offered}} };
\draw[dashed, popmuted] (c1) -- (c2);
\end{tikzpicture}
\legende{Syntax Tree: Analysis of a complex banking sentence. \textit{Main Clause: ``I will accept the loan''}. 1. Adverbial Clause of Condition (\textit{if the rate is below 12\,\%}). 2. Adverbial Clause of Cause (\textit{because I need capital urgently}). 3. Defining Relative Clause (\textit{which the agent offered}) modifying ``the loan''.}
\end{popfigure}

\begin{exoresolu}[Combining Simple Sentences into Complex Ones]
\textbf{Task:} Combine the following simple sentences into ONE complex sentence using the connector in brackets. Identify the type of subordinate clause created.

\begin{enumerate}
    \item The microfinance institution charges 30\,\%. The commercial bank charges 14\,\%. \textit{(while / concession)}
    \item You want a mortgage. You must prove stable income. \textit{(if / condition)}
    \item The customer closed his account. The fees were too high. \textit{(because / cause)}
    \item The ATM swallowed my card. I entered the wrong PIN three times. \textit{(after / time)}
    \item The new agency opened in Douala. It offers Islamic banking. \textit{(which / relative)}
\end{enumerate}

\tcblower
\textbf{Solutions:}
\begin{enumerate}
    \item \textit{While the microfinance institution charges 30\,\%, the commercial bank charges 14\,\%.} (Adverbial Clause of Concession/Contrast)
    \item \textit{If you want a mortgage, you must prove stable income.} (Adverbial Clause of Condition)
    \item \textit{The customer closed his account because the fees were too high.} (Adverbial Clause of Cause)
    \item \textit{The ATM swallowed my card after I entered the wrong PIN three times.} (Adverbial Clause of Time)
    \item \textit{The new agency, which opened in Douala, offers Islamic banking.} (Non-Defining Relative Clause --- note the commas) \quad OR \quad \textit{The new agency which opened in Douala offers Islamic banking.} (Defining --- implies there are other agencies elsewhere).
\end{enumerate}
\end{exoresolu}

\courssub{Expressing Preferences \& Strong Advice: \textit{Would Rather, Prefer, Had Better}}

In a professional banking context, you often need to advise a client or state your own choice clearly. Three structures are essential.

\begin{methode}[Choosing the Right Structure for Preferences \& Advice]
\textbf{Step 1: Identify the Function.}
\begin{itemize}
    \item \textbf{Specific Preference (Now/Future):} \cle{Would rather / Would sooner} + \textbf{Bare Infinitive} (no \textit{to}).
    \item \textbf{General Preference (Habit/Taste):} \cle{Prefer} + \textbf{Gerund (-ing)} \textit{or} \cle{Prefer} + \textbf{Full Infinitive (to + V)}.
    \item \textbf{Strong Advice / Warning (Consequence if ignored):} \cle{Had better} + \textbf{Bare Infinitive}.
\end{itemize}

\textbf{Step 2: Check the Grammar Form.}
\begin{center}
\begin{tabular}{|l|c|c|c|}\hline
\rowcolor{popblueL}
\textbf{Structure} & \textbf{Form} & \textbf{Negative} & \textbf{Time Reference} \\ \hline
\textit{I'd rather} & + V (bare inf.) & \textit{I'd rather not V} & Present / Future (Specific) \\ \hline
\textit{I prefer} & + V-ing / to V & \textit{I prefer not V-ing} & General / Habitual \\ \hline
\textit{You'd better} & + V (bare inf.) & \textit{You'd better not V} & Immediate Future (Warning) \\ \hline
\end{tabular}
\end{center}

\textbf{Step 3: Compare Alternatives (Comparison).}
\begin{itemize}
    \item \textit{I'd rather open a savings account \textbf{than} invest in stocks.} (Comparing two specific actions now).
    \item \textit{I prefer saving \textbf{to} investing.} (General habit, Gerund).
    \item \textit{I prefer to save \textbf{rather than} (to) invest.} (Infinitive structure).
\end{itemize}
\end{methode}

\begin{popfigure}
\begin{tikzpicture}[
    every node/.style={font=\small},
    setA/.style={circle, draw=popblue, thick, fill=popblue!20, minimum width=5cm},
    setB/.style={circle, draw=poporange, thick, fill=poporange!20, minimum width=5cm},
    setC/.style={circle, draw=popgreen, thick, fill=popgreen!20, minimum width=5cm},
    labelStyle/.style={font=\bfseries\small, text=popdark}
]
\begin{scope}[blend mode=multiply]
    \node[setA] (A) at (0,0) {};
    \node[setB] (B) at (3,0) {};
    \node[setC] (C) at (1.5,2.2) {};
\end{scope}

% Labels for sets
\node[labelStyle, popblue] at (-2.2, 0) {Would Rather / Sooner};
\node[labelStyle, poporange] at (5.2, 0) {Prefer};
\node[labelStyle, popgreen] at (1.5, 4.2) {Had Better};

% Specific zones
\node[align=center, text width=3.5cm] at (-2.2, -1.5) {\textbf{Specific Choice}\\\textit{I'd rather save than spend.}\\\textit{Form: Bare Inf.}};
\node[align=center, text width=3.5cm] at (5.2, -1.5) {\textbf{General Taste}\\\textit{I prefer saving to spending.}\\\textit{Form: Gerund / to+V}};
\node[align=center, text width=3.5cm] at (1.5, 3.5) {\textbf{Strong Advice}\\\textit{You'd better check the TEG.}\\\textit{Form: Bare Inf.}};

% Intersections
\node[align=center, font=\tiny, text=popdark] at (0.5, 0.5) {Express\\Preference};
\node[align=center, font=\tiny, text=popdark] at (1.5, 1.2) {\textbf{All Three}\\\textit{Modal / Semi-modal}\\\textit{No `to' after 'd better / 'd rather}};
\node[align=center, font=\tiny, text=popdark] at (2.5, 0.5) {Compare\\Options};
\end{tikzpicture}
\legende{Venn Diagram: Comparison of \textit{Would Rather}, \textit{Prefer}, and \textit{Had Better}. Central overlap: all express preference or advice and are followed by a verb form without `to' (except \textit{prefer to V}). Left: Specific/Immediate. Right: General/Habitual. Top: Warning/Consequence.}
\end{popfigure}

\begin{exemplebox}[Banking Context Examples]
\begin{itemize}
    \item \textbf{Would rather (Specific):} \textit{``I'd \textbf{rather} keep my money in a term deposit \textbf{than} risk it on the stock market.''} (Decision for \textit{now}).
    \item \textbf{Prefer (General):} \textit{``Most Cameroonian SMEs \textbf{prefer} borrowing from microfinance \textbf{to} waiting for bank approval.''} (General observation).
    \item \textbf{Had better (Warning):} \textit{``You'd \textbf{better} read the fine print \textbf{before} signing the loan agreement.''} (Implied threat: \textit{otherwise you will lose money}).
\end{itemize}
\end{exemplebox}

\begin{attention}[Common Errors with Preferences]
\begin{enumerate}
    \item \textbf{Wrong:} \textit{I'd rather to save.} \quad \textbf{Right:} \textit{I'd rather save.} (Bare infinitive).
    \item \textbf{Wrong:} \textit{I prefer save than spend.} \quad \textbf{Right:} \textit{I prefer saving to spending} (Gerund) \textit{OR} \textit{I prefer to save rather than (to) spend}.
    \item \textbf{Wrong:} \textit{You had better to go.} \quad \textbf{Right:} \textit{You'd better go.} (\textit{Had} is a past form but the meaning is present/future; no \textit{to}).
    \item \textbf{Contraction:} Always use \textit{'d} in speech and informal writing: \textit{I'd rather, You'd better, He'd prefer}.
\end{enumerate}
\end{attention}

\begin{savaistu}[La BEAC et le Taux d'Usure]
La \textbf{BEAC} (Banque des États de l'Afrique Centrale), la banque centrale de la CEMAC dont le siège est à Yaoundé, encadre le \textbf{taux d'usure} (maximum legal interest rate), c'est-à-dire le taux maximal auquel une banque ou un établissement de microfinance peut prêter. Il est fixé par catégorie de crédit à partir du TEG (Taux Effectif Global) moyen pratiqué par les établissements, majoré d'une marge.
\begin{itemize}
    \item \textit{Exemple chiffré (pédagogique) :} Si le TEG moyen constaté pour les crédits aux PME est de 12\,\% et que la marge autorisée est d'un tiers, le taux d'usure = 12\,\% + 4\,\% = \textbf{16\,\%}.
    \item Une banque qui prête à 18\,\% dans cette catégorie commet un délit d'usure. \textit{``Unless the rate respects the usury ceiling, the contract is void.''}
\end{itemize}
\end{savaistu}

\courssub{Transformation Workshop: Informal Advice $\rightarrow$ Formal Recommendations}

In professional writing (reports, official emails, loan proposals), we often transform direct, informal structures (\textit{You should, You must, I think}) into \textbf{impersonal, passive, or subjunctive structures}. This is called \cle{hedging} or using a \cle{formal register}.

\begin{propriete}[Formal Transformation Patterns]
\begin{center}
\begin{tabularx}{\linewidth}{|l|X|X|}\hline
\rowcolor{popblueL}
\textbf{Informal / Direct} & \textbf{Formal / Impersonal Structure} & \textbf{Example (Banking)} \\ \hline
\textit{You should / must / have to} & \textit{It is recommended / essential / mandatory that + Subjunctive (Base form)} & \textit{It is \textbf{mandatory that the borrower provide} collateral.} \\ \hline
\textit{I think / We believe} & \textit{It is considered / believed / assumed that + Indicative} & \textit{It is \textbf{assumed that the rate will drop}.} \\ \hline
\textit{You need to} & \textit{Requirement: + Noun phrase / Gerund} & \textit{\textbf{Providing} a business plan is a \textbf{requirement}.} \\ \hline
\textit{If you don't...} & \textit{Failure to + V / Non-compliance with... results in...} & \textit{\textbf{Failure to repay} results in legal action.} \\ \hline
\end{tabularx}
\end{center}
\textbf{Note on the Subjunctive:} After \textit{It is essential/vital/mandatory/recommended that}, use the \textbf{base form} (no \textit{s}, no \textit{to}) for all persons. \textit{``It is essential that he \textbf{sign} (not signs) the cheque.''}
\end{propriete}

\begin{exoresolu}[Rewriting Informal Advice into Formal Recommendations]
\textbf{Task:} Rewrite the sentences using the formal structure indicated.

\begin{enumerate}
    \item \textit{Informal:} ``You should check your credit score before applying.'' \quad $\rightarrow$ \quad \textbf{Formal (It is recommended that...)}
    \item \textit{Informal:} ``You must repay the loan on time.'' \quad $\rightarrow$ \quad \textbf{Formal (It is mandatory that...)}
    \item \textit{Informal:} ``If you don't provide the documents, we will reject the file.'' \quad $\rightarrow$ \quad \textbf{Formal (Failure to...)}
    \item \textit{Informal:} ``I advise you to compare the TEG of different banks.'' \quad $\rightarrow$ \quad \textbf{Formal (It is advisable to... / Comparing... is advisable)}
    \item \textit{Informal:} ``The manager says you have to sign here.'' \quad $\rightarrow$ \quad \textbf{Formal (The manager requires that...)}
\end{enumerate}

\tcblower
\textbf{Solutions:}
\begin{enumerate}
    \item \textit{It is \textbf{recommended that the applicant check} his credit score before applying.}
    \item \textit{It is \textbf{mandatory that the borrower repay} the loan on time.}
    \item \textit{\textbf{Failure to provide} the documents \textbf{will result in} the rejection of the file.}
    \item \textit{\textbf{Comparing} the TEG of different banks \textbf{is advisable}.} \quad OR \quad \textit{It is \textbf{advisable to compare} the TEG...}
    \item \textit{The manager \textbf{requires that the client sign} here.}
\end{enumerate}
\end{exoresolu}

\courssub{Practice Exercises}

\begin{exercice}[Subordination \& Connectors]
\textbf{A. Complete the sentences with the correct conjunction: \textit{although, unless, since, as soon as, while, provided that}.}
\begin{enumerate}
    \item \_\_\_\_\_\_\_\_ the mobile money network is down, you cannot withdraw cash via USSD.
    \item \_\_\_\_\_\_\_\_ the interest rate is variable, the monthly payment remains fixed for the first year.
    \item You can get the preferential rate \_\_\_\_\_\_\_\_ you maintain a minimum balance of 500\,000 FCFA.
    \item \_\_\_\_\_\_\_\_ you receive the OTP (One Time Password), enter it immediately.
    \item Microfinance is faster, \_\_\_\_\_\_\_\_ commercial banks offer lower rates.
    \item \_\_\_\_\_\_\_\_ you have a guarantor, the microfinance institution won't lend you the money.
\end{enumerate}

\textbf{B. Join the pairs of sentences using a relative pronoun (\textit{who, which, where, whose}). Add commas if necessary.}
\begin{enumerate}
    \item The teller helped me. She speaks English and French.
    \item The branch is in Bonanjo. I opened my first account there.
    \item The client missed the deadline. His file was incomplete.
    \item The new ATM accepts Visa cards. It was installed yesterday.
\end{enumerate}
\end{exercice}

\begin{exercice}[Preferences \& Advice]
\textbf{C. Choose the correct form.}
\begin{enumerate}
    \item I'd rather \textbf{(keep / keeping / to keep)} my savings in FCFA than \textbf{(convert / converting / to convert)} to Euros right now.
    \item Generally, customers \textbf{(prefer / would rather)} digital banking \textbf{(to / than)} queuing in the hall.
    \item You'd better \textbf{(to notify / notify / notifying)} the bank immediately if you lose your card.
    \item The director would sooner \textbf{(resign / resigning / to resign)} \textbf{(than / to)} approve an illegal loan.
\end{enumerate}

\textbf{D. Transform the following informal tips into formal banking recommendations (use \textit{It is... that}, \textit{Failure to}, \textit{Gerund subjects}).}
\begin{enumerate}
    \item ``Don't share your PIN with anyone.''
    \item ``You must update your KYC (Know Your Customer) documents every year.''
    \item ``If you miss a payment, your credit score drops.''
    \item ``I suggest you use the bank's official app only.''
\end{enumerate}
\end{exercice}

\begin{corrige}[Exercice 4.4]
\textbf{A. Conjunctions}
\begin{enumerate}
    \item \textbf{Since / If} the mobile money network is down, you cannot withdraw cash via USSD. (Cause / Condition)
    \item \textbf{Although} the interest rate is variable, the monthly payment remains fixed for the first year. (Concession)
    \item You can get the preferential rate \textbf{provided that} you maintain a minimum balance of 500\,000 FCFA. (Condition)
    \item \textbf{As soon as} you receive the OTP, enter it immediately. (Time)
    \item Microfinance is faster, \textbf{while} commercial banks offer lower rates. (Contrast)
    \item \textbf{Unless} you have a guarantor, the microfinance institution won't lend you the money. (Negative Condition)
\end{enumerate}

\textbf{B. Relative Clauses}
\begin{enumerate}
    \item The teller \textbf{who} helped me speaks English and French. (Defining --- no commas)
    \item The branch \textbf{where} I opened my first account is in Bonanjo. (Defining --- no commas)
    \item The client \textbf{whose} file was incomplete missed the deadline. (Defining --- no commas)
    \item The new ATM, \textbf{which} was installed yesterday, accepts Visa cards. (Non-defining --- commas needed: extra information about the ATM).
\end{enumerate}

\textbf{C. Preferences \& Advice}
\begin{enumerate}
    \item \textbf{keep / convert} (\textit{Would rather + bare infinitive ... than + bare infinitive})
    \item \textbf{prefer / to} (\textit{General preference: Prefer + Gerund + to + Gerund})
    \item \textbf{notify} (\textit{Had better + bare infinitive})
    \item \textbf{resign / than} (\textit{Would sooner + bare infinitive ... than + bare infinitive})
\end{enumerate}

\textbf{D. Formal Transformations}
\begin{enumerate}
    \item \textit{\textbf{Sharing} your PIN with anyone \textbf{is strictly prohibited}.} \quad OR \quad \textit{It is \textbf{mandatory that clients not share} their PIN.}
    \item \textit{It is \textbf{mandatory that the client update} KYC documents annually.} \quad OR \quad \textit{\textbf{Updating} KYC documents annually \textbf{is a requirement}.}
    \item \textit{\textbf{Failure to make} a payment \textbf{results in} a drop in the credit score.}
    \item \textit{It is \textbf{advisable that clients use} only the bank's official app.} \quad OR \quad \textit{\textbf{Using} the bank's official app \textbf{is strongly recommended}.}
\end{enumerate}
\end{corrige}

\begin{aretenir}[Key Takeaways for Section 4]
\begin{itemize}
    \item \textbf{Subordination} creates hierarchy: Main Idea (Main Clause) + Supporting Detail (Subordinate Clause).
    \item \textbf{Adverbial Clauses:} Cause (\textit{because/since}), Concession (\textit{although}), Condition (\textit{if/unless/provided that}), Time (\textit{when/before/after/as soon as}).
    \item \textbf{Relative Clauses:} Defining (no commas, identifies) vs Non-defining (commas, extra information). Never use \textit{that} in non-defining clauses.
    \item \textbf{Preferences:} \textit{Would rather + V} (specific, now) vs \textit{Prefer + V-ing/to V} (general). Comparison: \textit{rather... than}, \textit{prefer... to}.
    \item \textbf{Strong Advice:} \textit{Had better + V} (warning of a bad consequence).
    \item \textbf{Formal Writing:} Transform \textit{You must/should} $\rightarrow$ \textit{It is mandatory/recommended that + Subjunctive (Base Verb)}.
\end{itemize}
\end{aretenir}

\coursec{Writing Workshop: Composing About Banking Services}

Building on the complex sentences and preference structures mastered in Section 4, we now move from \emph{sentence-level} accuracy to \emph{text-level} coherence. In professional banking contexts, a single grammatical error or a poorly structured paragraph can lead to a rejected loan application or an unresolved complaint. This workshop teaches you to craft complete, purpose-driven texts -- formal letters, complaint emails, analytical reports, and blog articles -- that meet the standards of Cameroonian financial institutions and the expectations of the \textsc{Probatoire}/\textsc{Baccalauréat} technique examination.

\courssub{Analysing Model Texts: Formal Banking Correspondence}

Before writing, we must read like writers. Professional banking communication in Cameroon follows strict conventions: a formal tone, precise vocabulary, logical organisation, and a clear \cle{call to action}. Let us dissect the anatomy of a formal banking letter.

\begin{definition}[Call to Action]
A \textbf{call to action (CTA)} is the final sentence of a formal letter or email that states exactly what you expect the recipient to do. It removes ambiguity and prompts a response.
\begin{itemize}
    \item \emph{Weak:} I hope you will help me.
    \item \textbf{Strong:} \emph{I kindly request that you schedule a meeting within the next ten working days to discuss the revision of my account maintenance fees.}
\end{itemize}
\end{definition}

\begin{popfigure}
\begin{tikzpicture}[
    box/.style={draw=popblue, thick, rounded corners=4pt, fill=popblue!5, minimum height=1.8cm, text width=0.92\linewidth, align=left, inner sep=8pt},
    arrow/.style={-Stealth, thick, popdark},
    label/.style={font=\small\bfseries, popdark, anchor=west}
]
% Header block
\node[box] (header) {
    \textbf{1. SENDER'S ADDRESS (Top Right)}\par
    \textit{Your Name / Business Name}\par
    \textit{P.O. Box 1234, Yaoundé, Cameroon}\par
    \textit{Tel: (+237) 6XX XX XX XX | Email: name@email.cm}
};
% Date
\node[box, below=0.3cm of header.south west, anchor=north west] (date) {
    \textbf{2. DATE (Below Sender, Right)}\par
    \textit{Yaoundé, 15 October 2025}
};
% Recipient
\node[box, below=0.3cm of date.south west, anchor=north west] (recipient) {
    \textbf{3. RECIPIENT'S ADDRESS (Left)}\par
    \textit{The Branch Manager}\par
    \textit{Afriland First Bank, Bastos Branch}\par
    \textit{P.O. Box 6789, Yaoundé}
};
% Subject
\node[box, below=0.3cm of recipient.south west, anchor=north west] (subject) {
    \textbf{4. OBJET / SUBJECT LINE (Bold, Centred)}\par
    \textbf{REQUEST FOR REDUCTION OF ACCOUNT MAINTENANCE FEES}\par
    \textit{Account No: 10002-0345678-98}
};
% Salutation
\node[box, below=0.3cm of subject.south west, anchor=north west] (salut) {
    \textbf{5. SALUTATION}\par
    \textit{Dear Sir / Madam,} \quad (\emph{never "Dear Friend" or "Hi"})
};
% Body Para 1
\node[box, below=0.3cm of salut.south west, anchor=north west] (p1) {
    \textbf{6. PARAGRAPH 1: CONTEXT \& PURPOSE}\par
    \textit{I am writing as a loyal customer of 5 years (Account \#10002-0345678-98) to formally request a review of the quarterly maintenance fees currently charged at 6,250 FCFA.}
};
% Body Para 2
\node[box, below=0.3cm of p1.south west, anchor=north west] (p2) {
    \textbf{7. PARAGRAPH 2: ARGUMENTATION (FACTS + COMPARISON)}\par
    \textit{Comparing your fee structure with competitors (e.g., UBA: 4,500 FCFA; BICEC: 5,000 FCFA), I find the current rate 25\% above the market average. Given my average monthly turnover of 2,000,000 FCFA and zero overdraft history, I believe a reduction is justified.}
};
% Body Para 3
\node[box, below=0.3cm of p2.south west, anchor=north west] (p3) {
    \textbf{8. PARAGRAPH 3: PREFERENCE + CALL TO ACTION}\par
    \textit{I would rather maintain my relationship with Afriland First Bank than switch providers. \textbf{I kindly request a reduction to 4,500 FCFA per quarter} and propose a meeting next week to finalise the amendment.}
};
% Closing
\node[box, below=0.3cm of p3.south west, anchor=north west] (close) {
    \textbf{9. FORMAL CLOSING (Long Formula)}\par
    \textit{I look forward to your favourable response.}\par
    \textit{Yours faithfully,}\par
    \vspace{0.5cm}
    \textit{(Signature)}\par
    \textit{John N. Doe}
};
% Attachments
\node[box, below=0.3cm of close.south west, anchor=north west] (attach) {
    \textbf{10. PIÈCES JOINTES / ENCLOSURES}\par
    \textit{Enc.: 1) Bank statements (last 6 months)}\par
    \textit{2) Fee comparison table (3 banks)}\par
    \textit{3) National ID copy}
};
% Arrows
\draw[arrow] (header.south) -- (date.north);
\draw[arrow] (date.south) -- (recipient.north);
\draw[arrow] (recipient.south) -- (subject.north);
\draw[arrow] (subject.south) -- (salut.north);
\draw[arrow] (salut.south) -- (p1.north);
\draw[arrow] (p1.south) -- (p2.north);
\draw[arrow] (p2.south) -- (p3.north);
\draw[arrow] (p3.south) -- (close.north);
\draw[arrow] (close.south) -- (attach.north);
% Side labels
\node[label, xshift=-2cm] at (header.west) {STRUCTURE:};
\node[label, xshift=-2cm] at (date.west) {};
\node[label, xshift=-2cm] at (recipient.west) {};
\node[label, xshift=-2cm] at (subject.west) {};
\node[label, xshift=-2cm] at (salut.west) {};
\node[label, xshift=-2cm] at (p1.west) {BODY:};
\node[label, xshift=-2cm] at (p2.west) {};
\node[label, xshift=-2cm] at (p3.west) {};
\node[label, xshift=-2cm] at (close.west) {END:};
\node[label, xshift=-2cm] at (attach.west) {};
\end{tikzpicture}
\legende{Squelette d'une lettre formelle bancaire camerounaise : 10 blocs obligatoires. Respectez l'alignement (expéditeur à droite, destinataire à gauche) et l'espacement.}
\end{popfigure}

\textbf{Four Model Text Types for the Exam.}\par
You must recognise and reproduce the distinct \emph{register} and \emph{structure} of each:

\begin{center}
\begin{tabularx}{\linewidth}{|l|X|X|X|}\hline
\textbf{Model Type} & \textbf{Purpose} & \textbf{Key Vocabulary} & \textbf{Grammar Focus} \\\hline
\textbf{Formal Loan Request Letter} & Persuade bank to grant credit & \textit{collateral (garantie), repayment schedule (échéancier), interest rate (taux d'intérêt), feasibility study (étude de faisabilité), cash flow (flux de trésorerie)} & Subordination (\textit{provided that, on condition that}), Passive (\textit{It is requested that...}) \\\hline
\textbf{Complaint Email (Undue Fees)} & Obtain refund / correction & \textit{discrepancy (écart), unauthorized debit (prélèvement non autorisé), statement (relevé), refund (remboursement), investigate (enquêter)} & Coordination (\textit{not only... but also}), Adverbs of time (\textit{recently, subsequently, previously}) \\\hline
\textbf{Bank Product Evaluation Report} & Compare \& recommend & \textit{pros/cons (avantages/inconvénients), benchmarking, yield (rendement), liquidity (liquidité), risk profile (profil de risque)} & Double comparison (\textit{the higher..., the lower...}), Equality (\textit{as competitive as}) \\\hline
\textbf{Blog Article: ``Choosing a Bank in Cameroon''} & Inform / Guide peers & \textit{mobile money (Mobile Money), branch network (réseau d'agences), digital banking (banque en ligne), KYC (Know Your Customer), financial inclusion (inclusion financière)} & Expressing preferences (\textit{I would recommend, It is advisable to}), Imperatives for advice \\\hline
\end{tabularx}
\end{center}

\begin{exemplebox}[Analysing a Complaint Email -- Structure Breakdown]
\textbf{Subject:} Dispute Regarding Unauthorised SMS Alert Charges -- Acct 10002-0345678-98\par
\textbf{Dear Customer Service Manager,}\par
\textbf{(Context)} I am writing to dispute a charge of 500 FCFA for ``SMS Alert Services'' debited on 03/10/2025. I never subscribed to this service.\par
\textbf{(Evidence)} My contract (Clause 4.2) states: ``Optional services require explicit written consent.'' My signed contract has no tick-box for SMS alerts. Furthermore, the charge appears \textbf{neither} on my initial fee schedule \textbf{nor} on the tariff guide displayed in the branch.\par
\textbf{(Action)} \textbf{I request an immediate refund of 500 FCFA and the permanent deactivation of this service.} I have attached a screenshot of the debit and the relevant contract page.\par
\textbf{(Closing)} I expect a resolution within 5 working days.\par
\textbf{Yours faithfully,}\par
\textbf{Marie K.}
\end{exemplebox}

\courssub{The Standard Structure of a Banking Composition}

Every high-scoring composition -- whether a letter, report, or essay -- follows a macro-structure that guides the reader logically from \emph{context} to \emph{action}.

\begin{propriete}[Cohésion Lexicale -- Lexical Chains]
\textbf{Cohésion lexicale} is the use of semantically related words across paragraphs to create texture and flow. In banking texts, activate the \textbf{Banking Lexical Chain}:
\[
\text{bank} \rightarrow \text{account} \rightarrow \text{deposit} \rightarrow \text{interest} \rightarrow \text{loan} \rightarrow \text{repayment} \rightarrow \text{credit score} \rightarrow \text{collateral}
\]
\textbf{Rule:} Do not repeat ``bank'' ten times. Use \textit{the institution, the lender, the branch, the financial partner}. Do not repeat ``money''. Use \textit{funds, capital, liquidity, balance, turnover}.
\end{propriete}

\begin{methode}[Planifier avant d'écrire -- 5 Steps to a Band 20/20]
\begin{enumerate}
    \item \textbf{Analyser le sujet (Public, But, Ton).} \textit{Ex: ``Lettre à la banque pour négocier des frais'' $\rightarrow$ Public = Directeur d'agence (formel), But = Négocier (persuasif), Ton = Respectueux mais ferme.}
    \item \textbf{Brainstormer arguments + vocabulaire clé.} Créez deux colonnes: \textit{Arguments POUR} (ancienneté, volume transactions, fidélité) et \textit{Vocabulaire bancaire} (tenue de compte, commission, découvert, échéancier).
    \item \textbf{Structurer en paragraphes (1 idée = 1 paragraphe).} Paragraphe 1: Contexte. Paragraphe 2: Argument 1 (comparaison). Paragraphe 3: Argument 2 (historique client). Paragraphe 4: Préférence + Call to Action.
    \item \textbf{Rédiger le brouillon.} Écrivez sans vous arrêter. Utilisez les connecteurs préparés (\textit{furthermore, however, consequently, provided that}).
    \item \textbf{Relire (Checklist 5 minutes).} [✓] Objet clair ? [✓] Salutation/Clôture correctes ? [✓] 0 contraction (\textit{I'm $\rightarrow$ I am}) ? [✓] Au moins 2 phrases complexes (subordination) ? [✓] Vocabulaire technique (5+ termes) ? [✓] Call to action précis ?
\end{enumerate}
\end{methode}

\begin{attention}[Formal Register Traps -- Automatic Mark Deductions]
In formal banking writing, the following are \textbf{penalised} at the \textsc{Bac} Technique:
\begin{itemize}
    \item \textbf{Contractions:} \textit{I'm, don't, can't, won't, it's} $\rightarrow$ \textbf{Always write:} \textit{I am, do not, cannot, will not, it is}.
    \item \textbf{Abbreviations:} \textit{etc., e.g., i.e., approx., info} $\rightarrow$ \textbf{Write out:} \textit{and so on, for example, that is, approximately, information}.
    \item \textbf{Colloquialisms:} \textit{kids, guys, cool, okay, a lot of, get} $\rightarrow$ \textbf{Use:} \textit{children, colleagues, satisfactory, acceptable, a significant number of, obtain/receive}.
    \item \textbf{Active voice for objectivity:} Prefer \textit{It is requested that the fees be revised} over \textit{I want you to revise the fees}.
    \item \textbf{``Although... but'' error:} \textcolor{red}{Although I am a loyal customer, but I am unhappy.} $\rightarrow$ \textcolor{popgreen}{\textbf{Although I am a loyal customer, I am unhappy.}} \textbf{OR} \textcolor{popgreen}{\textbf{I am a loyal customer, but I am unhappy.}} (Never both connectors).
\end{itemize}
\end{attention}

\courssub{Writing Practice: Two Examination Topics}

You will choose \textbf{ONE} topic. Both require 250--300 words, formal register, banking vocabulary, and complex sentences (coordination + subordination).

\begin{exemplebox}[Données comparatives réelles -- Cameroun 2024/2025]
Utilisez ces chiffres concrets pour étayer vos arguments (sources : sites web banques, MINCOMMERCE, GIMAC).
\begin{center}
\begin{tabularx}{\linewidth}{|l|X|X|}\hline
\textbf{Frais / Service} & \textbf{Compte Bancaire Classique (Moyenne 3 banques)} & \textbf{Mobile Money (Orange Money / MTN MoMo)} \\\hline
Ouverture de compte & 5~000 -- 10~000 FCFA (souvent offert si dépôt initial) & \textbf{0 FCFA} (gratuite, carte d'identité seulement) \\\hline
Frais de tenue de compte (annuel) & \textbf{12~000 -- 25~000 FCFA} & \textbf{0 FCFA} \\\hline
Retrait au guichet / agent & Gratuit (propre banque) / 1~000--2~000 FCFA (autre) & \textbf{1\% -- 2\% du montant} (min 50 FCFA, max 10~000 FCFA) \\\hline
Virement interne & 500 -- 1~500 FCFA & Gratuit (même opérateur) \\\hline
Virement interbancaire / inter-opérateur & 2~000 -- 5~000 FCFA + délai 24--48h & 100 -- 500 FCFA (instantané) \\\hline
Plafond journalier retrait & Selon convention (souvent 500~000 -- 1~000~000 FCFA) & \textbf{500~000 FCFA} (Orange) / \textbf{1~000~000 FCFA} (MTN) \\\hline
Accès au crédit & \textbf{Oui (prêt, découvert, leasing)} -- nécessite historique & \textbf{Limité (nano-crédit: 5~000--500~000 FCFA, court terme)} \\\hline
Sécurité / Assurance dépôts & \textbf{Couvert par le Fonds de Garantie (jusqu'à 50~M FCFA)} & Non couvert par fonds de garantie bancaire \\\hline
\end{tabularx}
\end{center}
\textbf{Calcul sur 5 ans (exemple):} Tenue compte 18~000 FCFA/an $\times$ 5 = 90~000 FCFA. Mobile Money: 0 FCFA frais fixes, mais si 2 retraits/mois de 50~000 FCFA $\times$ 1,5\% $\times$ 60 mois = 90~000 FCFA frais variables. \textbf{Le point d'équilibre dépend de votre volume de retraits.}
\end{exemplebox}

\begin{exoresolu}[Sujet 1: ``Avantages et inconvénients du Mobile Money vs Compte bancaire classique pour un jeune entrepreneur camerounais'']
\textbf{Plan détaillé + Introduction \& Conclusion modèles.}

\textbf{Title:} \textbf{Mobile Money vs Traditional Bank Account: A Strategic Choice for Young Cameroonian Entrepreneurs}

\textbf{Introduction (40 mots):}
In Cameroon's dynamic digital economy, young entrepreneurs face a critical financial decision: relying solely on Mobile Money services like Orange Money or MTN MoMo, or opening a traditional bank account. While both offer transactional capabilities, they serve distinct strategic needs. This essay compares their respective advantages and limitations to guide an informed choice.

\textbf{Body Paragraph 1: Accessibility \& Costs (70 mots):}
Mobile Money excels in accessibility. Opening an account requires only a national ID card and a phone number, with \textbf{zero opening fees} and \textbf{no maintenance charges}. For a startup with irregular cash flow, this eliminates the \textbf{12~000--25~000 FCFA annual fees} typical of banks (UBA, Afriland, BICEC). \textbf{However}, withdrawal fees of \textbf{1--2\%} accumulate rapidly; an entrepreneur withdrawing 500~000 FCFA monthly pays 7~500--15~000 FCFA monthly -- exceeding annual bank fees.

\textbf{Body Paragraph 2: Credit Access \& Business Growth (70 mots):}
\textbf{Conversely}, a traditional bank account builds a \textbf{credit history} -- the gateway to \textbf{business loans, overdraft facilities, and leasing}. Banks require \textbf{6--12 months of statements} to assess \textbf{repayment capacity}. Mobile Money offers only \textbf{nano-credits (up to 500~000 FCFA)} at high daily rates, insufficient for equipment purchase or stock financing. \textbf{Provided that} the entrepreneur maintains clean statements, the bank becomes a \textbf{growth partner}, not just a custodian.

\textbf{Body Paragraph 3: Security \& Professional Image (50 mots):}
Bank deposits are \textbf{insured up to 50 million FCFA} by the Deposit Guarantee Fund (Fonds de Garantie des Dépôts), whereas Mobile Money balances lack this protection. Moreover, a corporate bank account (\textit{compte professionnel}) enhances \textbf{credibility with suppliers and institutional clients} who often reject Mobile Money numbers for B2B payments.

\textbf{Conclusion + Recommendation (40 mots):}
\textbf{In conclusion}, Mobile Money is \textbf{ideal for daily micro-transactions and financial inclusion}, but a traditional account is \textbf{indispensable for scaling}. \textbf{I would strongly advise} a \textbf{hybrid strategy}: use Mobile Money for retail collections and the bank account for savings, supplier payments, and credit building.

\textbf{Total: $\approx$ 270 mots.}
\textbf{Grammar highlights:} \textit{While both offer... (subordination)}, \textbf{However / Conversely} (coordination/adverbials), \textbf{Provided that} (condition), \textbf{I would strongly advise} (preference), \textbf{ideal for / indispensable for} (prepositional collocations).
\tcblower
\textbf{Solution détaillée :} Ce modèle montre la structure attendue : introduction avec thèse, trois paragraphes de développement (accès/coûts, crédit/croissance, sécurité/image), conclusion avec recommandation claire. Le vocabulaire bancaire est précis (turnover, overdraft, credit history, nano-credits, deposit guarantee fund). Les connecteurs variés assurent la cohésion. Les structures complexes incluent la subordination (While, Provided that), la coordination (However, Conversely), l'expression de préférence (I would strongly advise).
\end{exoresolu}

\begin{exercice}[Sujet 2 -- À vous de jouer !]
\textbf{Rédigez la lettre formelle complète (250--300 mots).}
\textbf{Sujet:} \textit{``Vous êtes titulaire d'un compte courant à la Société Générale Cameroun depuis 3 ans. Votre chiffre d'affaires mensuel moyen est de 1~500~000 FCFA. Vous jugez les frais de tenue de compte (22~000 FCFA/trimestre) excessifs par rapport à la concurrence. Écrivez au Directeur d'Agence pour négocier une baisse à 12~000 FCFA/trimestre.''}

\textbf{Contraintes obligatoires:}
\begin{enumerate}
    \item Structure complète (En-tête, Objet, 3 paragraphes corps, Call to action, Formule longue, Pièces jointes).
    \item Minimum 5 termes de vocabulaire bancaire (surlignez-les).
    \item Minimum 2 phrases de subordination (\textit{although, provided that, on condition that, in order that}).
    \item Minimum 2 connecteurs de coordination (\textit{not only... but also, furthermore, consequently}).
    \item Zéro contraction, zéro abréviation, zéro langage familier.
\end{enumerate}
\end{exercice}

\courssub{Self-Assessment Grid \& Peer Review}

Autonomous learners evaluate their own work \emph{before} the teacher does. Use this radar chart criteria -- each axis scored 0--4 (total 20).

\begin{popfigure}
\begin{tabular}{|l|cc|}\hline
\textbf{Critère} & \textbf{Score (0--4)} \\ \hline
Content / Relevance & \dots\,/4 \\ \hline
Organisation \& Connectors & \dots\,/4 \\ \hline
Banking Vocabulary & \dots\,/4 \\ \hline
Complex Grammar & \dots\,/4 \\ \hline
Spelling \& Mechanics & \dots\,/4 \\ \hline
\end{tabular}

\smallskip
\textbf{Scoring Guide (0--4 each) :} 4 = Excellent ; 3 = Good ; 2 = Fair ; 1 = Weak ; 0 = Absent. \textbf{Total /20} = somme des 5 critères.
\legende{Grille d'auto-évaluation (5 critères $\times$ 4 points = 20).}
\end{popfigure}

\begin{center}
\begin{tabularx}{\linewidth}{|l|X|}\hline
\textbf{Critère} & \textbf{Indicateurs de réussite (Score 4/4)} \\\hline
\textbf{Content Relevance} & All prompt requirements addressed; arguments are specific, quantified (using data from \S5.3), and logically developed; clear stance/preference expressed. \\\hline
\textbf{Organisation $\&$ Connectors} & 4--5 distinct paragraphs (1 idea each); logical flow (Context $\rightarrow$ Arguments $\rightarrow$ CTA); varied connectors (\textit{furthermore, however, consequently, provided that, in addition}); no ``and/but/so'' only. \\\hline
\textbf{Banking Vocabulary} & 8+ precise terms used correctly (\textit{maintenance fees, overdraft, collateral, repayment schedule, creditworthiness, KYC, turnover, discrepancy, refund, statement}). No vague words (\textit{money, paper, thing}). \\\hline
\textbf{Complex Grammar} & 3+ complex sentences: at least 1 subordination (\textit{although, provided that, in order that}), 1 coordination (\textit{not only... but also}), 1 preference structure (\textit{would rather, would prefer, it is advisable that}). Passive used for objectivity. \\\hline
\textbf{Spelling $\&$ Mechanics} & Zero contractions, zero abbreviations, correct formal salutations/closings, accurate punctuation, consistent British/UK spelling (\textit{cheque, programme, organise}). \\\hline
\end{tabularx}
\end{center}

\courssub{Guided Rewriting: Improving Cohesion \& Variety}

Writing is rewriting. The gap between a ``passable'' (10/20) and an ``excellent'' (18+/20) text is almost always \textbf{cohesion} (flow) and \textbf{syntactic variety}. Let us transform a typical student draft.

\begin{popfigure}
\begin{tikzpicture}[
    textnode/.style={text width=0.45\linewidth, align=left, inner sep=4pt},
    comment/.style={text width=0.45\linewidth, align=left, inner sep=4pt, fill=poporange!10, draw=poporange, rounded corners=3pt, font=\small},
    arrow/.style={-Stealth, poporange, thick, shorten <=2pt, shorten >=2pt}
]
% Original text (left)
\node[textnode, anchor=north west] (orig) at (0,0) {
    \textbf{Original Draft (Student)}\par
    \textit{Dear Sir,}\par
    \textit{I write this letter because I want to complain about fees. My account number is 10002-034. I am a customer since 3 years. The fees are too high. Other banks are cheaper. I want you to reduce them. If you don't reduce, I will close my account. I hope you answer me.}\par
    \textit{Yours faithfully,}\par
    \textit{Paul}
};
% Annotated comments (right)
\node[comment, anchor=north east] (com1) at (\linewidth, 0) {
    \textbf{Commentaires enseignant:}\par
    \textcolor{red}{$\times$} \textit{``I write this letter because''} $\rightarrow$ trop oral. \textbf{Use:} \textit{I am writing to formally request...}\par
    \textcolor{red}{$\times$} \textit{``I am a customer since 3 years''} $\rightarrow$ erreur grammaire (since + point de départ). \textbf{Use:} \textit{I have been a customer \textbf{for} 3 years / \textbf{since} 2022.}\par
    \textcolor{red}{$\times$} \textit{``Other banks are cheaper''} $\rightarrow$ vague. \textbf{Use:} \textit{Competitor analysis shows UBA charges 4,500 FCFA vs your 6,250 FCFA.}\par
    \textcolor{red}{$\times$} \textit{``If you don't reduce...''} $\rightarrow$ menace, condition mal formulée. \textbf{Use:} \textit{Unless a revision is granted, I \textbf{will be compelled to} transfer my account.}\par
    \textcolor{red}{$\times$} \textit{``I hope you answer me''} $\rightarrow$ pas de Call to Action. \textbf{Use:} \textit{I \textbf{look forward to your written response within 10 working days.}}
};
% Improved version (bottom)
\node[textnode, anchor=north west, fill=popgreen!10, draw=popgreen, rounded corners=3pt] (improved) at (0, -7.5) {
    \textbf{Rewritten Version (Target 18+/20)}\par
    \textit{Dear Sir,}\par
    \textit{I am writing to formally request a review of the quarterly maintenance fees applied to my current account (No. 10002-034), which I have held for three years.}\par
    \textit{\textbf{Not only} does my average monthly turnover reach 1,800,000 FCFA, \textbf{but} I have also never used an overdraft facility. \textbf{However}, a benchmarking of three major banks reveals your fee of 6,250 FCFA/quarter exceeds the market average of 4,800 FCFA by 30\%.}\par
    \textit{\textbf{Provided that} you agree to align your tariff to 4,500 FCFA/quarter, I \textbf{would rather} maintain my relationship with your institution than switch providers. \textbf{Otherwise}, I \textbf{will be compelled to} transfer my domiciliation.}\par
    \textit{\textbf{I kindly request a meeting next week to finalise this amendment} and look forward to your written confirmation within ten working days.}\par
    \textit{Yours faithfully,}\par
    \textit{Paul M.}\par
    \textit{Enc.: Fee comparison table (3 banks), Last 6 months statements}
};
% Arrows
\draw[arrow] (orig.east) -- (com1.west);
\draw[arrow] (orig.south) -- (improved.north);
\draw[arrow] (com1.south) -- (improved.north east);
\end{tikzpicture}
\legende{Processus de réécriture guidée : texte original $\rightarrow$ commentaires ciblés (grammaire, vocabulaire, registre, cohésion) $\rightarrow$ version améliorée. Notez l'intégration de \textit{not only... but also, however, provided that, would rather, otherwise, compelled to}.}
\end{popfigure}

\begin{exoresolu}[Exercice de réécriture -- Appliquez la méthode]
\textbf{Consigne:} Réécrivez le paragraphe suivant en améliorant : (1) la cohésion lexicale (chaînes bancaires), (2) la variété syntaxique (subordination + coordination), (3) le registre formel.

\textbf{Texte élève:}
\textit{``Banks are important for business. They give loans. You need a good account. If you have a good account, the bank gives you money. But fees are high. I think banks should lower fees. Because small businesses don't have much money. Also mobile money is cheaper. But mobile money doesn't give big loans. So you need both. I prefer bank for big things and mobile money for small things.''}

\textbf{Version révisée (modèle):}
\textit{``Financial institutions play a pivotal role in enterprise development by providing \textbf{credit facilities} tailored to \textbf{cash flow} cycles. \textbf{Provided that} a business maintains a \textbf{healthy account history} -- characterised by consistent \textbf{turnover} and zero \textbf{default} -- it unlocks access to \textbf{working capital loans} and \textbf{overdraft facilities}. \textbf{Nevertheless}, \textbf{exorbitant maintenance fees} (often 20,000+ FCFA/quarter) \textbf{disproportionately burden} micro-enterprises. \textbf{While} Mobile Money offers \textbf{cost-effective transactionality} for \textbf{retail collections}, it \textbf{lacks the capacity} for \textbf{long-term asset financing}. \textbf{Consequently}, a \textbf{dual-channel strategy} -- leveraging Mobile Money for \textbf{liquidity management} and a traditional account for \textbf{creditworthiness building} -- \textbf{is advisable}.''}

\textbf{Analyse des améliorations:}
\begin{itemize}
    \item \textbf{Cohésion lexicale:} \textit{financial institutions $\rightarrow$ credit facilities $\rightarrow$ cash flow $\rightarrow$ account history $\rightarrow$ turnover $\rightarrow$ default $\rightarrow$ working capital $\rightarrow$ overdraft $\rightarrow$ maintenance fees $\rightarrow$ micro-enterprises $\rightarrow$ transactionality $\rightarrow$ asset financing $\rightarrow$ dual-channel $\rightarrow$ liquidity management $\rightarrow$ creditworthiness}. (17 maillons !)
    \item \textbf{Syntaxe:} \textit{Provided that... (subordination conditionnelle), While... (concession), Consequently... (conséquence), Passive: is advisable.}
    \item \textbf{Registre:} Zéro contraction, vocabulaire précis (\textit{pivotal, disproportionately, exorbitant, retail collections, asset financing}).
\end{itemize}
\tcblower
\textbf{Solution détaillée :} La révision transforme un texte élémentaire (phrases courtes, vocabulaire vague, répétitions) en un texte académique cohérent. La chaîne lexicale bancaire crée la cohésion. Les subordonnées (Provided that, While) et les connecteurs (Nevertheless, Consequently) structurent l'argumentation. Le registre formel élimine les contractions et utilise le passif pour l'objectivité.
\end{exoresolu}

\courssub{Publication \& Real-World Application}

The ultimate test of a banking text is its effect in the real world. The Cameroonian legal framework empowers you to use these writing skills concretely.

\begin{savaistu}[Vos droits réels -- Loi camerounaise]
La \textbf{Loi n°2011/012 du 6 mai 2011 (Code de la Consommation)}, en son Article 17, interdit les \textbf{clauses abusives} dans les contrats bancaires (frais disproportionnés, modification unilatérale sans préavis, renonciation aux recours). L'Article 42 vous donne le droit de saisir la \textbf{Commission Nationale de la Consommation (CNC)} auprès du \textbf{MINCOMMERCE} pour tout litige non résolu par la banque en 30 jours.
\textbf{Application pratique:} Votre lettre de réclamation (Sujet 2, \S5.3) n'est pas un exercice scolaire -- c'est une \textbf{mise en demeure} juridique. Conservez une copie signée + accusé de réception. Si pas de réponse sous 30 jours $\rightarrow$ Saisine MINCOMMERCE (gratuit). C'est ce pouvoir réel qui donne du sens à votre apprentissage de l'anglais bancaire formel.
\end{savaistu}

\begin{exercice}[Tâche finale intégrée -- Portfolio Exam Ready]
\textbf{Choisissez UNE des trois options et produisez une version finale « prête à envoyer / publier ».}
\begin{enumerate}
    \item \textbf{Option A (Lettre réelle):} Finalisez votre lettre de négociation de frais (Sujet 2, \S5.3) après réécriture (\S5.5). Imprimez-la, signez-la, mettez-la sous enveloppe avec pièces jointes simulées. \textbf{Remettez au professeur pour évaluation sommative (sur 20 avec grille radar).}
    \item \textbf{Option B (Article de blog):} Rédigez l'article ``\textit{How to Choose Your Bank in Cameroon: A 2025 Guide for Tech Students}'' (300--350 mots). Cible: vos camarades de Terminale. Ton: informel \textbf{mais} vocabulaire technique précis. Publiez sur le blog de classe / groupe WhatsApp de la section.
    \item \textbf{Option C (Rapport d'évaluation):} En binôme, comparez \textbf{deux offres réelles} de « Compte Jeune Entrepreneur » (ex: Afriland \textit{Jeune Pro} vs UBA \textit{NextGen} vs BICEC \textit{Start}). Produisez un rapport structuré (Intro, Tableau comparatif, Analyse coûts/bénéfices sur 3 ans, Recommandation conditionnelle \textit{``If X, then Y; provided that Z''}). Présentez oralement 3 min (évaluation \textsc{Speaking}).
\end{enumerate}
\textbf{Critères communs:} Grille radar (\S5.4) $\ge$ 16/20 + zéro faute « Attention » (\S5.2) + Call to action explicite.
\end{exercice}

\begin{corrige}[Exercice Sujet 2 -- Lettre de négociation (Modèle de référence)]
\textbf{Yaoundé, 15 October 2025}\par
\textbf{The Branch Manager}\par
\textbf{Société Générale Cameroun, Akwa Branch}\par
\textbf{P.O. Box 123, Douala}\par
\vspace{0.3cm}
\textbf{OBJET: REQUEST FOR REDUCTION OF QUARTERLY ACCOUNT MAINTENANCE FEES}\par
\textbf{Account No: 10002-0345678-98}\par
\vspace{0.3cm}
\textbf{Dear Sir,}\par
I am writing as a loyal corporate client of three years (Account No. 10002-0345678-98) to formally request a revision of the quarterly maintenance fees currently charged at 22,000 FCFA.\par
\textbf{Not only} does my account record an average monthly turnover of 1,500,000 FCFA, \textbf{but} I have also maintained a clean history with zero overdraft incidents and full KYC compliance. \textbf{However}, a comparative benchmarking of three major banks in Douala reveals that your fee exceeds the market average of 14,000 FCFA by 57\%. For instance, UBA charges 12,000 FCFA and BICEC 13,500 FCFA for equivalent service packages.\par
\textbf{Provided that} you agree to align the quarterly fee to 12,000 FCFA -- reflecting my volume and loyalty -- I \textbf{would strongly prefer} to consolidate my banking relationship with Société Générale rather than migrate to a competitor. \textbf{Otherwise}, I \textbf{will be compelled to} initiate account closure procedures.\par
\textbf{I kindly request a meeting within the next ten working days to formalise this amendment} and look forward to your favourable written response.\par
\vspace{0.3cm}
\textbf{Yours faithfully,}\par
\vspace{1cm}
\textbf{(Signature)}\par
\textbf{Jean-Pierre M.}\par
\textbf{General Manager, JPM Agro-Services SARL}\par
\vspace{0.3cm}
\textbf{Enc.: 1) Bank statements (Jan--Jun 2025)  2) Fee comparison table (3 banks)  3) Business registration (Extrait K-bis)}\par
\vspace{0.2cm}
\textit{Word count: 268. Grammar: Subordination (Provided that), Coordination (Not only... but also, However), Preference (would strongly prefer), Passive (be compelled to), Zero contractions/abbreviations. Vocabulary: 12+ banking terms.}
\end{corrige}

\par
\textbf{Teacher's Note:} This Writing Workshop completes the \textit{Economic Life and Occupations} unit. The integration activity (Section 6) will combine \textbf{Listening} (price comparison audio), \textbf{Speaking} (role-play bank appointment), \textbf{Reading} (bank contract analysis), and this \textbf{Writing} portfolio in a single simulated ``Financial Literacy Day'' assessment. Prepare your portfolio (letter + essay + self-assessment grid) for submission.

\coursec{Integration, Evaluation \& Remediation}

In the previous sections, you practised assessing consumer needs, comparing prices and quality, mastering banking vocabulary, building complex sentences (coordination and subordination), expressing preferences, and writing compositions about banking services. Now it is time to bring \emph{everything} together. In this final section, you will use all the skills of the lesson in one realistic project. Then you will be evaluated, evaluate yourself, and receive targeted remediation where needed. This is exactly the spirit of the \cle{integration activities} of the official programme.

\begin{definition}[What is integration?]
\cle{Integration} is the joint mobilisation of knowledge (\emph{savoirs}), know-how (\emph{savoir-faire}) and attitudes (\emph{savoir-être}) in a new, complex situation close to real professional life. It is not a simple revision exercise: you must \emph{transfer} what you learned in Sections 1 to 5 to solve an authentic problem you have never seen before.
\end{definition}

\courssub{Integrative Project: A Simplified Business Plan for a Village Cooperative}

\textbf{The situation.} Your village (or your school neighbourhood) wants to create a small \cle{cooperative} (\emph{une coopérative}) to buy food products in bulk and resell them at fair prices to members. Your class, divided into teams of four, must produce a \textbf{simplified business plan in English} and defend it orally in front of a ``banker'' (your teacher).

The project naturally recycles every skill of the lesson:

\begin{itemize}
\item \textbf{Section 1 — Assessing needs:} design a short questionnaire and survey 10 potential members about their consumer needs and the community resources available.
\item \textbf{Section 2 — Smart shopping:} compare prices and quality of three suppliers to find the \cle{best buy} for bulk purchases (\emph{achats groupés}).
\item \textbf{Section 3 — Banking services:} choose a bank or microfinance institution (\emph{institution de microfinance}), compare \cle{interest rates}, \cle{account fees} and \cle{loan conditions}, and justify your choice.
\item \textbf{Sections 4–5 — Argumentation and writing:} write a one-page report using complex sentences (coordination, subordination) and expressions of preference (\emph{We would rather open a savings account because...}).
\end{itemize}

\begin{methode}[The project approach — démarche projet]
Follow these five steps, in order:
\begin{enumerate}
\item \textbf{Define the real problem.} What do people in the community need? What is missing? Write one clear problem sentence: \emph{``Families in our quarter pay too much for rice and soap.''}
\item \textbf{Plan tasks and roles.} In each team, appoint a \cle{surveyor}, a \cle{price analyst}, a \cle{banking researcher} and a \cle{writer/presenter}. Every member speaks during the final defence.
\item \textbf{Collect information.} Use fieldwork (interviews, market visits) and the internet (bank websites, price lists). Record data in tables.
\item \textbf{Produce the deliverables.} A written report (250--300 words) with a comparison table, and a 5-minute oral presentation.
\item \textbf{Present and defend.} Answer the ``banker's'' questions using comparatives, adverbs of time (\emph{for, since, ago}) and preference structures.
\end{enumerate}
\end{methode}

\begin{popfigure}
\begin{tikzpicture}[x=1cm,y=0.85cm]
% Axes
\draw[->,thick,popdark] (0,0) -- (16,0) node[below left]{\small Days};
\draw[->,thick,popdark] (0,0) -- (0,6.5) node[left]{\small Tasks};
% Vertical grid lines for days
\foreach \d/\lab in {0/Mon,3/Wed,5/Fri,8/Wed,10/Fri,13/Wed,15/Fri}{
  \draw[popline,thin] (\d,0) -- (\d,6.2);
  \node[below,font=\small,popdark] at (\d,0) {\lab};
}
% Week labels background
\fill[popblue!10] (0,0) rectangle (5,6.2);
\fill[poporange!10] (5,0) rectangle (10,6.2);
\fill[popgreen!10] (10,0) rectangle (15,6.2);
% Week labels
\node[popblue,font=\small\bfseries,rotate=90] at (-0.5,2.5) {Week 1: Survey \& Comparison};
\node[poporange,font=\small\bfseries,rotate=90] at (-0.5,7.5) {}; % dummy for spacing
\node[poporange,font=\small\bfseries] at (7.5,6.0) {Week 2: Banking \& Writing};
\node[popgreen,font=\small\bfseries] at (12.5,6.0) {Week 3: Presentation};
% Bars (Top-down: y=5.5 to 1.0)
% 1. Needs survey (Week 1)
\fill[popblue] (0,5.5) rectangle (5,4.9); \node[left,font=\small,align=right,xshift=-0.2cm] at (0,5.2) {Needs survey};
% 2. Price comparison (Week 1)
\fill[popblue!60] (1,4.8) rectangle (5,4.2); \node[left,font=\small,align=right,xshift=-0.2cm] at (1,4.5) {Price comparison};
% 3. Bank research (Week 2)
\fill[poporange] (5,4.1) rectangle (10,3.5); \node[left,font=\small,align=right,xshift=-0.2cm] at (5,3.8) {Bank research};
% 4. Report writing (Week 2)
\fill[poporange!60] (6,3.4) rectangle (11,2.8); \node[left,font=\small,align=right,xshift=-0.2cm] at (6,3.1) {Report writing};
% 5. Rehearsal (Week 3)
\fill[popgreen] (11,2.7) rectangle (14,2.1); \node[left,font=\small,align=right,xshift=-0.2cm] at (11,2.4) {Rehearsal};
% 6. Final defence (Week 3)
\fill[popgreen!60] (14,2.0) rectangle (15.5,1.4); \node[left,font=\small,align=right,xshift=-0.2cm] at (14,1.7) {Final defence};
% Milestones (diamonds at end of bars)
\node[popdark,font=\small] at (5,5.2) {$\blacklozenge$ Survey report};
\node[popdark,font=\small] at (11,3.1) {$\blacklozenge$ Draft ready};
\node[popdark,font=\small] at (15.5,1.7) {$\blacklozenge$ Plan defended};
\end{tikzpicture}
\legende{Gantt chart of the integrative project over three weeks (15 working days). Bars show task duration and overlap. Diamonds mark the three deliverables: survey report, written draft, final oral defence.}
\end{popfigure}

\begin{exoresolu}[Sample extract from a team's report]
\textbf{Task:} write the ``banking choice'' paragraph of the business plan (about 80 words), using at least one comparative, one adverb of time and one expression of preference.
\tcblower
\textbf{Model answer:} \emph{``We have compared three banks \textbf{for} two weeks. Bank A offers \textbf{lower} account fees than Bank B, \textbf{but} its interest rate on loans is \textbf{twice as high as} Bank C's. \textbf{Since} January, Bank C has also provided free mobile banking, \textbf{which} is useful in our village. \textbf{We would rather} open our account at Bank C \textbf{because} it is cheaper and closer to our members.''}\\
\textbf{Why it works:} comparative of inferiority/superiority (\emph{lower... than}), double comparison (\emph{twice as high as}), adverbs of time (\emph{for, since}), coordination (\emph{but, because}), subordination (\emph{which}), preference (\emph{would rather}). All the grammar targets of the lesson appear naturally.
\end{exoresolu}

\courssub{Formative Evaluation}

Evaluation here is \cle{formative}: its goal is to show you what you master and what you must still improve \emph{before} the exam.

\textbf{Activity 1 — Digital quiz (Kahoot/Quizizz, 15 min).} Twenty questions on the vocabulary of banking and consumption (\emph{loan, overdraft, discount, value for money...}) and on grammar (comparisons of equality and degree, adverbs \emph{for/since/ago}, coordination vs subordination). Instant feedback allows immediate correction of misconceptions.

\textbf{Activity 2 — Assessed role-play: negotiating a loan.} In pairs, one student is the \cle{loan applicant}, the other the \cle{bank manager}. The applicant must explain the cooperative's needs, compare two offers and express a preference; the manager asks questions about guarantees and repayment. Assessment criteria: fluency (2 pts), correct grammar in context (2 pts), appropriate vocabulary (1 pt).

\textbf{Activity 3 — Marked written production.} A 250-word composition (letter or report) marked on the official-style grid below.

\begin{attention}[How writing is marked at the Bac]
At the Baccalauréat technique, written expression is marked out of 10 points: \textbf{4 points} for content and coherence (\emph{contenu et cohérence}), \textbf{3 points} for organisation and connectors (\emph{organisation et connecteurs}), \textbf{3 points} for linguistic accuracy (\emph{correction linguistique}: grammar, spelling, punctuation). Do not neglect any dimension — a brilliant idea written with poor grammar cannot score more than about 5/10.
\end{attention}

\begin{popfigure}
\begin{tikzpicture}
\begin{axis}[
  xbar stacked, width=\linewidth, height=6.5cm,
  xmin=0, xmax=100,
  symbolic y coords={Vocabulary,Grammar,Writing,Reading,Listening,Speaking},
  ytick=data, y dir=reverse,
  xlabel={Percentage of the class (\%)},
  legend style={at={(0.5,-0.28)},anchor=north,legend columns=3,font=\small},
  bar width=9pt, font=\small,
]
\addplot[fill=popgreen] coordinates {(66,Vocabulary) (40,Grammar) (48,Writing) (70,Reading) (55,Listening) (62,Speaking)};
\addplot[fill=poporange] coordinates {(22,Vocabulary) (33,Grammar) (30,Writing) (20,Reading) (28,Listening) (25,Speaking)};
\addplot[fill=poppink] coordinates {(12,Vocabulary) (27,Grammar) (22,Writing) (10,Reading) (17,Listening) (13,Speaking)};
\legend{Mastered, In progress, Not yet acquired}
\end{axis}
\end{tikzpicture}
\legende{Stacked-bar evaluation grid: class profile per skill after the formative evaluation (fictitious example). Grammar and Writing show the largest pink zones — they become the priority for remediation.}
\end{popfigure}

Reading the chart: if, like in this example, \textbf{Grammar} and \textbf{Writing} show the largest pink zones, the remediation workshops below will target them first.

\courssub{Targeted Remediation Workshops}

\begin{propriete}[Principles of effective remediation]
Effective remediation is: \textbf{targeted} (one single notion at a time), \textbf{short} (15--20 minutes), \textbf{interactive} (a game or a contextualised exercise, never a lecture repeated), and \textbf{assessed} (a mini exit-test to check progress).
\end{propriete}

Students are assigned to one workshop according to the formative results. Each workshop lasts 20 minutes and ends with a 5-question exit quiz.

\textbf{Workshop A — Grammar difficulties (comparisons and adverbs).} Card game: match sentence halves (\emph{``This loan is twice...'' / ``...as expensive as the other.''}). Then a ``find the mistake'' race: \emph{*I live here since five years ago} $\rightarrow$ \emph{I have lived here \textbf{for} five years / \textbf{since} 2020.}

\textbf{Workshop B — Banking vocabulary.} Word-wall and bingo with definitions: \cle{deposit} (\emph{dépôt}), \cle{withdrawal} (\emph{retrait}), \cle{overdraft} (\emph{découvert}), \cle{interest rate} (\emph{taux d'intérêt}), \cle{statement} (\emph{relevé de compte}). Exit task: label a mini bank statement with five of these words.

\textbf{Workshop C — Writing structure.} Rebuild a scrambled letter of complaint (paragraphs in envelopes), then highlight the connectors (\emph{first, moreover, however, finally}) and the subordinate clauses. Exit task: write the opening and closing paragraphs of a complaint letter.

\begin{exercice}[Exit quiz — Workshop A]
Choose the correct option: (1) I have been a customer here \emph{for / since} 2019. (2) This account is \emph{as / so} convenient as the mobile one. (3) The more you save, \emph{the more / the most} interest you earn. (4) He left two hours \emph{ago / since}. (5) Rice is \emph{cheaper / cheapest} than imported pasta.
\end{exercice}

\begin{corrige}[Exit quiz — Workshop A]
(1) \emph{since} (a starting date); (2) \emph{as} (equality: \emph{as... as}); (3) \emph{the more} (double comparative); (4) \emph{ago} (a finished past moment); (5) \emph{cheaper} (comparison of two items).
\end{corrige}

\courssub{Student Self-Assessment}

Each student completes the grid below honestly, then writes a personal action plan.

\begin{center}
\begin{tabularx}{\linewidth}{|X|c|c|}
\hline
\rowcolor{popblueL} \textbf{Lesson objective} & \textbf{I can do it} & \textbf{Not yet} \\
\hline
I can ask questions to assess consumer needs and community services. & $\square$ & $\square$ \\
\hline
I can compare prices and quality (comparatives, equality, degree). & $\square$ & $\square$ \\
\hline
I can use adverbs of time correctly (\emph{for, since, ago, afterwards}). & $\square$ & $\square$ \\
\hline
I can understand a text about banking services. & $\square$ & $\square$ \\
\hline
I can build complex sentences (coordination and subordination). & $\square$ & $\square$ \\
\hline
I can express my preferences (\emph{would rather, prefer... to}). & $\square$ & $\square$ \\
\hline
I can write a 250-word composition about banking services. & $\square$ & $\square$ \\
\hline
\end{tabularx}
\end{center}

\textbf{Personal action plan — example.} \emph{``I ticked `not yet' for subordinate clauses and banking vocabulary. My plan: (1) I will revise relative clauses (\emph{who, which, that}) for 15 minutes every evening this week; (2) I will learn 10 banking words per week with flashcards; (3) I will rewrite my complaint letter and ask the teacher for feedback before Friday.''}

\courssub{Cross-Curricular Links}

\begin{itemize}
\item \textbf{Economics:} every purchase involves an \cle{opportunity cost} (\emph{coût d'opportunité}) — choosing supplier A means giving up the advantages of supplier B. The cooperative's budget is a direct application of your economics course.
\item \textbf{Mathematics:} comparing offers requires percentages (a 10\,\% \cle{discount}), and banking requires interest calculations. Simple interest: $I = C \times t \times n$; for example, a deposit of \num{100000} FCFA at 5\,\% per year earns $100000 \times 0{,}05 \times 2 = \num{10000}$ FCFA in two years. Compound interest (\emph{intérêts composés}) grows faster because each year's interest is added to the capital.
\item \textbf{Citizenship education:} responsible consumption (avoiding waste, checking quality) and \cle{financial inclusion} — giving every citizen, even in rural areas, access to a bank account and credit — are pillars of community development in Cameroon.
\end{itemize}

\begin{popfigure}
\begin{tikzpicture}[
  mindmap, grow cyclic,
  every node/.style={concept,font=\small},
  concept color=popblue,
  level 1/.append style={level distance=3.2cm,sibling angle=72},
  level 2/.append style={level distance=1.9cm,sibling angle=50},
]
\node[concept color=popdark,text=white, align=center]{Integration\\ \emph{Intégration}}
  child[concept color=popblue] { node[align=center]{Needs\\ \emph{Besoins}} 
    child { node[concept color=popblueL] {survey} }
    child { node[concept color=popblueL] {resources} } }
  child[concept color=poporange] { node[align=center]{Comparison\\ \emph{Comparaison}}
    child { node[concept color=poporangeL] {best buy} }
    child { node[concept color=poporangeL] {as...as} } }
  child[concept color=popgreen] { node[align=center]{Banking\\ \emph{Banque}}
    child { node[concept color=popgreenL] {loan} }
    child { node[concept color=popgreenL] {interest} } }
  child[concept color=poppurple] { node[align=center]{Writing\\ \emph{Écriture}}
    child { node[concept color=poppurpleL] {connectors} }
    child { node[concept color=poppurpleL] {subordination} } }
  child[concept color=poppink] { node[align=center]{Evaluation\\ \emph{Évaluation}}
    child { node[concept color=poppinkL] {self-grid} }
    child { node[concept color=poppinkL] {remediation} } };
\end{tikzpicture}
\legende{Mind map of the lesson's key concepts: needs lead to comparison, comparison to banking choices, banking to written argumentation, and everything converges towards integration and evaluation.}
\end{popfigure}

\courssub{Preparing for the Probatoire and Baccalauréat}

\begin{exemplebox}[A real-type Bac subject (Technical series, 2023)]
\emph{``Write a letter to your bank manager complaining about unauthorized charges on your account (about 250 words).''} A model answer is available on the MINESEC website. Notice the expected structure: address and date, formal greeting (\emph{Dear Sir/Madam}), statement of the problem with dates and amounts, request for a refund, polite closing (\emph{Yours faithfully}).
\end{exemplebox}

\begin{savaistu}[The ``Case Study'' at the Probatoire technique]
The Probatoire technique often includes a banking \textbf{case study}: you receive an account statement (\emph{relevé de compte}), you must detect errors (a duplicated charge, a wrong balance), and propose solutions in correct English. This is exactly the skill trained in the integrative project when your team checks the cooperative's figures.
\end{savaistu}

\begin{methode}[Analysing a typical exam paper]
\begin{enumerate}
\item \textbf{Reading comprehension:} identify the text type (letter, statement, advert), then answer with words from the text and justify with line references.
\item \textbf{Grammar in context:} cloze exercises often test \emph{for/since/ago} and comparatives — re-read the time marker before choosing (\emph{since + starting point}, \emph{for + duration}, \emph{ago + finished past}).
\item \textbf{Written expression:} spend 5 minutes planning (3 paragraphs), 20 minutes writing, 5 minutes proofreading connectors and verb tenses.
\end{enumerate}
\end{methode}

\begin{aretenir}[Key points of Section 6]
Integration means using \emph{all} the lesson's skills together in a realistic project. Evaluation is a tool for progress, not a punishment: read your results, choose the right remediation workshop, and build a personal action plan. At the Bac, balance content, organisation and accuracy — and practise with real-type subjects (complaint letters, banking case studies) until the format becomes automatic.
\end{aretenir}

\newpage

\coursec{Activité d'intégration}

\coursec{LEÇON N02 : Economic Life and Occupations}

\courssub{Objectifs de l'activité}

This integrative project (\emph{projet intégrateur}) brings together all the skills of the unit \textbf{Economic Life and Occupations}. By the end of the activity, you will be able to:

\begin{itemize}
\item \textbf{Assess consumer needs} in your community using a questionnaire in English (\emph{évaluer les besoins des consommateurs}).
\item \textbf{Compare prices and quality} of suppliers to determine the best buys, using double comparatives and comparisons of equality and degree (\emph{comparer prix et qualité}).
\item \textbf{Analyse real banking offers} (interest rates, guarantees, fees) and write a formal loan application letter using coordination, subordination and expressions of preference (\emph{rédiger une demande de prêt}).
\item \textbf{Give and receive information orally} during a five-minute pitch in front of a ``bank jury'' (\emph{soutenance orale devant un jury}).
\item \textbf{Use adverbs of time} (for, since, ago, after, afterwards) correctly in speech and writing.
\end{itemize}

\begin{aretenir}[Competences mobilised]
Speaking: giving and receiving information. Vocabulary: consumer needs, best buys, banking services. Grammar: comparisons (double, equality, degree), adverbs of time, coordination, subordination, expressing preferences. Reading: banking documents. Writing: a formal letter and a business plan summary.
\end{aretenir}

\courssub{Situation}

\textbf{Project: ``Coopérative Agropastorale Jeunesse'' -- Business Plan \& Bank Pitch}

You are a group of five young school-leavers in \textbf{Bafoussam, West Region of Cameroon}. You want to create a youth cooperative (\emph{coopérative}) that will raise broiler chickens (\emph{poulets de chair}) and process cassava into gari (\emph{transformation du manioc}). To obtain a loan (\emph{prêt}) from a bank or a microfinance institution, you must prepare a complete file and defend it orally.

\textbf{Data collected during your community survey} (50 households interviewed):

\begin{center}
\begin{tabularx}{\linewidth}{|l|X|c|}\hline
\rowcolor{popblueL} \textbf{Rank} & \textbf{Community need} & \textbf{Households concerned} \\ \hline
1 & Affordable source of animal protein (chicken meat) & 41 out of 50 (82 \%) \\ \hline
2 & Good-quality gari at stable prices & 36 out of 50 (72 \%) \\ \hline
3 & Local jobs for young people & 29 out of 50 (58 \%) \\ \hline
\end{tabularx}
\end{center}

\textbf{Three suppliers of inputs} (\emph{fournisseurs d'intrants}) for 1 000 chicks and feed for one production cycle of 8 weeks:

\begin{center}
\begin{tabularx}{\linewidth}{|l|X|X|X|}\hline
\rowcolor{popblueL} \textbf{Supplier} & \textbf{Price per chick (FCFA)} & \textbf{Feed per 50 kg bag (FCFA)} & \textbf{Quality / reliability} \\ \hline
SODEPA Yaoundé & 600 & 19 500 & Vaccinated chicks, delivery in 2 days, very reliable \\ \hline
Ferme Nkoulou (Bafoussam) & 550 & 18 000 & Vaccinated chicks, local, sometimes late \\ \hline
Marché Mfoundi (informal) & 400 & 16 500 & Not always vaccinated, no guarantee \\ \hline
\end{tabularx}
\end{center}

\textbf{Two financing offers} for a loan of \num{1500000} FCFA over 24 months:

\begin{center}
\begin{tabularx}{\linewidth}{|l|X|X|}\hline
\rowcolor{popblueL} \textbf{Feature} & \textbf{Crédit Agricole (Offer A)} & \textbf{Microfinance ``Confiance'' (Offer B)} \\ \hline
Annual interest rate (\emph{taux d'intérêt}) & 8 \% & 15 \% \\ \hline
Guarantee (\emph{garantie}) & Land title or guarantor & Group solidarity only \\ \hline
File fees (\emph{frais de dossier}) & 25 000 FCFA & 10 000 FCFA \\ \hline
Decision time & 6 weeks & 1 week \\ \hline
Grace period (\emph{différé de remboursement}) & 3 months & None \\ \hline
\end{tabularx}
\end{center}

\textbf{Expected results per cycle} (8 weeks): 950 chickens sold at 3 500 FCFA each; cost of production about \num{2400000} FCFA per cycle of 1 000 chicks including feed, vaccines and labour.

\begin{savaistu}[Contexte bancaire camerounais]
Au Cameroun, les jeunes entrepreneurs peuvent s'adresser aux banques commerciales (ex. Crédit Agricole, Afriland First Bank, UBA) ou aux institutions de microfinance (IMF) agréées par la COBAC. Les banques exigent souvent des garanties réelles (titre foncier, caution) mais offrent des taux plus bas (7--10 \%). Les IMF acceptent la garantie solidaire de groupe, mais leurs taux sont plus élevés (12--18 \%). Le choix dépend de la capacité de garantie et de l'urgence du financement.
\end{savaistu}

\courssub{Consignes}

Work in groups of four or five. Answer the following tasks in English. Each task uses skills from the unit.

\begin{enumerate}
\item \textbf{Community needs.} Using the survey table, write three sentences identifying the three priority needs. Use comparisons of degree (e.g. ``Chicken meat is the most urgent need because...''). Then design two extra questionnaire questions you would ask households.
\item \textbf{Best buy.} Compare the three suppliers in a short paragraph. Use at least: one double comparative (``the cheaper the chicks, the higher the risk''), one comparison of equality (``as reliable as''), and one superlative. Justify your final choice of supplier.
\item \textbf{Unit price calculation.} Compute the total cost of 1 000 chicks plus 40 bags of feed for each supplier. Which supplier is the cheapest? Which one offers the best value for money (\emph{le meilleur rapport qualité-prix})? Explain the difference.
\item \textbf{Banking analysis.} Read the two financing offers. Extract the interest rate, the guarantee and the fees for each. Compute the total interest paid over two years for Offer A (simple interest: $I = C \times t \times n$). Say which offer you prefer and why, using expressions of preference (``We would rather...'', ``We prefer... to...'').
\item \textbf{Formal letter.} Write the loan application letter (150--180 words) to the institution you chose. Use at least two coordinated sentences (and, but, so) and two subordinated sentences (because, although, if, so that).
\item \textbf{Pitch preparation.} Prepare a five-minute oral pitch with four parts: needs, market, profitability, repayment plan. Use adverbs of time correctly: ``We have worked on this project \emph{for} six months'', ``We started our survey two months \emph{ago}'', ``\emph{After} the first cycle, we will repay...''.
\item \textbf{Jury simulation.} Two classmates and the teacher play the bankers. They ask questions with conditionals and concession: ``What will you do \emph{if} the price of feed increases?'' ``\emph{Although} your project is interesting, why is it not risky?'' Answer politely and justify.
\item \textbf{Self-evaluation.} Fill in the skills grid and identify your two weakest points for remediation.
\end{enumerate}

\courssub{Ressources à mobiliser}

\begin{aretenir}[Grammar toolbox]
\begin{itemize}
\item \textbf{Double comparative:} the + comparative..., the + comparative... (``The more we produce, the more we earn.'')
\item \textbf{Equality:} as + adjective + as (``Ferme Nkoulou is almost as reliable as SODEPA.'')
\item \textbf{Degree:} too + adjective, adjective + enough (``The informal market is too risky.'')
\item \textbf{Adverbs of time:} for + duration, since + starting point, ago (past), after / afterwards.
\item \textbf{Coordination:} and, but, or, so, yet. \textbf{Subordination:} because, although, if, when, so that, whereas.
\item \textbf{Preferences:} prefer X to Y, would rather + base verb, would prefer + to-infinitive.
\end{itemize}
\end{aretenir}

\begin{aretenir}[Key vocabulary]
\cle{consumer needs} (\emph{besoins des consommateurs}), \cle{survey} (\emph{enquête}), \cle{supplier} (\emph{fournisseur}), \cle{unit price} (\emph{prix unitaire}), \cle{value for money} (\emph{rapport qualité-prix}), \cle{loan} (\emph{prêt}), \cle{interest rate} (\emph{taux d'intérêt}), \cle{guarantee} (\emph{garantie}), \cle{fees} (\emph{frais}), \cle{repayment} (\emph{remboursement}), \cle{profit} (\emph{bénéfice}), \cle{savings account} (\emph{compte épargne}), \cle{grace period} (\emph{différé de remboursement}), \cle{broiler chickens} (\emph{poulets de chair}).
\end{aretenir}

\begin{methode}[Simple interest formula]
For a loan of capital $C$ at annual rate $t$ (as a decimal) over $n$ years:
\[ I = C \times t \times n \]
The total amount to repay is $C + I + \text{fees}$.
\end{methode}

\courssub{Démarche et corrigé commenté}

\begin{exoresolu}[Model answer -- Group ``Jeunesse Agropastorale'']
\textbf{Task 1 -- Community needs.}\\
Énoncé : identify the three priority needs and write two extra questions.
\tcblower
Solution : ``Chicken meat is the \textbf{most urgent} need because 82 \% of households mentioned it. Good-quality gari is \textbf{more important than} local jobs in our survey, but jobs are \textbf{the most useful} result for young people.'' Extra questions: ``How much do you spend on chicken every week?'' ``Would you buy gari from a youth cooperative if the price were lower?''

\textbf{Task 2 -- Best buy (comparisons).}\\
Énoncé : compare the three suppliers.
\tcblower
Solution : ``Marché Mfoundi is \textbf{the cheapest} supplier, but it is \textbf{too risky}: the chicks are not always vaccinated. \textbf{The cheaper the chicks, the higher the risk} of disease. Ferme Nkoulou is almost \textbf{as reliable as} SODEPA, and it is cheaper. SODEPA is \textbf{the most reliable} of the three, but also \textbf{the most expensive}. We choose Ferme Nkoulou because it offers the best balance between price and quality.''

\textbf{Task 3 -- Unit price calculation.}\\
Énoncé : total cost of 1 000 chicks + 40 bags of feed.
\tcblower
Solution :
\begin{align*}
\text{SODEPA} &: 1000 \times 600 + 40 \times 19500 = 600000 + 780000 = \num{1380000} \text{ FCFA}\\
\text{Ferme Nkoulou} &: 1000 \times 550 + 40 \times 18000 = 550000 + 720000 = \num{1270000} \text{ FCFA}\\
\text{Marché Mfoundi} &: 1000 \times 400 + 40 \times 16500 = 400000 + 660000 = \num{1060000} \text{ FCFA}
\end{align*}
Marché Mfoundi is the cheapest, but Ferme Nkoulou gives the \textbf{best value for money}: it costs 210 000 FCFA more, yet the chicks are vaccinated and the supplier is local.

\textbf{Task 4 -- Banking analysis.}\\
Énoncé : extract the data, compute the interest for Offer A, express a preference.
\tcblower
Solution : Offer A: rate 8 \%, guarantee = land title or guarantor, fees 25 000 FCFA. Offer B: rate 15 \%, guarantee = group solidarity, fees 10 000 FCFA.
\[ I = C \times t \times n = 1500000 \times 0,08 \times 2 = \num{240000} \text{ FCFA} \]
Total to repay with Offer A: $1500000 + 240000 + 25000 = \num{1765000}$ FCFA. With Offer B, the interest would be $1500000 \times 0,15 \times 2 = \num{450000}$ FCFA, almost double. ``\textbf{We would rather} borrow from Crédit Agricole \textbf{because} the rate is lower, \textbf{although} the decision takes longer. \textbf{We prefer} a three-month grace period \textbf{to} a fast decision.''

\textbf{Task 5 -- Formal letter (extract).}\\
Énoncé : loan application letter, 150--180 words.
\tcblower
Solution (full letter) :
\begin{quote}
Bafoussam, 15 October 2025

The General Manager\\
Crédit Agricole du Cameroun\\
Bafoussam Branch

\textbf{Subject: Loan Application -- Coopérative Agropastorale Jeunesse}

Dear Sir,

We are five young graduates from Bafoussam, \textbf{and} we have created the cooperative ``Jeunesse Agropastorale''. We are writing to apply for a loan of 1 500 000 FCFA \textbf{so that} we can buy 1 000 chicks and feed for our first production cycle. \textbf{Although} we are young, we have completed a serious market survey, \textbf{because} we want our project to succeed. \textbf{If} you accept our request, we will repay the loan in 24 months with the grace period of three months.

Our projected revenue per cycle is 3 325 000 FCFA for a production cost of 2 400 000 FCFA, leaving a net profit of 925 000 FCFA. We have chosen Ferme Nkoulou as our supplier for the best value for money.

We remain at your disposal for any further information.

Yours faithfully,\\
The Board of ``Jeunesse Agropastorale''
\end{quote}

\textbf{Task 6 -- Pitch (key sentences).}\\
Énoncé : five-minute oral pitch with adverbs of time.
\tcblower
Solution : ``We have prepared this project \textbf{for} six months. We have lived in this community \textbf{since} 2015. We finished our survey two months \textbf{ago}. \textbf{After} the first cycle, we expect a revenue of $950 \times 3500 = \num{3325000}$ FCFA, so a profit of about 925 000 FCFA. \textbf{Afterwards}, we will start repaying the loan every month.''

\textbf{Task 7 -- Jury questions (sample answers).}\\
Énoncé : answer the bankers.
\tcblower
Solution : ``\textbf{If} the price of feed increases, we \textbf{will} buy maize locally and mix our own feed. \textbf{Although} our project is new, it is not very risky \textbf{because} the demand for chicken is very high in Bafoussam.''
\end{exoresolu}

\begin{attention}[Common mistakes to avoid]
Do not say ``more cheaper'' (double marking of the comparative). Use \emph{for} with a duration and \emph{since} with a starting point: ``for six months'' but ``since January''. ``Ago'' always comes after the duration: ``two months ago''. In the formal letter, avoid contractions (write \emph{we are}, not \emph{we're}). Do not confuse ``after'' (preposition/conjunction) and ``afterwards'' (adverb): ``After the cycle'' vs ``We will repay afterwards''.
\end{attention}

\courssub{Bilan de l'activité}

This project has shown that English is a real working tool (\emph{un outil de travail}) for economic life in Cameroon. You have used the whole unit in one authentic task:

\begin{itemize}
\item \textbf{Speaking:} you gave and received information during the survey, the pitch and the jury simulation.
\item \textbf{Comparisons:} double comparatives, equality and degree helped you justify the best buy and the best bank offer.
\item \textbf{Adverbs of time:} for, since, ago, after and afterwards structured your project timeline.
\item \textbf{Complex sentences:} coordination and subordination made your letter and your answers clear and professional; expressions of preference justified your choices.
\item \textbf{Reading and writing:} you read real banking documents and produced a formal loan application letter.
\item \textbf{Calculation:} unit prices and simple interest proved that a good decision must be justified with figures (\emph{justifiée par des chiffres}).
\end{itemize}

The self-evaluation grid identifies your weak points (for example: subordination, adverbs of time, or pronunciation of numbers). The remediation session will focus on these points with short targeted exercises before the final evaluation. A well-defended business plan can really convince a bank: your English is now part of your professional future.

\newpage

\coursec{Méthodes et exercices résolus}

\courssub{Applying the lesson to real-life situations}

\begin{methode}[Comparing prices and quality to find the best buy]
To compare products and choose the best buy, follow these steps:
\begin{enumerate}
\item Identify the \cle{price} and the \cle{quality} (size, weight, durability, brand) of each product.
\item Use the correct form of comparison: \textbf{as ... as} (equality), \textbf{more ... than / -er than} (superiority), \textbf{less ... than} (inferiority), \textbf{the most / the -est} (superlative).
\item Use the \cle{double comparison} for linked changes: \emph{the cheaper, the better}.
\item Compare the \textbf{price per unit} (per kilo, per litre) when quantities differ: divide the price by the quantity.
\item Conclude with a clear recommendation and justify it with one argument of price and one of quality.
\end{enumerate}
\textbf{Key phrases:} \emph{value for money} (rapport qualité-prix), \emph{a bargain} (une bonne affaire), \emph{to be worth it} (en valoir la peine).
\end{methode}

\begin{exoresolu}[Choosing the best buy at the market]
\textbf{Task (Bac-type):} At Mokolo market, shop A sells a 5~kg bag of rice for \SI{3}{}000 FCFA and shop B sells a 10~kg bag of the same rice for \SI{5}{}500 FCFA. Write a short paragraph (5--6 lines) comparing the two offers and recommending the best buy. Use at least two comparative structures.
\tcblower
\textbf{Step 1 -- Compute the price per kilo.}
\begin{align*}
\text{Shop A: } & 3000 \div 5 = 600 \text{ FCFA per kg}\\
\text{Shop B: } & 5500 \div 10 = 550 \text{ FCFA per kg}
\end{align*}
\textbf{Step 2 -- Compare.} Shop B's rice is cheaper per kilo: $550 < 600$. The difference is $600 - 550 = 50$ FCFA per kg.
\textbf{Step 3 -- Write the paragraph with comparatives.}

\emph{Shop A sells rice at 600 FCFA per kilo, whereas shop B sells the same rice at 550 FCFA per kilo. Shop B's rice is therefore cheaper than shop A's, although the bag is bigger and heavier to carry. The quality is as good as shop A's, since it is the same brand. The more you buy at shop B, the more you save: 500 FCFA on a 10 kg bag. I recommend shop B because it offers the best value for money.}

\textbf{Conclusion:} Shop B is the best buy: the unit price is lower (550 against 600 FCFA/kg) for an equal quality.
\end{exoresolu}

\begin{methode}[Using time adverbs: for, since, ago, after, afterwards]
\begin{enumerate}
\item \cle{for} + duration (a period): \emph{for two years, for six months}. Used with the present perfect: \emph{I have banked with UBA for five years.}
\item \cle{since} + starting point (a date or event): \emph{since 2020, since last January}. Also with the present perfect.
\item \cle{ago} + past point in time, counted back from now: \emph{two days ago}. Used with the simple past, always placed \textbf{after} the time expression.
\item \cle{after} + noun or clause (order of events): \emph{After opening the account, he received a card.}
\item \cle{afterwards} (adverb, alone): \emph{We signed the contract; afterwards, we went home.}
\end{enumerate}
\textbf{Reflex:} ask ``duration or starting point?'' If duration $\rightarrow$ \textbf{for}; if starting point $\rightarrow$ \textbf{since}.
\end{methode}

\begin{exoresolu}[Completing a dialogue with time adverbs]
\textbf{Task:} Complete with \emph{for, since, ago, after} or \emph{afterwards}.
\begin{enumerate}
\item Mme Etoundi has been a customer of Afriland \dots{} 2019.
\item She opened her first savings account ten years \dots{}.
\item \dots{} filling in the form, she showed her ID card.
\item She has used mobile banking \dots{} three years.
\item He deposited the cheque and left the bank \dots{}.
\end{enumerate}
\tcblower
\textbf{Solutions and justifications:}
\begin{enumerate}
\item \textbf{since} -- 2019 is a starting point, and the verb is in the present perfect (\emph{has been}).
\item \textbf{ago} -- ``ten years'' is counted back from now; the verb \emph{opened} is in the simple past.
\item \textbf{After} -- it introduces a sequence of two actions and is followed by a verb in -ing (\emph{filling}).
\item \textbf{for} -- ``three years'' is a duration with the present perfect (\emph{has used}).
\item \textbf{afterwards} -- it stands alone at the end of the sentence, meaning ``then, after that''.
\end{enumerate}
\textbf{Conclusion:} \emph{since; ago; After; for; afterwards}.
\end{exoresolu}

\begin{methode}[Building complex sentences: coordination and subordination]
\begin{enumerate}
\item \textbf{Coordination} links two independent clauses with \cle{and} (addition), \cle{but} (contrast), \cle{or} (choice), \cle{so} (result): \emph{The bank is safe, but the fees are high.}
\item \textbf{Subordination} links a main clause and a dependent clause with \cle{because} (cause), \cle{although/though} (concession), \cle{when, while} (time), \cle{if} (condition), \cle{so that} (purpose), \cle{who/which/that} (relative).
\item Check that the dependent clause \textbf{cannot stand alone}: if it starts with \emph{because/although/when}, it needs a main clause.
\item Punctuation: when the subordinate clause comes first, add a comma: \emph{Although the fees are high, the bank is safe.}
\end{enumerate}
\end{methode}

\begin{exoresolu}[Combining sentences about banking services]
\textbf{Task:} Combine each pair of sentences using the connector in brackets.
\begin{enumerate}
\item Many Cameroonians use mobile money. It is fast and convenient. (because)
\item The interest rate is low. A savings account is still useful. (although)
\item You can withdraw money. The ATM is working. (if)
\end{enumerate}
\tcblower
\textbf{Step 1 -- Identify the logical relation:} (1) cause, (2) concession, (3) condition.
\textbf{Step 2 -- Choose the main clause} (the most important idea) and attach the subordinate clause.
\textbf{Step 3 -- Write and punctuate.}
\begin{enumerate}
\item \emph{Many Cameroonians use mobile money \textbf{because} it is fast and convenient.} (No comma: the subordinate clause comes second.)
\item \emph{\textbf{Although} the interest rate is low, a savings account is still useful.} (Comma after the fronted subordinate clause.)
\item \emph{You can withdraw money \textbf{if} the ATM is working.}
\end{enumerate}
\textbf{Conclusion:} each sentence now has one main clause and one subordinate clause, correctly linked and punctuated.
\end{exoresolu}

\begin{methode}[Expressing preferences]
\begin{enumerate}
\item \cle{prefer + noun/-ing + to}: \emph{I prefer cash to cheques. / I prefer saving to borrowing.}
\item \cle{would rather + base verb + than}: \emph{I would rather save than borrow.} (No -ing, no \emph{to} before the verb.)
\item \cle{would prefer + to-infinitive}: \emph{I would prefer to open a joint account.}
\item \cle{like ... better than / favourite}: \emph{I like mobile banking better than queuing at the counter.}
\item To justify a preference, add \emph{because} + reason: \emph{... because it saves time.}
\end{enumerate}
\textbf{Reflex:} after \emph{would rather}, never write \emph{to} or -ing.
\end{methode}

\begin{exoresolu}[Expressing your banking preferences]
\textbf{Task:} Rewrite using the structure in brackets.
\begin{enumerate}
\item I like mobile money more than bank cheques. (prefer)
\item She wants to save money, not to take a loan. (would rather ... than)
\item We want to use an ATM. (would prefer)
\end{enumerate}
\tcblower
\textbf{Solutions:}
\begin{enumerate}
\item \emph{I \textbf{prefer} mobile money \textbf{to} bank cheques.} -- Structure: prefer + noun + to + noun.
\item \emph{She \textbf{would rather} save money \textbf{than} take a loan.} -- Base verbs \emph{save} and \emph{take}, no \emph{to}, no -ing.
\item \emph{We \textbf{would prefer to use} an ATM.} -- would prefer + to-infinitive.
\end{enumerate}
\textbf{Conclusion:} the three structures are \emph{prefer X to Y}, \emph{would rather V than V}, \emph{would prefer to V}.
\end{exoresolu}

\begin{methode}[Writing a composition about banking services]
\begin{enumerate}
\item \textbf{Plan} in three parts: introduction (present the topic), body (2--3 arguments with examples from Cameroon: mobile money, savings accounts, loans, ATMs), conclusion (your opinion).
\item Use \textbf{linkers}: \emph{first, moreover, however, for example, in addition, to conclude}.
\item Include at least one \textbf{comparison}, one \textbf{complex sentence} (subordination) and one \textbf{preference}.
\item Use the vocabulary of the unit: \emph{account, deposit, withdrawal, loan, interest, fees, transfer, balance}.
\item \textbf{Proofread}: subject-verb agreement, verb tenses, spelling of numbers and currency (FCFA).
\end{enumerate}
\end{methode}

\begin{exoresolu}[Guided composition: ``Why I opened a bank account'']
\textbf{Task (Bac-type):} Write a 120--150 word composition entitled ``Why I opened a bank account''. Use at least one comparative, one subordinate clause and one expression of preference.
\tcblower
\textbf{Model answer:}

\emph{Last year, I opened a savings account at a microfinance bank in Yaoundé because I wanted to keep my money safe. First, a bank account is safer than keeping cash at home, where it can be stolen. Moreover, my savings earn interest, so my money grows little by little. Although the bank charges small fees, the service is worth it. I prefer mobile banking to queuing at the counter, since it is faster and I can check my balance at any time. For example, I transfer money to my mother in the village in one minute. The more I save, the more confident I feel about my future. To conclude, opening a bank account was one of the best decisions of my life, and I advise every young Cameroonian to do the same.} (128 words)

\textbf{Checklist:} comparative (\emph{safer than, faster}), subordination (\emph{because, although, since}), preference (\emph{I prefer ... to}), double comparison (\emph{the more ... the more}), linkers (\emph{first, moreover, for example, to conclude}).
\end{exoresolu}

\begin{methode}[Assessing consumer needs and community resources (speaking)]
\begin{enumerate}
\item \textbf{Ask} for information with polite questions: \emph{What do you need most? How much can you spend? Which services are available in your neighbourhood?}
\item \textbf{Give} information clearly: \emph{Our community has a health centre, two schools and a market, but there is no clean water supply.}
\item \textbf{Assess} needs by ranking them: \emph{The most urgent need is... / Water is more important than...}
\item \textbf{Suggest} solutions with \emph{should, could, why don't we...}: \emph{The council should build a borehole.}
\item \textbf{Conclude} the exchange by summarising the decision.
\end{enumerate}
\end{methode}

\begin{exoresolu}[Dialogue: assessing community needs]
\textbf{Task:} Complete the dialogue between a development worker (W) and a villager (V).
\begin{enumerate}
\item W: \dots{} (ask what the village needs most)
\item V: We need clean water more than anything else.
\item W: \dots{} (ask if there is a health centre)
\item V: Yes, but it is smaller and less equipped than the one in town.
\item W: \dots{} (suggest building a borehole)
\end{enumerate}
\tcblower
\textbf{Possible answers:}
\begin{enumerate}
\item \emph{What does your village need most?} -- open question with \emph{What ... most}.
\item \emph{Is there a health centre in your village?} -- yes/no question, inversion of \emph{there is}.
\item \emph{Why don't we build a borehole near the market?} or \emph{The community should build a borehole.} -- suggestion with \emph{why don't we} or \emph{should}.
\end{enumerate}
\textbf{Conclusion:} the exchange assesses needs (water first), compares resources (\emph{smaller, less equipped than}) and proposes a solution -- exactly the speaking function required by the programme.
\end{exoresolu}

\begin{attention}[Les 5 erreurs qui coûtent des points au Bac]
\begin{enumerate}
\item \textbf{Double marque de comparaison:} ne jamais écrire \emph{more cheaper} ou \emph{the most easiest}. On choisit \emph{cheaper} OU \emph{more expensive}, jamais les deux ensemble.
\item \textbf{Confondre \emph{for} et \emph{since}:} \emph{for} + durée (\emph{for two years}), \emph{since} + point de départ (\emph{since 2022}). Et \emph{ago} se place \textbf{après} la durée avec le prétérit: \emph{two years ago}, jamais \emph{since two years ago}.
\item \textbf{\emph{Would rather} suivi de \emph{to} ou de -ing:} la bonne forme est \emph{would rather + base verbale}: \emph{I would rather save} (pas \emph{to save}, pas \emph{saving}).
\item \textbf{Phrase subordonnée laissée seule:} \emph{Because the bank was closed.} n'est pas une phrase complète. Il faut une proposition principale: \emph{Because the bank was closed, I used the ATM.}
\item \textbf{Comparaison de prix sans ramener à l'unité:} comparer deux prix sans calculer le prix au kilo ou au litre fausse le raisonnement. Toujours diviser le prix par la quantité avant de conclure, et justifier la recommandation par un argument de prix ET un argument de qualité.
\end{enumerate}
\end{attention}

\newpage

\coursec{Exercices}

\courssub{Je vérifie mes connaissances}

\begin{exercice}[QCM : Vocabulary \& Grammar Basics]
Choose the correct option (a, b, or c) to complete each sentence.

\begin{enumerate}
    \item The new smartphone is \_\_\_\_\_\_ expensive \_\_\_\_\_\_ the old model. (Equality comparison)
    \begin{itemize}
        \item[a)] as / than
        \item[b)] so / as
        \item[c)] as / as
    \end{itemize}

    \item We have been waiting for the bank manager \_\_\_\_\_\_ 10:00 a.m. (Adverb of time)
    \begin{itemize}
        \item[a)] since
        \item[b)] for
        \item[c)] ago
    \end{itemize}

    \item A \_\_\_\_\_\_ is an asset that a borrower offers to a lender to secure a loan. (Banking vocabulary)
    \begin{itemize}
        \item[a)] overdraft
        \item[b)] collateral
        \item[c)] interest rate
    \end{itemize}

    \item \_\_\_\_\_\_ the high interest rates, many small businesses still apply for loans. (Subordination / Concession)
    \begin{itemize}
        \item[a)] Because
        \item[b)] Although
        \item[c)] So
    \end{itemize}

    \item I \_\_\_\_\_\_ save money in a credit union rather than keep it at home. (Preference)
    \begin{itemize}
        \item[a)] would rather
        \item[b)] prefer to
        \item[c)] like
    \end{itemize}
\end{enumerate}
\end{exercice}

\begin{exercice}[True / False Justified]
Read the statements below. Write \textbf{True} or \textbf{False} and justify your answer with a sentence from the lesson or your knowledge.

\begin{enumerate}
    \item Assessing consumer needs means identifying what the seller wants to produce.
    \item Community resources include only financial institutions like banks.
    \item ``Double comparison'' structure (The more..., the more...) is used to show a proportional relationship.
    \item A savings account generally offers a higher interest rate than a current (checking) account.
    \item ``Since'' is used with a duration (e.g., since three years) and ``for'' with a starting point (e.g., for 2020).
\end{enumerate}
\end{exercice}

\begin{exercice}[Definitions \& Key Concepts]
Define the following terms in your own words (in English). Give a short example for each.

\begin{enumerate}
    \item Consumer needs assessment (\emph{évaluation des besoins des consommateurs})
    \item Collateral (\emph{garantie / collatéral})
    \item Overdraft (\emph{découvert bancaire})
    \item Double comparison (\emph{comparaison double / proportionnelle})
    \item Mobile banking (\emph{banque mobile})
\end{enumerate}
\end{exercice}

\begin{exercice}[Grammar Quick Practice]
Rewrite the following sentences using the words in brackets. Do not change the meaning.

\begin{enumerate}
    \item The loan term is long. The interest paid is high. \quad (\textbf{The longer... the higher...})
    \item This bank offers good services. That bank offers better services. \quad (\textbf{not as... as})
    \item He opened his account five years ago. He still uses it. \quad (\textbf{for / since})
    \item She has a low income. She got a microfinance loan. \quad (\textbf{Although / Despite})
    \item I think saving in a bank is better than mobile money. \quad (\textbf{would rather})
\end{enumerate}
\end{exercice}

\courssub{Je m'entraîne}

\begin{exercice}[Reading \& Table Completion: Smart Shopping]
Read the text comparing two laptop offers for a student association. Complete the comparison table below.

\textbf{Text:} \emph{The association ``Tech Future'' needs a laptop. \textbf{Offer A (Shop X)}: Price 450,000 FCFA. Warranty: 2 years. Specs: 8GB RAM, 256GB SSD, Core i5. \textbf{Offer B (Shop Y)}: Price 400,000 FCFA. Warranty: 1 year. Specs: 4GB RAM, 500GB HDD, Core i3. The treasurer thinks Offer A is the best buy because the processor is faster and the SSD is more reliable than HDD, even if it is more expensive.}

\begin{center}
\begin{tabularx}{\linewidth}{|l|X|X|}
\hline
\textbf{Criteria} & \textbf{Offer A (Shop X)} & \textbf{Offer B (Shop Y)} \\ \hline
Price (FCFA) & & \\ \hline
Warranty & & \\ \hline
Processor (CPU) & & \\ \hline
RAM & & \\ \hline
Storage Type / Size & & \\ \hline
\textbf{Best Buy? (Justify in 1 sentence)} & \multicolumn{2}{c|}{} \\ \hline
\end{tabularx}
\end{center}
\end{exercice}

\begin{exercice}[Grammar: Complex Sentences \& Preferences]
Combine the pairs of sentences using the linking word in brackets. Make any necessary changes.

\begin{enumerate}
    \item The interest rate dropped. Many people borrowed money. \quad (\textbf{As / Because})
    \item The client signed the contract. He did not read the terms. \quad (\textbf{without / -ing})
    \item You must provide collateral. The bank will not approve the loan. \quad (\textbf{Unless})
    \item Mobile money is convenient. Traditional banks offer more security. \quad (\textbf{While / Whereas})
    \item I want to open a current account. I want a debit card. \quad (\textbf{not only... but also})
    \item He prefers investing in land. He does not trust the stock market. \quad (\textbf{prefer... to... / would rather... than})
\end{enumerate}
\end{exercice}

\begin{exercice}[Vocabulary in Context: Banking Services]
Fill in the blanks with the appropriate word from the box. Two words are extra.

\begin{center}
\textbf{overdraft \quad interest rate \quad collateral \quad withdrawal \quad deposit \quad statement \quad loan \quad savings account \quad standing order \quad mobile banking}
\end{center}

\begin{enumerate}
    \item Every month, the bank sends me a \_\_\_\_\_\_ showing all my transactions.
    \item To get a large \_\_\_\_\_\_, the entrepreneur used his land as \_\_\_\_\_\_.
    \item I set up a \_\_\_\_\_\_ to pay my rent automatically on the 1st of each month.
    \item The \_\_\_\_\_\_ on this savings account is 3.5\% per year.
    \item She made a large \_\_\_\_\_\_ yesterday to pay her suppliers.
    \item With \_\_\_\_\_\_, I can check my balance on my phone without going to the agency.
    \item Be careful! If you spend more than you have, you will go into \_\_\_\_\_\_ and pay fees.
\end{enumerate}
\end{exercice}

\begin{exercice}[Speaking / Writing: Assessing Community Resources]
\textbf{Scenario:} You are a member of a Youth Agricultural Cooperative in your village. You need to assess the community resources available to start a poultry farm.

\textbf{Task:} Write a short report (80--100 words) answering these questions:
\begin{enumerate}
    \item What are the \textbf{human resources} available (skills, labour)?
    \item What \textbf{financial services} (banks, microfinance, tontines) exist nearby?
    \item What \textbf{material resources} (land, water, feed suppliers) are accessible?
    \item What is the \textbf{main need} you must request help for (training, funding, equipment)?
\end{enumerate}
Use linking words: \emph{Firstly, Moreover, However, Therefore}.
\end{exercice}

\courssub{Je me prépare au Bac}

\begin{exercice}[Bac Type: Formal Email -- Banking Services]
\textbf{Context:} You are the Secretary General of the ``Young Entrepreneurs Association'' (YEA) in Douala. Your association wants to open a \textbf{savings account} to manage membership fees and project funds.

\textbf{Task:} Write a \textbf{formal email} to the Manager of \textbf{Afriland First Bank}, Douala Branch.
\begin{itemize}
    \item \textbf{Subject line:} Request for Information -- Opening Association Savings Account.
    \item \textbf{Body:} Introduce the association (legal status, number of members, purpose).
    \item Ask specific questions about: \textbf{minimum opening deposit}, \textbf{monthly maintenance fees}, \textbf{interest rate} on savings, \textbf{conditions for withdrawal}, and \textbf{mobile banking / USSD options} for members in remote areas.
    \item Request an appointment to finalize the procedure.
    \item Use formal register: \emph{Dear Sir/Madam, I am writing to inquire..., Could you please provide..., We would appreciate..., Yours faithfully.}
\end{itemize}
\textbf{Word count:} 200--250 words.
\end{exercice}

\begin{exercice}[Bac Type: Grammar Synthesis \& Reading Comprehension]
\textbf{Part A: Grammar Transformation (5 marks)}
Rewrite the sentences as instructed.

\begin{enumerate}
    \item The repayment period is long. The total cost of the credit is high. \quad (Begin with: \textbf{The longer...})
    \item The microfinance institution approved the loan. The applicant had no collateral. \quad (Begin with: \textbf{Despite...})
    \item ``When did you open this account?'' the clerk asked me. \quad (Reported speech: \textbf{The clerk asked me...})
    \item I started using mobile money in 2019. I still use it daily. \quad (Use: \textbf{since / for} in two separate sentences)
    \item My father prefers keeping cash at home. He does not trust banks. \quad (Use: \textbf{would rather... than})
\end{enumerate}

\textbf{Part B: Reading Comprehension (10 marks)}
Read the extract from a CRTV ``Eco d'ici'' report on agricultural financing and answer the questions.

\emph{``In the Centre Region, the `Green Valley' cooperative struggled to expand their maize production. Local banks required land titles as collateral, which most members lacked. Then, a partner NGO introduced a `Warehouse Receipt System'. Farmers store harvested maize in a certified warehouse. They receive a receipt proving ownership and quantity. This receipt serves as collateral for a bank loan up to 70\% of the crop's value. The interest rate is 5\% per annum, subsidized by the state. The loan duration is 6 months, matching the storage period. The advantage: farmers sell when prices are high, not at harvest. The risk: if prices drop, the sale might not cover the loan.''}

\begin{enumerate}
    \item Why couldn't the cooperative get a traditional bank loan initially? (2 marks)
    \item Explain how the Warehouse Receipt System works. (3 marks)
    \item What are the \textbf{loan conditions} (rate, duration, percentage of value)? (3 marks)
    \item State \textbf{one advantage} and \textbf{one risk} mentioned in the text. (2 marks)
\end{enumerate}
\end{exercice}

\begin{exercice}[Bac Type: Argumentative Essay + Practical Case]
\textbf{Part A: Argumentative Essay (10 marks)}
\textbf{Topic:} \emph{``Mobile money services (Orange Money, MTN MoMo) will completely replace traditional banks in Cameroon within the next ten years.''}

\textbf{Instructions:} Write a structured essay (Introduction, Body, Conclusion).
\begin{itemize}
    \item \textbf{Introduction:} Hook, define the topic, announce plan.
    \item \textbf{Body Paragraph 1 (For):} 2 arguments why mobile money could replace banks (accessibility, low cost, financial inclusion for unbanked).
    \item \textbf{Body Paragraph 2 (Against):} 2 arguments why banks remain essential (large loans, security/regulation, business services, international trade).
    \item \textbf{Conclusion:} Summarize, give \textbf{your personal opinion} (nuanced).
\end{itemize}
\textbf{Requirements:} Use linking words (\emph{Furthermore, However, On the one hand, Conversely, In conclusion}). Use vocabulary: \emph{unbanked population, transaction fees, creditworthiness, regulatory framework, financial inclusion}. \textbf{Length:} 250--300 words.

\textbf{Part B: Practical Case -- Bank Statement Analysis (10 marks)}
Below is a fictitious bank statement for ``Mme Ngo Bakang'' (Current Account n° 1002345) for January 2026.

\begin{center}
\begin{tabular}{|c|l|r|r|r|}
\hline
\textbf{Date} & \textbf{Description} & \textbf{Debit (FCFA)} & \textbf{Credit (FCFA)} & \textbf{Balance (FCFA)} \\ \hline
01/01 & Opening Balance & & & 150,000 \\ \hline
05/01 & Salary Deposit & & 300,000 & 450,000 \\ \hline
10/01 & ATM Withdrawal & 50,000 & & 400,000 \\ \hline
12/01 & Transfer to Supplier (Invoice 44) & 120,000 & & 280,000 \\ \hline
15/01 & \textbf{Bank Fees (Monthly)} & \textbf{5,500} & & \textbf{274,500} \\ \hline
20/01 & POS Payment Supermarket & 45,000 & & 229,500 \\ \hline
25/01 & \textbf{Interest Credit} & & \textbf{1,200} & \textbf{230,700} \\ \hline
28/01 & \textbf{Unknown Debit ``COM*''} & \textbf{15,000} & & \textbf{215,700} \\ \hline
31/01 & Closing Balance & & & 215,700 \\ \hline
\end{tabular}
\end{center}

\textbf{Questions:}
\begin{enumerate}
    \item Verify the \textbf{final balance} (215,700 FCFA) by calculating step by step. Show your work. (3 marks)
    \item Identify \textbf{two anomalies} or suspicious entries in the statement. Explain why. (4 marks)
    \item Write a \textbf{formal complaint letter} (email format) to the Bank Manager regarding the anomalies.
        \begin{itemize}
            \item Use: \emph{``I would like to draw your attention to...''}, \emph{``I would rather...''}, \emph{``Unless this is rectified...''}.
            \item Request a reversal and an explanation.
        \end{itemize}
    (3 marks)
\end{enumerate}
\end{exercice}



\newpage

\coursec{Corrigés des exercices}

\begin{corrige}[Exercice 1 — QCM : Vocabulary \& Grammar Basics]
\begin{enumerate}
    \item \textbf{c)} \cle{as / as}. Equality comparison: \emph{as + adjective + as}. ``So... as'' is only used in negative sentences (\emph{not so expensive as}).
    \item \textbf{a)} \cle{since}. \emph{Since} + starting point (10:00 a.m.). \emph{For} + duration; \emph{ago} is placed after the duration (\emph{two hours ago}).
    \item \textbf{b)} \cle{collateral} (\emph{garantie}). An \cle{overdraft} is a negative balance; an \cle{interest rate} is the cost of borrowing.
    \item \textbf{b)} \cle{Although}. Concession: the result is unexpected. ``Because'' would express cause, ``so'' consequence.
    \item \textbf{a)} \cle{would rather}. Structure: \emph{would rather + base verb ... rather than + base verb}. ``Prefer to'' cannot be followed by ``rather than'' in this pattern.
\end{enumerate}
\textbf{Points de vigilance :} ne pas confondre \emph{since} (point de départ) et \emph{for} (durée) ; la structure \emph{would rather + BV} exige le verbe de base, sans \emph{to}.
\end{corrige}

\begin{corrige}[Exercice 2 — True / False Justified]
\begin{enumerate}
    \item \textbf{False.} Assessing \cle{consumer needs} means identifying what \emph{consumers} (the community) want or need, not what the seller wants to produce.
    \item \textbf{False.} \cle{Community resources} also include human resources (skills, labour), material resources (land, water, markets) and social services (schools, health centres), not only banks.
    \item \textbf{True.} \emph{The longer the loan term, the higher the total interest} shows that two quantities change together proportionally (\cle{double comparative}).
    \item \textbf{True.} \cle{Savings accounts} are remunerated (e.g., around 3\% per year in many Cameroonian institutions), while \cle{current accounts} usually earn little or no interest.
    \item \textbf{False.} It is the opposite: \emph{since} + starting point (\emph{since 2020}), \emph{for} + duration (\emph{for three years}).
\end{enumerate}
\textbf{Points de vigilance :} la justification est exigée pour chaque item ; une réponse ``True/False'' seule ne vaut pas la totalité des points.
\end{corrige}

\begin{corrige}[Exercice 3 — Definitions \& Key Concepts]
Expected elements (accept any correct formulation):
\begin{enumerate}
    \item \cle{Assessing consumer needs}: The process of finding out what goods or services the people of a community want, need and can pay for. \emph{Example: before opening a shop, the trader asks households which products they buy most.}
    \item \cle{Collateral}: A valuable asset (land title, vehicle, warehouse receipt) that a borrower gives the bank as security for a loan; the bank can seize it if the loan is not repaid. \emph{Example: a farmer uses his land title as collateral.}
    \item \cle{Overdraft}: A situation where an account holder spends more money than the balance; the bank allows a negative balance but charges fees and interest. \emph{Example: his account shows $-20\,000$ FCFA, so he is in overdraft.}
    \item \cle{Double comparative}: A structure with two comparatives showing that two things change together: \emph{the + comparative..., the + comparative...}. \emph{Example: The more you save, the richer you become.}
    \item \cle{Mobile banking}: Banking operations (balance check, transfer, payment) done on a mobile phone, by app or USSD code, without going to the agency. \emph{Example: she pays her supplier by mobile banking from her village.}
\end{enumerate}
\textbf{Points de vigilance :} une définition complète = nature/fonction + exemple ; éviter les définitions circulaires (``an overdraft is when you are in overdraft'').
\end{corrige}

\begin{corrige}[Exercice 4 — Grammar Quick Practice]
\begin{enumerate}
    \item The longer the loan term, the higher the interest paid.
    \item This bank does not offer services as good as that bank. (ou : This bank's services are not as good as that bank's.)
    \item He has used his account for five years. / He has used his account since 2021 (five years ago).
    \item Although she has a low income, she got a microfinance loan. / Despite her low income, she got a microfinance loan.
    \item I would rather save money in a bank than use mobile money.
\end{enumerate}
\textbf{Points de vigilance :} avec \emph{since/for}, utiliser le present perfect (\emph{has used}), pas le prétérit ; \emph{despite} est suivi d'un nom ou d'un V-ing, jamais d'une phrase complète avec sujet + verbe conjugué.
\end{corrige}

\begin{corrige}[Exercice 5 — Reading \& Table Completion: Smart Shopping]
\begin{center}
\begin{tabularx}{\linewidth}{|l|X|X|}
\hline
\rowcolor{popblueL}
\textbf{Criteria} & \textbf{Offer A (Shop X)} & \textbf{Offer B (Shop Y)} \\ \hline
Price (FCFA) & 450,000 & 400,000 \\ \hline
Warranty & 2 years & 1 year \\ \hline
Processor (CPU) & Core i5 (faster) & Core i3 (slower) \\ \hline
RAM & 8 GB & 4 GB \\ \hline
Storage Type / Size & 256 GB SSD (more reliable) & 500 GB HDD (larger but slower) \\ \hline
\rowcolor{popgreenL}
\textbf{Best Buy?} & \multicolumn{2}{l|}{Offer A: although it is 50,000 FCFA more expensive, it is faster,} \\
 & \multicolumn{2}{l|}{more reliable and has a longer warranty, so it is better value for money.} \\ \hline
\end{tabularx}
\end{center}
Vérification du calcul d'écart : $450\,000 - 400\,000 = 50\,000$ FCFA, soit $\dfrac{50\,000}{400\,000} \times 100 = 12,5\%$ de plus que l'offre B.
\textbf{Points de vigilance :} la justification doit comparer prix \emph{et} qualité (garantie, performance) ; un ``best buy'' n'est pas forcément le moins cher, mais le meilleur rapport qualité-prix.
\end{corrige}

\begin{corrige}[Exercice 6 — Grammar: Complex Sentences \& Preferences]
\begin{enumerate}
    \item As / Because the interest rate dropped, many people borrowed money. (\cle{Subordination of cause})
    \item The client signed the contract without reading the terms. (\cle{Preposition + V-ing})
    \item Unless you provide collateral, the bank will not approve the loan. (\cle{Condition: unless = if not})
    \item While / Whereas mobile money is convenient, traditional banks offer more security. (\cle{Contrast})
    \item I want not only to open a current account but also to get a debit card. (ou : I want to open not only a current account but also a debit card.) (\cle{Correlative coordination})
    \item He prefers investing in land to investing in the stock market. / He would rather invest in land than trust the stock market. (\cle{Preference structures})
\end{enumerate}
\textbf{Points de vigilance :} \emph{unless} = ``if... not'', donc ne pas ajouter une négation dans la subordonnée ; \emph{prefer V-ing to V-ing} (les deux verbes en -ing) ; \emph{without} + V-ing.
\end{corrige}

\begin{corrige}[Exercice 7 — Vocabulary in Context: Banking Services]
\begin{enumerate}
    \item \cle{statement} (\emph{relevé de compte})
    \item \cle{loan} ; \cle{collateral}
    \item \cle{standing order} (\emph{ordre de virement permanent})
    \item \cle{interest rate}
    \item \cle{withdrawal} (\emph{retrait}) --- elle sort de l'argent pour payer ses fournisseurs
    \item \cle{mobile banking}
    \item \cle{overdraft}
\end{enumerate}
Extra words: \textbf{deposit}, \textbf{savings account}.
\textbf{Points de vigilance :} distinguer \emph{withdrawal} (retrait) et \emph{deposit} (dépôt) ; \emph{standing order} = paiement automatique régulier, à ne pas confondre avec un simple transfer.
\end{corrige}

\begin{corrige}[Exercice 8 — Speaking / Writing: Assessing Community Resources]
\textbf{Model answer (about 95 words):}

\emph{Firstly, our cooperative has strong human resources: twelve active young members, and two of them have experience in poultry keeping. Moreover, the village has a microfinance office and a weekly tontine, so small loans are possible. Concerning material resources, we own a plot of land near a stream, and a feed supplier operates in the neighbouring town. However, we have no training in disease prevention and no money to buy chicks and vaccines. Therefore, our main need is funding and technical training. We request support from the MINADER delegation and a microfinance loan.}

\textbf{Barème indicatif :} contenu (4 questions traitées) : 4 pts ; connecteurs exigés (\emph{Firstly, Moreover, However, Therefore}) : 2 pts ; vocabulaire de l'unité : 2 pts ; grammaire et orthographe : 2 pts.
\textbf{Points de vigilance :} respecter la fourchette 80--100 mots ; répondre aux quatre questions dans l'ordre ; employer les quatre connecteurs demandés.
\end{corrige}

\begin{corrige}[Exercice 9 — Bac Type: Formal Email -- Banking Services]
\textbf{Model answer (about 220 words):}

\emph{Subject: Request for Information -- Opening Association Savings Account}

\emph{Dear Sir/Madam,}

\emph{I am writing to inquire about the conditions for opening a savings account at your Douala branch. I am the Secretary General of the Young Entrepreneurs Association (YEA), a legally registered association with forty-five active members. Our purpose is to promote youth entrepreneurship through training and small income-generating projects.}

\emph{We would like to open a savings account to manage membership fees and project funds safely. Could you please provide us with the following information? Firstly, what is the minimum opening deposit required for an association account? Secondly, what are the monthly maintenance fees? Thirdly, what interest rate do you offer on savings, and how often is the interest paid? Moreover, we would like to know the conditions for withdrawal: is a notice period necessary, and are there limits on the amount? Finally, since some of our members live in remote areas, we would appreciate details about your mobile banking and USSD services, including transaction fees.}

\emph{We would be grateful if you could also indicate the documents required, such as our registration certificate and the minutes of our general assembly.}

\emph{Could we arrange an appointment next week to finalize the procedure? Thank you in advance for your assistance.}

\emph{Yours faithfully,}

\emph{[Name], Secretary General, Young Entrepreneurs Association, Douala.}

\textbf{Barème indicatif (20 pts) :} format email (objet, salutation, formule finale) : 3 pts ; présentation de l'association : 3 pts ; les 5 questions précises posées : 5 pts ; demande de rendez-vous : 2 pts ; registre formel et expressions imposées : 4 pts ; exactitude linguistique : 3 pts.
\textbf{Points de vigilance :} \emph{Yours faithfully} (et non \emph{sincerely}) quand le destinataire n'est pas nommé ; registre formel : pas de contractions (\emph{I am}, pas \emph{I'm}) ; toutes les questions de la consigne doivent apparaître.
\end{corrige}

\begin{corrige}[Exercice 10 — Bac Type: Grammar Synthesis \& Reading Comprehension]
\textbf{Part A (5 marks, 1 mark each):}
\begin{enumerate}
    \item The longer the repayment period, the higher the total cost of the credit.
    \item Despite having no collateral, the applicant got the loan approved by the microfinance institution. (ou : Despite the fact that the applicant had no collateral, the microfinance institution approved the loan.)
    \item The clerk asked me when I had opened that account.
    \item I have used mobile money since 2019. / I have used mobile money for seven years (since 2019).
    \item My father would rather keep cash at home than trust banks.
\end{enumerate}
\textbf{Points de vigilance (Part A) :} en discours rapporté, recul des temps (\emph{did you open} $\rightarrow$ \emph{had opened}) et \emph{this} $\rightarrow$ \emph{that} ; avec \emph{since/for}, present perfect obligatoire.

\textbf{Part B (10 marks):}
\begin{enumerate}
    \item Because local banks required land titles as collateral, and most cooperative members did not have land titles. (2 marks)
    \item Farmers store their harvested maize in a certified warehouse (1 mark); they receive a receipt proving ownership and quantity (1 mark); this receipt is used as collateral to obtain a bank loan (1 mark).
    \item Interest rate: 5\% per annum, subsidized by the state (1 mark); duration: 6 months, matching the storage period (1 mark); amount: up to 70\% of the crop's value (1 mark).
    \item Advantage: farmers can sell when prices are high instead of selling cheaply at harvest (1 mark). Risk: if prices drop, the sale of the maize might not cover the loan repayment (1 mark).
\end{enumerate}
\textbf{Points de vigilance (Part B) :} répondre en phrases complètes en s'appuyant sur le texte ; ne pas recopier tout le passage ; pour la question 3, citer les \emph{trois} conditions (taux, durée, pourcentage).
\end{corrige}

\begin{corrige}[Exercice 11 — Bac Type: Argumentative Essay + Practical Case]
\textbf{Part A -- Marking guide and model outline (10 marks):}

\emph{Introduction:} In Cameroon, mobile money has grown spectacularly: millions of people now send and receive money with a simple telephone. Some observers claim that services like Orange Money and MTN Mobile Money will completely replace traditional banks within ten years. This essay will examine both sides of the argument before giving a personal opinion.

\emph{Body 1 (For):} On the one hand, mobile money is accessible everywhere, even in villages where there is no bank agency; a simple kiosk is enough. Furthermore, transaction fees are low and opening an account takes a few minutes, which favours the financial inclusion of the unbanked population.

\emph{Body 2 (Against):} However, traditional banks remain essential. Only banks can grant large loans, because they assess the creditworthiness of borrowers and operate under a strict regulatory framework. Conversely, mobile money offers neither important credit nor international trade services such as letters of credit, which businesses need.

\emph{Conclusion:} In conclusion, although mobile money will keep growing, it will probably \emph{complement} rather than replace banks, because the two systems answer different needs.

\textbf{Barème indicatif (Part A) :} structure Intro/Body/Conclusion : 2 pts ; 2 arguments ``for'' + 2 arguments ``against'' : 4 pts ; connecteurs et vocabulaire imposés : 2 pts ; opinion personnelle nuancée et exactitude linguistique : 2 pts.
\textbf{Points de vigilance :} respecter 250--300 mots ; donner une opinion \emph{personnelle} en conclusion ; employer au moins trois des mots imposés (\emph{unbanked population, transaction fees, creditworthiness, regulatory framework, financial inclusion}).

\textbf{Part B (10 marks):}
\begin{enumerate}
    \item Step-by-step verification:
    \begin{align*}
    150\,000 + 300\,000 &= 450\,000 \\
    450\,000 - 50\,000 &= 400\,000 \\
    400\,000 - 120\,000 &= 280\,000 \\
    280\,000 - 5\,500 &= 274\,500 \\
    274\,500 - 45\,000 &= 229\,500 \\
    229\,500 + 1\,200 &= 230\,700 \\
    230\,700 - 15\,000 &= 215\,700 \text{ FCFA}
    \end{align*}
    The final balance of 215,700 FCFA is \textbf{correct}. (3 marks)
    \item Two anomalies (any two, 2 marks each):
    \begin{itemize}
        \item The \textbf{unknown debit ``COM*'' of 15,000 FCFA} on 28/01: the description is unclear and the customer did not authorize it; it may be a fraudulent or erroneous debit.
        \item The \textbf{monthly bank fees of 5,500 FCFA} are unusually high for a current account and should be justified by the bank.
        \item The \textbf{interest credit of 1,200 FCFA} on a current account is surprising, since current accounts normally earn no interest; it needs an explanation.
    \end{itemize}
    \item \textbf{Model complaint email (3 marks):}

    \emph{Subject: Complaint -- Unauthorized Debit on Account n° 1002345}

    \emph{Dear Sir/Madam,}

    \emph{I am writing to complain about anomalies on my January 2026 statement. I would like to draw your attention to an unknown debit labelled ``COM*'' of 15,000 FCFA on 28/01, which I never authorized, and to monthly fees of 5,500 FCFA that seem excessive. I would rather have these amounts reversed immediately and receive a written explanation. Unless this is rectified within ten days, I will be obliged to refer the matter to the banking regulator. I look forward to your prompt reply.}

    \emph{Yours faithfully,}

    \emph{Mme Ngo Bakang.}
\end{enumerate}
\textbf{Points de vigilance (Part B) :} question 1 : montrer \emph{toutes} les étapes du calcul, pas seulement le résultat ; question 2 : identifier l'anomalie \emph{et} expliquer pourquoi elle est suspecte ; question 3 : utiliser les trois expressions imposées (\emph{I would like to draw your attention to...}, \emph{I would rather...}, \emph{Unless this is rectified...}) et demander explicitement le remboursement (reversal) et une explication.
\end{corrige}

\newpage

\coursec{Fiche bilan}

\begin{popfigure}
\begin{tikzpicture}[
    mindmap,
    concept color=popblue,
    level 1 concept/.append style={level distance=4.5cm, sibling angle=360/7},
    level 2 concept/.append style={level distance=3cm},
    every node/.style={font=\small, text=white, align=center},
    branch1/.style={concept color=poporange, grow=30},
    branch2/.style={concept color=popgreen, grow=90},
    branch3/.style={concept color=poppurple, grow=150},
    branch4/.style={concept color=poppink, grow=210},
    branch5/.style={concept color=popgold, grow=270},
    branch6/.style={concept color=popdark, grow=330},
    branch7/.style={concept color=popblue!80!black, grow=-30}
]
\node[concept, text width=3.2cm, minimum size=3.2cm] {Economic Life\\ \& Occupations}
    child[branch1] { node[concept, text width=2.8cm] {1. Assessing Consumer Needs\\ \& Community Resources\\ \textit{(Speaking/Vocab)} } }
    child[branch2] { node[concept, text width=2.8cm] {2. Smart Shopping:\\ Comparing Prices \& Quality\\ \textit{(Listening/Vocab/Grammar: Adverbs)} } }
    child[branch3] { node[concept, text width=2.8cm] {3. Banking Services\\ in Cameroon\\ \textit{(Reading/Vocab/Writing)} } }
    child[branch4] { node[concept, text width=2.8cm] {4. Grammar:\\ Forms of Comparison\\ \textit{(Double, Equality, Degree)} } }
    child[branch5] { node[concept, text width=2.8cm] {5. Grammar:\\ Complex Sentences\\ \textit{(Coordination \& Subordination)} } }
    child[branch6] { node[concept, text width=2.8cm] {6. Grammar:\\ Expressing Preferences\\ \textit{(Would rather, prefer, etc.)} } }
    child[branch7] { node[concept, text width=2.8cm] {7. Integration:\\ Evaluation \& Remediation\\ \textit{(Real-life Tasks)} } };
\end{tikzpicture}
\legende{Mindmap of Lesson 02: Economic Life and Occupations (Terminale Technique)}
\end{popfigure}

\begin{bilan}

\courssub{Key Definitions \& Vocabulary}
\begin{itemize}
    \item \cle{Consumer Needs} (Besoins du consommateur) : Goods or services required to satisfy wants (e.g., food, healthcare, banking).
    \item \cle{Community Resources} (Ressources communautaires) : Local assets available to residents (markets, hospitals, credit unions, vocational centers).
    \item \cle{Best Buy} (Meilleur achat) : Product offering the best \emph{value for money} (rapport qualité-prix), not just the lowest price.
    \item \cle{Banking Services} (Services bancaires) : Savings accounts (compte épargne), current accounts (compte courant), loans (prêts), transfers (virements), mobile money (Mobile Money / Orange Money / MTN MoMo).
    \item \cle{Comparison} (Comparaison) : Evaluating similarities/differences in price, quality, durability, warranty (garantie).
\end{itemize}

\courssub{Grammar Reference Tables}

\courssub{Forms of Comparison}
\begin{center}
\begin{tabularx}{\linewidth}{|l|X|X|}\hline
\textbf{Type} & \textbf{Structure} & \textbf{Examples} \\ \hline
\textbf{Equality} & as + adjective/adverb + as & \textit{This bank is \textbf{as reliable as} that one.} \\ \hline
\textbf{Inequality} & not as/so + adjective/adverb + as & \textit{Mobile money is \textbf{not as fast as} a wire transfer.} \\ \hline
\textbf{Comparative (Short)} & adj+-er + than / more + long adj + than & \textit{Cash is \textbf{cheaper than} credit cards.} \\ \hline
\textbf{Superlative} & the + adj+-est / the most + long adj & \textit{This is \textbf{the best} interest rate.} \\ \hline
\textbf{Double Comparative} & The + comparative ..., the + comparative & \textit{\textbf{The more} you save, \textbf{the richer} you become.} \\ \hline
\textbf{Parallel Increase} & comparative + and + comparative & \textit{Prices get \textbf{higher and higher}.} \\ \hline
\end{tabularx}
\end{center}

\courssub{Adverbs of Time (for, since, ago, after, afterwards)}
\begin{center}
\begin{tabularx}{\linewidth}{|l|l|X|}\hline
\textbf{Adverb} & \textbf{Use / Tense} & \textbf{Example} \\ \hline
\textbf{For} & Duration (Period) + Perfect/Cont. & \textit{I have had this account \textbf{for} 5 years.} \\ \hline
\textbf{Since} & Starting Point + Perfect & \textit{She has worked here \textbf{since} 2020.} \\ \hline
\textbf{Ago} & Past Point from Now + Past Simple & \textit{He opened the account 2 months \textbf{ago}.} \\ \hline
\textbf{After} & Preposition / Conjunction (Then) & \textit{\textbf{After} signing, you get the card.} \\ \hline
\textbf{Afterwards} & Adverb (Later, then) & \textit{We applied; \textbf{afterwards}, we waited.} \\ \hline
\end{tabularx}
\end{center}

\courssub{Complex Sentences: Coordination vs. Subordination}
\begin{center}
\begin{tabularx}{\linewidth}{|l|l|X|}\hline
\textbf{Type} & \textbf{Linkers (Conjunctions)} & \textbf{Structure / Example} \\ \hline
\textbf{Coordination} & FANBOYS: For, And, Nor, But, Or, Yet, So & \textit{[Ind. Clause] , \textbf{but} [Ind. Clause]}\\
& & \textit{I need a loan, \textbf{but} my score is low.} \\ \hline
\textbf{Subordination} & \textbf{Time}: when, while, as soon as, before, after, until & \textit{\textbf{When} I receive my salary, I save.} \\
& \textbf{Reason}: because, since, as & \textit{\textbf{Because} rates are high, I wait.} \\
& \textbf{Condition}: if, unless, provided that & \textit{\textbf{If} you sign today, you get a bonus.} \\
& \textbf{Contrast}: although, even though, whereas & \textit{\textbf{Although} it's far, the service is good.} \\
& \textbf{Purpose}: so that, in order that & \textit{I save \textbf{so that} I can buy a car.} \\ \hline
\end{tabularx}
\end{center}

\courssub{Expressing Preferences}
\begin{center}
\begin{tabularx}{\linewidth}{|l|X|}\hline
\textbf{Structure} & \textbf{Example} \\ \hline
\textbf{Prefer + noun/gerund + to + noun/gerund} & \textit{I \textbf{prefer cash to credit cards.}} \\ \hline
\textbf{Prefer + full infinitive + rather than + bare infinitive} & \textit{I \textbf{prefer to save rather than spend.}} \\ \hline
\textbf{Would rather (’d rather) + bare infinitive (+ than)} & \textit{I \textbf{’d rather use} Mobile Money \textbf{than go} to the bank.} \\ \hline
\textbf{Would prefer + full infinitive (+ rather than)} & \textit{I \textbf{would prefer to talk} to a manager.} \\ \hline
\textbf{Better + bare infinitive (advice)} & \textit{You \textbf{’d better check} the interest rate.} \\ \hline
\end{tabularx}
\end{center}

\courssub{Express Methods (Méthodes Express)}

\begin{methode}[Comparing Prices for Best Buys (Smart Shopping)]
\begin{enumerate}
    \item \textbf{Identify needs:} List exact specs (quantity, quality, features).
    \item \textbf{Gather data:} Check 3+ sources (market, supermarket, online, wholesale).
    \item \textbf{Calculate Unit Price:} $\frac{\text{Total Price}}{\text{Quantity}}$ (e.g., FCFA/kg, FCFA/litre).
    \item \textbf{Check Hidden Costs:} Transport, VAT (TVA), warranty fees, maintenance.
    \item \textbf{Evaluate Quality:} Durability, brand reputation, expiry date, certification (ANOR).
    \item \textbf{Decide:} Choose best \emph{Value for Money} (Quality/Price ratio).
\end{enumerate}
\end{methode}

\begin{methode}[Writing a Formal Letter/Email to a Bank]
\begin{enumerate}
    \item \textbf{Header:} Your Address (Top Right) $\rightarrow$ Date $\rightarrow$ Bank Address (Left).
    \item \textbf{Subject Line (Objet):} Clear, bold, caps (e.g., \textsc{Application for Loan / Complaint: Failed Transfer}).
    \item \textbf{Salutation:} \textit{Dear Sir/Madam,} (or \textit{Dear Mr. [Name],}).
    \item \textbf{Opening:} State purpose immediately (\textit{I am writing to apply for... / to complain about...}).
    \item \textbf{Body (2-3 paragraphs):} Give details (Account n°, dates, amounts). Use \textbf{Subordination} (\textit{Because, Although, If}) and \textbf{Coordination} (\textit{But, So}).
    \item \textbf{Call to Action:} \textit{I would appreciate it if you could... / I look forward to your reply.}
    \item \textbf{Closing:} \textit{Yours faithfully,} (if name unknown) / \textit{Yours sincerely,} (if name known).
    \item \textbf{Signature + Printed Name + Phone Number.}
\end{enumerate}
\end{methode}

\end{bilan}

\begin{aretenir}[Checklist avant le Bac]
\begin{itemize}
    \item $\square$ I can define \cle{Consumer Needs} and \cle{Community Resources} with local examples (Cameroon context).
    \item $\square$ I master the \textbf{5 forms of comparison}: Equality, Comparative, Superlative, Double Comparative, Parallel Increase.
    \item $\square$ I distinguish \textbf{For} (duration) vs \textbf{Since} (start point) vs \textbf{Ago} (past simple marker).
    \item $\square$ I distinguish \textbf{After} (prep/conj) vs \textbf{Afterwards} (adverb).
    \item $\square$ I can build \textbf{Compound sentences} using FANBOYS (Coordination) correctly punctuated.
    \item $\square$ I can build \textbf{Complex sentences} using subordinators (Time, Reason, Condition, Contrast, Purpose).
    \item $\square$ I use \textbf{Preference structures} accurately: \textit{Prefer X to Y}, \textit{Would rather V}, \textit{Would prefer to V}.
    \item $\square$ I know key \textbf{Banking Vocabulary}: Deposit/Withdrawal, Interest Rate, Overdraft, Standing Order, Mobile Money, Collateral.
    \item $\square$ I can read a bank statement / loan offer and extract key info (APR/TAEG, maturity, fees).
    \item $\square$ I can write a formal \textbf{letter/email to a bank} (Application, Complaint, Inquiry) with correct layout.
    \item $\square$ I can role-play \textbf{assessing needs} (interviewing a consumer) and \textbf{comparing prices} (shop dialogue).
    \item $\square$ I can justify a \textbf{Best Buy} decision using \textit{Value for Money} arguments (not just price).
\end{itemize}
\end{aretenir}

\findecours

\end{document}
